% ============================================================
%  R. Nielsen, Om personlig Sandhed og sand Personlighed
%  [On Personal Truth and True Personality]
%  Twelve lectures for an educated audience of both sexes,
%  delivered at the University in the winter of 1854.
%  Copenhagen: Gyldendalske Boghandling (F. Hegel), 1854.
%  ([1] + 144 pp., Fraktur)
%  TRANSLATION (English)
%  Source of truth: transcription.tex (Danish), COMPLETE.
%  Mirrors transcription.tex structure 1:1. See TRANSLATION-PLAYBOOK.md.
% ============================================================
\documentclass[12pt,a4paper]{book}
\usepackage[utf8]{inputenc}
\usepackage[T1]{fontenc}
\usepackage{libertinus}
\usepackage{libertinust1math}
\usepackage{textalpha}   % Greek typed directly, if/when it turns up
\usepackage{graphicx}    % for \reflectbox (the old Danish »ɔ:« id-est mark, if it turns up)
\DeclareUnicodeCharacter{0254}{\reflectbox{c}}  % ɔ (U+0254) = reversed c
\usepackage[english]{babel}
\usepackage[protrusion=true,expansion=true]{microtype}
\usepackage{setspace}
\usepackage[margin=1.2in]{geometry}
\usepackage{enumitem}
\usepackage{fancyhdr}
\setlength{\headheight}{28pt}
\addtolength{\topmargin}{-13.5pt}

% Footnotes: the Danish uses a constant *) mark (see transcription.tex preamble).
% Mirror the same convention here — one footnote in the whole book, printed p.49.
\renewcommand{\thefootnote}{*)}

\usepackage{hyperref}

\hypersetup{
  pdftitle={On Personal Truth and True Personality},
  pdfauthor={R. Nielsen},
  hidelinks
}

\pagestyle{fancy}
\fancyhf{}
\fancyhead[LE,RO]{\thepage}
\fancyhead[RE]{\textit{On Personal Truth and True Personality}}
\fancyhead[LO]{\textit{\leftmark}}

\setstretch{1.3}

\title{\textbf{On Personal Truth and True Personality}\\[6pt]
  \large Twelve Lectures for an Educated Audience of Both Sexes,\\
  Delivered at the University in the Winter of 1854}
\author{Lic.\ theol.\ R.\ Nielsen\\\normalsize Professor of Philosophy}
\date{Copenhagen: Gyldendalske Boghandling (F.\ Hegel), 1854 \\ \normalsize (translated from the Danish)}

\begin{document}
\maketitle
\thispagestyle{empty}
\newpage

\tableofcontents
\newpage

%% ============================================================
%%  FORORD / PREFACE (one unfoliated leaf; synthetic marker 0--0,
%%  mirroring transcription.tex)
%% ============================================================
\chapter*{Preface}
\addcontentsline{toc}{chapter}{Preface}
\label{ch:forord}

To speak about personality is an easy matter; but to speak personally and truly about personal truth is not merely a difficult undertaking — strictly speaking, it is a superhuman task.

That there are writings in our literature which solve the task by a standard that does not apply to my contributions, I have already, and, as it seems to me, emphatically, pointed out before.

The fellow-striving reader who holds fast to the rule that every honest endeavor is best measured by its own measure will doubtless not mistake the good will in my contributions.

Should some researcher or other then unhappily discover that the present writing has miscarried, since the author has not yet found himself able to furnish the purely scientific; such a discovery will not matter, if one only bears in mind that what is here offered is not meant to be a fragment either of a new or of an old science, but only something personal about the personal.

\begin{flushright}
Copenhagen, 1 May 1854.

R. Nielsen.
\end{flushright}

%% ============================================================
%%  I. INTRODUCTION: A FANTASY (printed pp. 1--11)
%% ============================================================
\chapter*{I. Introduction: a Fantasy.}
\addcontentsline{toc}{chapter}{I. Introduction: a Fantasy}
\label{ch:forel1}
\markboth{I. Introduction: a Fantasy}{}

If there is any question that can rightly be called a question of life, it must surely be the question of the eternal life; if there is any cognition of which it indisputably holds that it concerns every human being personally, it must surely be the cognition of personal truth; if every human being in the race is truly disposed toward true personal development, then true personality must indeed be the prototype of each individual human being and, consequently, of the whole race.

``The eternal life,'' ``personal truth,'' ``true personality'': these different expressions point to one and the same thing. If I only knew that there truly was an eternal life, a judgment, a transfiguration, a reunion awaiting us all, and knew it on sufficient grounds and with complete assurance: then I would also know what dwells within human personality --- that the Eternal itself is there within, however quietly and silently, so that it is not perceptible to the senses.

And could I see so sharply that my gaze penetrated to the ground of personality, and discovered the hidden kernel; and could I draw this kernel forth into the light of day, and examine it closely, and satisfy myself whether it is
% --- p. 2 ---
perishable or imperishable: then I would at the same time know what I should think and judge concerning the eternal life.

Is it then possible to see through the essence of personality? Is there not, after all, something enigmatic in the inward part of our being, something that is not exhausted in a concept, something that no human knowledge can rightly clarify? Is it not better, as regards this point, at once to abandon every attempt at comprehending knowledge, and to remember what is written: ``we walk by faith, and not by sight''?

If I bring these questions before the tribunal of the spirit of the age in the present, a sharp reproof is heard: ``Do not be so mysterious; the so-called unfathomable has already been fathomed. What the Eternal is, we know; what the Personal is, we know also: the Eternal is not personal, the Personal is not eternal. The Personal is born in time, and must therefore die in time; it has a beginning, and must therefore have an end. Were you a true thinker, with a sense for the real (and not an unclear, sluggish brooder), you would know this too; for we all know it.'' --- All? --- ``Yes, of course, all the knowledgeable, all the well-informed. That the great, half-educated crowd does not yet know it, or at least prefers to act as though one could not know it, is natural enough. Life has its dark sides; the terrors of transience ought not to be inhumanly bared to human eyes: so let religion, for as long as it can, spread its veil over this. But do notice how this veil grows ever more airy and more transparent. Notice how, in the present, the interests of reality more and more lay claim to every thought and feeling: the time is not far off when the race as a whole will find itself so at home on earth that there can no longer, in full earnest, be any question of a heaven.''

So speak the spirits of the age, and what shall we answer? Is the eternal life a fantasy; is the resurrection and the judgment and
% --- p. 3 ---
the heaven of reunion --- is all this now really to be explained as fantasies: then let us begin with a fantasy.

``There are,'' says Hamlet, ``more things in heaven and earth than are dreamt of in your philosophy.'' Let us then begin by imagining that there was a hitherto undiscovered place, a place the philosophers do not dream of, a place far off on the earth, from which one could see the land of eternity, look into the regions of transfiguration, and hear wondrous tones. Let us imagine that a solitary traveler once came, as if by chance, to this place, saw the vision, heard the tones, and afterward returned, and tried to describe what he had heard and seen: how would such a one be received? --- ``As an enthusiast, a madman!''

Yes, but if the traveler were also a man who knew the world, who had a sharp eye for the real, and who could give an exact account of the impressions he had otherwise gathered on his many travels (and he had seen much, and been much about: on the heights of the Himalayas, at the sources of the Nile, by the rivers of Babylon, among the ruins of Nineveh he had been, and could speak of it, so that those who listened to him seemed themselves to be making the journey with him, and to experience what he described) --- only, when he came back to that inexplicable point, indeed when he came back to ``the vision,'' and wished to describe the indescribable, he became as though transformed; then he spoke as if to himself, then his listeners could well sense that there was something attractive, something wondrously gripping in his words and thoughts, but no one could hold on to it, no one could grasp it, for he spoke as in riddles; --- if the traveler, we say, were such a man, what then?

It could surely never fail that there would be one or another among the listeners who felt a desire to travel after him, to seek out the disputed place, and to satisfy themselves how it really stood with the alleged aerial vision. And when this one then returned with the same account,
% --- p. 4 ---
and yet a third, and yet a fourth, and ever more: what then?

Would the enlightened of this age, who have seen behind the veil, and seen that there is nothing but the transient, receive these reports with perfect coldness, only repeating to themselves: ``the Eternal is not personal, the Personal is not eternal''? I do not believe it; I believe that these certain ones would soon grow somewhat uncertain. Still less can I suppose that these certain ones would have spiritual superiority enough to govern the movement of the current. For is it not far more natural, even judging by the examples we see in the present, to suppose that the tidings of that wondrous vision, confirmed day by day by more and more people, would transform itself into a spiritual power, soon stronger than all the powers of the spirit of the age combined? Is it not likely that something would awaken within the interior of human beings, something that would rise the more overwhelmingly the longer it had lain dormant, a drive toward the supersensible, a yearning for the beyond, a romantic longing, fantastical, unconquerable? That there would begin a pilgrimage, an adventurous thronging, a stirring, a movement, as in those times when the hermit inflamed the multitudes for the Holy Sepulchre, and the Pope pronounced the blessing?

The thoughts of the age change with the age; but the nature of man is the same at all times.

Let us then carry the image further, and imagine the numerous throngs gathered at that unknown place where the vision appeared: would they greet the supernatural with rejoicing, as one greets a smiling prospect? By no means. They would fall silent with awe, they would lose themselves in quiet contemplation. The usual rule, that where the crowd is, there also is the noise, would not hold here. Here there might be millions, and yet not a disturbing cry be heard; here there might be a world of human beings, and yet a world as still as when Adam was created, and the first breath of air passed through Paradise.
% --- p. 5 ---
In this silence humankind would then listen to the tones, and the more they listened, the more the harmonies would become intelligible to them; it would resound again within them: they would be drawn involuntarily toward the region from which the tones came; they would move forward as if in time and rhythm: the further it went forward, the stronger the movement would become. If nothing stopped them, they would surely forget everything that lay behind, and, strangely seized by a blessed presentiment, would journey, each with their own, straight into the land of eternity.

But if they were now stopped; if a sea of fire were now suddenly laid between them and the shining shores, and it were told them, as though from heaven, that no one could enter the land of the blessed who did not have courage to cast themselves into the flames, and let the earthly part be burned to ashes: what should we then suppose? Should we suppose that these human beings, struck with despondency, would at once turn back, and say to one another: ``For us, then, the eternal life does not exist at all; no, such a thing exists only for gods and demigods. We are earthly, and know, with unerring certainty, that when the earthly burns, we ourselves become nothing. Let us then remain by the earth, and cling to the earthly, for the while we live, we know not how long.''

Should we indeed think so basely of our race, our nature, and ourselves, and dishonor Him who formed us? Is it not, rather, far more fitting to suppose that human beings, inspired by the blessed vision and lured by the holy tones, would cast themselves into the flames in multitudes, assured that the earthly was nonetheless not their essential part, but that the true self was consecrated to the Eternal?

And if the trial became still more severe; if it were not now a sea of flames that made the separation between here and there, but an insurmountable wall with countless narrow entrances, one entrance for each person; and if at the entrance there stood a guide, appointed to warn each one that no one should venture heedlessly out over the
% --- p. 6 ---
finite; and if the guide now began to set before the solitary one how long and thorny and toilsome the way was; that it stood with those shining shores as with the distant mountains that appear to the traveler: they seem to be so near, and yet are far off; that it might well take fifty, perhaps seventy years, before the wanderer reached the goal: what should we then suppose? Should we suppose that the human being, disheartened at the thought of the narrow way through so many years, would at once turn back, saying within himself: ``It is too much. A leap into the flames I would gladly have ventured, for glorious is the prospect, and there is something wondrously blessed in these tones; a leap into the flames, and it would be done at once; but --- so long a way, 50, perhaps 70 years under burdens and hardships, no, it is too much! Now I feel for the first time how finitude weighs; it is temporal time that presses on every child of temporality: I am then dust, and shall become dust again, and with that it is over.''

Should we think so little of our race, our nature, and ourselves, and dishonor Him who formed us? Or is it not, rather, far fairer to suppose that the human being would understand that 70 years with burdens and hardships count for nothing against an eternity? For the human being is, after all, a thinking being. The thinking one reflects, and reflection tells us indeed how the years vanish. What, then, are 70 years, once they have vanished? Everything then gathers into a sigh, into an unspeakable melancholy over what has vanished. But the self that sighs, and is in melancholy, thereby confesses that it lacks something, and desires something, something that does not vanish: it lacks and desires the Eternal.

Ought we not, then, rather to suppose that the human being, inspired by the vision and animated by the tones, would say to the guide: ``Then I give thanks that there is a distance between here and there; for then there is indeed a preparation in the temporal. And I feel it, I am in need of preparation, I could not possibly, as I am, enter into the land of blessedness: what
% --- p. 7 ---
stirs within me here is so great and mighty that it would surely burst my being, were I not granted time for preparation.''

And if the preparation were now exceedingly strict, if the guide had still further scruples, and said to the human being: ``Take good care that you do not regret what you do. You perhaps imagine that you are to walk the long way with this blessed vision before your eyes. That will not happen. On this height you cannot remain. What you now see will again hide itself behind the horizon. The first step goes downward; the way goes downward, and the valley is dark. On this way the sun of day does not shine; only the star yonder, the twinkling star, shall be your light and your guide.'' And if now the human being, who saw the star, how brightly it shone --- as brightly as one can imagine it once shone for those kings in the holy night --- would then at once and confidently set out on the journey, and yet was stopped by still further scruples, if the guide continued to speak and repeat: ``Test yourself! On the long way weariness comes. I tell you: you will grow weary on the way. The weary wanderer so gladly wishes to stop, and look back a little; but the one who, on this way, stops and looks back, is lost. Hear an old legend: Orpheus descended into the underworld; he longed for his wife. In the kingdom of the dead the tones of the lyre were heard; it moved Proserpina, it stirred Pluto, it lulled Cerberus: the pleader's lament was heard; but when the one who was heard was then to lead his beloved, he looked back, and Eurydice vanished like a shadow. Hear a sacred legend: the angels came to Sodom at the hour of judgment, to warn the righteous. And at the angels' word the chosen departed from the city and from the destruction, that they might be saved from the judgment; but Lot's wife looked back, and Lot's wife was turned into a pillar of salt. Hear Truth itself, personal truth: ``No one who puts his hand to the plow, and looks back, enters into the kingdom of heaven.'' The one who looks back, hesitates; the one who
% --- p. 8 ---
hesitates, becomes doubtful, the doubtful one doubts and is confounded. Should it befall you that you look back, then the star's light could so easily be extinguished. The star of the future shines only for the one who has his all in the future; but the one who has his all in the future dare not turn, and seek counsel from the past. And one thing more. As brightly as the star shines in this present moment, it will not always shine on your way. I tell you: you will grow weary on the way. From weariness comes despondency; in despondency hope wanes, and as hope wanes, the star's brightness dims: it veils itself in mist. In heavy hours, in dark moments, when you most need light, the light will be as if extinguished. When there is darkness within you, there will also be darkness around you, and then, just then, the way will be rugged, steep, and trackless. Consider this now, test yourself, and take good care what you do. If it is only a mirage, a distant image of light, after which you walk, then turn back, for on this way you will perish. But if you are assured that what shines for you yonder also shines within you, as the essence of your own being, then be of good courage, take up your walking-staff, and you shall reach your goal.''

If now all this were seen as in reality, and heard as in reality, what should we then suppose? Should we suppose that the human being could in the end confess to himself that the personal ideal of eternity is nonetheless only a mirage, a phantom image, and find peace in the recognition that the Eternal is not personal, and the Personal not eternal? Should we think so basely of our race, our nature, and ourselves, and dishonor Him who formed us?

There is a difficulty that applies to the present age above all preceding ones, and it is a difficulty with fantasy.

Without fantasy, no personality; as fantasy is in its fundamental form and governing direction, so also is personality. If fantasy, with its ideal activity, is taken away, what then remains in a human being? Five senses, a
% --- p. 9 ---
dry understanding, and a few finite passions. No impression of an idea, no spiritual stirring is possible without a corresponding activity of fantasy.

It is not only the poet and the artist who must be aided by fantasy: every human being personally needs the aid of fantasy, in order to be human.

What is honor, when we take fantasy away? What becomes of the sense of propriety, of modesty, and of the nobler self-feeling, once fantasy has first been dulled, or even corrupted? And now joy! Joy, as it sparkles in the eye and plays in the smile --- does it not weave its lovely, animating images in fantasy? And now sorrow, the quiet, earnest, and yet so beneficial, mind-purifying sorrow --- it would indeed depart along with fantasy, and leave behind the pains, the purely physical pains. For physical pains can be felt without any fantasy at all; but an inward sorrow cannot be felt without a corresponding inwardness of fantasy: fantasy nourishes sorrow, and sorrow in turn cultivates fantasy. And now love itself, and longing and hope, and every breath toward the infinite that expands the breast: what becomes of all this, when fantasy departs?

``But,'' one may perhaps ask, ``is there then really anyone in our age who thinks of abolishing fantasy? True enough, our age is not fantastical: it knows how to distinguish between dream and reality; but that is precisely an advantage, for which we may thank the Enlightenment. The Enlightenment clarifies fantasy, and drives out fancifulness; but fancifulness is, after all, precisely the abuse of fantasy.''

I answer: here is a difficulty in the present age that has not existed in the same way in any earlier age, and it is a difficulty with fantasy. It must be granted that the modern age has clarified a distinction between the real and the fantastical which earlier ages did not see through in the same way; but that the distinction between fantasy's ideal right and its
% --- p. 10 ---
real, material right has thereby come to be understood in a manner just as one-sided as it is unsatisfactory, I maintain; and I prove my claim by pointing to the religious.

All fantasy is, according to the prevailing opinion of the present day, illusory; but illusion is of different kinds. There are empty, unreasonable, idle, fruitless illusions; they all belong under the fantastical. There are essential, beautiful, uplifting, beneficial illusions; they all belong under the aesthetic.

Where, then, does the religious belong? Naturally, under the aesthetic. The difference between religion and poetry is, accordingly, only this, that poetry addresses itself most directly to our sense of beauty, religion to our moral sense, both through fantasy. If I am to be edified by the illusions of beauty, I must go to the theater; if I am to be edified by the illusions of the good, I must go to church.

With such a juxtaposition the religious cannot stand. The kingdom of religion is at once a kingdom of reality, and yet a supersensible kingdom. Either there must be a fantasy that grasps the real and is grasped by it, in that it is grasped by the supersensible; or we must simply grant that those are right who maintain that the Enlightenment, which has abolished fancifulness, has at the same time transformed religion into an edifying, humane illusion.

If they are right in this claim, then they are right in more; for then they are also right when they flatly declare the eternal life, resurrection, retribution, and reunion to be illusions, nothing but illusions; then they stand firm on the ground of reality, who stand on the proposition: the Personal is not eternal, the Eternal is not personal.

Should this, then, be the highest personal truth? Should the highest personal truth be so personally revolting?

The Middle Ages were, in their main direction, religious; a proof of this is that they had a distinctive, actively creative religious fantasy.
% --- p. 11 ---
The Reformation was, at its root, religious, and what wondrous powers of fantasy stirred in Luther.

Can the present age, then, likewise be said to have a distinctive, creative religious fantasy?

Just as little as the individual human being can give himself the gift of fantasy, just as little can an age do so. Nor is it needed, for the gift is given by Him who gives every good and perfect gift. Poetic fantasy, which corresponds to a distinctive natural endowment, is not given to all, and is not directly necessary for all. Religious fantasy, by contrast, which is necessary for all, does not correspond to a peculiar natural endowment; it corresponds to a distinctive personal direction of spirit, and is given in proportion to the striving of the will and the earnestness of the mind.

Whether the productions of religious fantasy are, in their phenomena, more or less beautiful, is no principal question; the principal question is whether these productions personally correspond to personal truth. In nature the flowers are beautiful, that is to say: most of them; yet not all, for the flower is natural not for the sake of beauty, but for the sake of the fruit. \emph{By their fruits shall ye know them.}

%% ============================================================
%%  II. AESTHETIC AND RELIGIOUS FANTASY: A PERSONAL DIFFERENCE
%%       (printed pp. 12--22)
%% ============================================================
\chapter*{II. Aesthetic and Religious Fantasy: a Personal Difference.}
\addcontentsline{toc}{chapter}{II. Aesthetic and Religious Fantasy: a Personal Difference}
\label{ch:forel2}
\markboth{II. Aesthetic and Religious Fantasy: a Personal Difference}{}

Poetry and religion resemble each other in this, that both address our feeling, and act upon our fantasy, in that both address us personally. Nevertheless there is a very essential difference between the aesthetic and the religious direction of fantasy.

% NB: the Danish heading reads "Forskjel" but the body text on this page reads
% "Forskiel" (no j) -- a genuine printer's spelling variant on the same page,
% noted in transcription.tex; the two spellings are indistinguishable in
% English translation ("difference").
This difference may briefly be expressed thus: for poetry the personal is a game, in religion it is a holy earnestness; poetry looks to personal illusion, religion to personal truth; poetry conjures up an illusory personality, in religion the true personality is revealed.

If we choose examples from poetry, everyone knows that its creations are illusory. Axel and Valborg, Hakon and Thora are now as if at home with us; we know these personalities; we are familiar with them; we have almost lived with them as with real friends; and yet they all belong to illusion: Axel and Hakon, as well as Thora and Valborg.

That there once really lived in Norway a real Hakon Jarl is a historical remark that does not concern our Hakon. When our Hakon is on the stage, he has for us all the reality he shall and will have. We need,
% --- p. 13 ---
then, no historical testimony of his existence; we see him himself, the lofty figure, vividly alive in the illusion.

As these figures now show themselves to us in illusion, they invite us to follow them, and to share their passions, their sorrow and joy, their delight and pain, personally with them. We do so gladly, we give ourselves over: the invitation is welcome to us. Then we forget ourselves and our own reality. Our thoughts, our feelings, our whole personality now belong, for a moment, to these exalted ones; only in that we thus belong to them, do they belong to us.

So Hakon storms, wild in battle, and when Hakon storms, I am with him in the battle: it surges in my mind. I see him everywhere, I follow him, I hear him, for I am caught up in the mood. If I am not caught up in the mood, as I ought to be, I do not see him, and do not hear him at all. I hear a man shout, but Hakon I do not hear.

When Thora stands by Hakon's coffin in quiet sorrow, bowed down and melancholy; then I too am there by Hakon's coffin, caught up in the mood. I then weave a memorial wreath, as Thora weaves it for me: ``he was a rare hero;'' --- we lay the wreath upon his coffin. ``He was, in the North, a rare hero,'' --- we write the hero's name into the wreath, that this name shall never be erased, and never forgotten.

So strong is the illusion; it seems earnest, and yet is not earnest. We seem to be caught up in the struggle and reconciliation of the personalities; in such moments we live through so much of sorrow and joy, delight and pain, fear and danger, so much as could not otherwise be lived through in many years; and yet it is all illusion.

We become clear-sighted; it seems to us --- so glorious is poetry --- as though we saw human beings and existence into the very heart, saw life's riddles and their solution; and yet it is all illusion. It seems to us as though we saw into ourselves, and saw how deeply, at bottom, we love the good and hate the wicked; how we indeed sympa-
% --- p. 14 ---
thize both with courage and valor and steadfastness, and with that faithfulness of love which will in the end prevail. Yes, so glorious is poetry, so humane and so ingratiating, that it acts upon us as though we ourselves had become both courageous and magnanimous, and pious and virtuous, indeed even witty. And yet it is all illusion, alas, nothing but illusion.

But there is nonetheless a freedom of the spirit, an inexhaustible richness of personality, in poetry. The poetic spirit is personally free; its fantasy is so deft, supple, exuberant, daring; it is like a juggler who plays lightly with sharp knives: we grow anxious, we are entertained, when it flashes.

Poetic fantasy is personally rich, and freely audacious, mighty, like the wealthy: here is Aladdin himself, with the spirit of the ring and the wondrous lamp.

Yes, the poetic spirit is personally rich; --- I could go on, --- for it is rich in adventure without end. Its watchword is that life is an adventure. So it is itself an adventurer, an enterprising, yet lighthearted idler, who never lays anything by, but lives on what the moment gives. It is a musician who travels about; it is that strange fiddler who once came to the King's court.

``He touched the strings, and the King laughed, and the whole court laughed; it was so merry to hear. He touched the strings, and the King wept, and the whole court wept; it was so sorrowful to hear. He seized the strings with fire, a wild, a daring grasp, and, as if maddened, the King started up, and drew his sword, and slew his men where they stood. Then the strange fiddler vanished. The King stood alone in his hall; he brooded and was so dark; he was so deeply grieved. He sought comfort, but nothing, no one could comfort him. Then he seized the pilgrim's staff, summoned an assembly, and spoke to the people. And the people heard him take his leave, and the people mourned, and prayed with tears: ``he was, after all, so good, so truly good; the good King must stay, --- stay with his people.'' But alas,
% --- p. 15 ---
he could not; the staff burned in his hand. The good King could not stay; he had to go away, far away --- to the river Jordan, to Golgotha, and to that wellspring which washes away guilt, and so he went away.''

In this image, recalled from our own Middle Ages, we meet poetic fantasy at the crossroads where it fantastically crosses the religious. It seems almost as though it were earnest, for the King departs; he departs truly out of repentance. Nevertheless it is not full earnestness; it is fantastical earnestness, and so it is an adventure.

For religion, the adventurous gives way. Religion begins in plainness; it takes reality as it is, and lays its burden upon us. It looks so everyday. We feel the burden, where it weighs, where it presses; but this pressure of temporality signifies precisely an emphasis of eternity upon our personality: here is earnestness, holy earnestness.

By personal earnestness revelation is distinguishable from mythology. In the various myths the human spirit has adventurously expounded its own natural content; but religious fantasy has thereby strayed at the same time into the aesthetic: it has not become earnest with the personal.

The mythical gods are not personalities; they are only personifications. They borrow the semblance of personality, but personal earnestness and personal truth are not in them.

Zeus knits his brows, Olympus trembles: it looks like earnestness. He threatens Hera, it looks dangerous: the Queen of Heaven turns pale and trembles. Yes, it looks dangerous enough; but it is not dangerous, it passes over. It is a whim of the good Zeus, it is a cloud before the sun. It lasts only a moment; the cloud passes by: it clears up. The limping Hephaestus trips about, and speaks fair words; he comforts Hera with the cup, and all the gods laugh, and Zeus laughs along; for the whim is at play: fantasy goes on an adventure.

In no mythology, neither the Greek nor the Norse, is there personal earnestness; in none of mythology's
% --- p. 16 ---
divinities, in Odin as little as in Zeus, in Thor as little as in Ares, is there personal truth.

Yet, it is true: I ought to mention here a divinity, this place's exalted protectress, Minerva, --- Pallas Athena with spear and helmet, and with the Gorgon on her shield; yes, it is true, in her there is earnestness, in her being there is truth. She does not jest, does not trifle; she is severe, as she is pure and beautiful. She represents the earnestness of science. She is in armor: for the truth of science she will contend. But science itself is, like the goddess, --- impersonal. See, therefore she received this strange, sharp gaze, but no heart. Never does she speak the language of personality; she cannot, nor does she wish to; she speaks the pure, clear language of knowledge; but toward personality she is cold and mute: she is of marble.

The perfect religion begins with a revelation; to revelation, faith corresponds. These two are inseparable: faith comes only through revelation; revelation is only for faith. for the believer; both meet in personal truth.
% [printed: period after "faith" ("Troen"), not a comma -- printer's defect;
% sense requires "for faith, for the believer," transcribed as printed.]

What, then, is it that is revealed? Many visions, many strange signs! Here the eye may easily be confused; here the outer is easily mixed with the inner. Where, then, is the center, around which everything moves, the One for whose sake everything else appears and happens? I answer: \emph{personal truth, true personality.}

It is, after all, God himself, in his own holy person, who reveals himself to the world. ``Is it possible, is it real?'' --- For faith, for the believer!

Yes, for the believer this is real; for the God of revelation is not a fantastical being that hovers in the clouds, is not a mythical personification of some single natural force or some single passion, nor is he a God of marble; he is a human being in flesh and blood, like one of us, and yet --- it was indeed said of him: \emph{behold, what a man!}
% --- p. 17 ---
But, the rational people of this age may perhaps answer: ``that may indeed be very fine for whoever merely wishes to believe; but this age is not an age of faith. We have thought too much, seen and heard too much, to believe in such things.''

It is, however, a harm, an eternal harm to personality, if anyone becomes so rational that he can no longer believe in such things. For it takes only a little reflection to see that, when this point is violated, then in the whole of existence there is nothing inviolable. If there is not full earnestness about the holy right of personality, God's right over human beings, then in the whole world there is no right and no earnestness; for all that the passing day otherwise happens to call right and earnestness, what is it, after all? A passion, a fantasy, a play of illusions born of the moment's whims! If He who says: \emph{I am the Truth,} is not Truth itself in his own holy person, but only a phenomenon long since vanished: what, then, is the personal? A fleeting wave-beat in the ocean of the race. The excellent one is then like the high, the shining wave. But the shining wave is leveled out just as much as the dark one: the high wave, puffed up by the storm, and foaming, is leveled out, and comes to rest, just as much as the low one.

It is, then, a harm, an eternal harm to personality, if anyone becomes so enlightened that the strong light in the end darkens the truth. The true personality, in whom I am to believe, does not, after all, come to me from outside as a stranger, one otherwise unconnected to me. To look upon him is at the same time to look into oneself; to look up to him is to look down into the depth of the mind; for he is the depth.

The true personality is the one who reveals himself, who opens and reveals each human being's inward part to that human being himself. This, then, is the center of revelation, the kernel of faith. We now return to religious fantasy.

In order to keep as close as possible to the so-called real, we will begin by imagining that in the phenomenon of revelation there was as little for fantasy as possible.
% --- p. 18 ---
The true personality would then, in the fullness of time, have appeared as a human being among human beings; and it would nonetheless be a wonder how this human being had come into the world. But no sign, no portent would have been given outwardly. No angel of annunciation would then have greeted ``the favored one,'' no multitude of the heavenly host proclaimed his birth to the shepherds, no star have risen for the wise men of the East.

As an isolated individual, ``without father, without mother, without genealogy,'' the ideal would suddenly have stood among human beings, a true wonder, and yet without wonder and sign. As he then was in his appearance, so he would presumably also have been in his word. ``I am personal truth; you cannot comprehend it. But now follow my example, obey my commands, and believe in me, and you shall live forever.'' Thus, presumably, he would have spoken. His word would then nonetheless have been what he himself was, the eternal, personal truth; but human beings surely could not have grasped this cold word: fantasy, after all, could see nothing, and grasp nothing.

Therefore he does not come in this way either, and does not speak in this way. The true personality speaks the language of personality; but the language of personality is at the same time the language of fantasy.

The true personality is not content to say: ``Among the many who hear my word, there are only a few who grasp it.'' In awakening reflection, he speaks to fantasy; so he speaks in parables. ``A sower goes out to sow his seed''; we see that sower before our eyes. ``And some falls among thorns, and thorns grow up and choke it''; then the word penetrates, and troubles and warns. ``And some falls on good ground, and bears much fruit''; then the word penetrates, and comforts and edifies.

We do not merely hear, we see --- every trait in this parable, and the truth in each of his divine parables: when the shepherd seeks the lost; when the woman lights a lamp, and searches for her coin; when the father rises, and opens his
% --- p. 19 ---
arms for the prodigal son, then it becomes alive for us, then fantasy can stir, and the mind be gripped.

These parables, however beautiful they may be, are nonetheless not for the aesthetic; they exist essentially only for religious fantasy. The difference is known by the personal.

If I understand the parables merely aesthetically, as beautiful, lovely poems, then I blaspheme the holy. If they are only beautiful poems, then it is quite indifferent who he was who gave us these poems. A poem stands on its own beauty: the poet's personality is indifferent.

Religious fantasy sees the matter the other way around. The personality of the speaker is everything: I grasp him, him himself in the word, I believe him on his word, in that I believe the parable.

For aesthetic fantasy the parable is a lovely image for contemplation, for enjoyment in contemplation. The observer's personality is indifferent. He merely enjoys the beautiful impression, but is not himself obligated to anything..
% [printed: double full stop after "anything" ("Noget"), transcribed as
% printed -- printer's defect, not corrected to a single period.]
The one who enjoys stands outside the whole, and can otherwise do as he pleases.

Religious fantasy sees the matter the other way around. In the parable it is, after all, spoken to me, to you; it is precisely me, precisely you, who is thereby to be awakened, and obligated, and edified. In religion, then, it is from both sides, both the speaker's and the hearer's, the highest earnestness with the personal. What holds true of the parable holds, in a corresponding way, of the miracle.

The one who reveals himself is himself the miracle, and for the sake of personality I must then believe the miracle in its essence. But this essence is invisible. If it now remained in its invisibility, then I would have to attempt, on my own, to form for myself a conception of the incomprehensible's origin, essence, and nature. How false, how misleading such a self-made conception would have to become, if this incomprehensible were nonetheless to be held fast as in a cold concept.
% --- p. 20 ---
It is otherwise, when the incomprehensible gives the sign, stirs fantasy, and thereby looses the wing of the spirit. Then the angel comes and announces; then the star is kindled over Bethlehem; and the shepherds in the field are illumined, in the still night, with radiance.

The wonders of revelation, in and of themselves, prove nothing. The proof is personality itself; these wonders do not exist for sense and understanding; they exist for fantasy, well to note, not the aesthetic, but the religious fantasy. Religious fantasy takes everything personally and in earnest; it works everything up for ethical use in faith.

When heaven is cold and empty, and eternity like a thought without life; then it is no wonder if a human being becomes a stranger to the Most High. This he knew, who knows our nature, and therefore he opened the heaven of fantasy.

The real heaven of the spirit, with its blessedness, cannot be seen with earthly eyes. Everything must first be transformed: we must be transformed. And yet, we need a sign; this he knew, who knows our nature, and therefore he gave us a sign from heaven.

If I understand him in faith, then I understand why he speaks so wondrously of these things. In some places it is as though everything became sensuous, in flesh and blood. It is a banquet prepared in heaven. ``Many shall come from east and west, and sit down at table with Abraham, Isaac, and Jacob.'' In other places all such images are again effaced; we are reminded that his word is spirit, and that he speaks of the spiritual.

If I understand him in faith, then I understand why he speaks in this way. A vision shall appear, and I shall see it clearly with my eyes; and yet in the same instant the vision shall vanish, like a portent. He speaks to the disciples on the way. And one among the disciples asks him, and says: ``See, we have left all things, and followed you, what reward awaits us?'' And the Master does not answer in a dead concept, as though he would say: ``the good has its reward in itself''; he answers in an image, a portent: ``In the regeneration, when the Son of Man
% --- p. 21 ---
shall sit upon the throne of his glory, you also shall sit upon twelve thrones, and judge the twelve tribes of Israel.''

--- Twelve thrones, and twelve such judges; so heaven, then, is not empty! but the vision vanishes even as it is seen: it is a portent.

He speaks of the last things, of death, of the resurrection and the judgment; he does not speak in a dead concept, as though he would say: ``The fleshly dies, only the spirit lives; in eternity the good shall be rewarded, and the wicked punished.'' He speaks to fantasy: ``the hour is coming in which all who are in the graves shall hear the voice of the Son of God.''

It is visible, it comes quite near; it is as though we saw graves open, and the dead step forth in the solemn hour. And yet there is here nothing for the sensuous eye, and nothing for a sensuous concept; but faith holds fast to the word, and the word surely keeps what it has promised. It is personality's word concerning right and judgment, and what are right and judgment, if the voice of justice cannot, in the end, open the graves?

``They shall come forth, those who have done good, unto the resurrection of life; but those who have done evil, unto the resurrection of judgment.''

Unto the resurrection of life --- unto the resurrection of judgment! For aesthetic fantasy, which loves illusion, this image is too dim. Religious fantasy, by contrast, holds fast to the real, ethical opposition, so that the representation carries the thought, and awakens reflection. And what a difference!

A painter paints us a Day of Judgment; a Michelangelo terrifies us with a picture. It seems earnest; but it is nonetheless not earnest. The picture stands still, all the figures in their semblance of motion stand, after all, quite still. Thereby the sensuous becomes even more sensuous. In art it can be a masterpiece; but religious fantasy turns away from such an illusion with reluctance. The religious image never stands still; it is a thought-image; it flashes like lightning through
% --- p. 22 ---
the soul, and awakens our conscience; it comes like a portent, and turns our attention to the ethical.

But the aesthetic attention that is turned upon the painted picture lingers with pleasure over each individual figure in the various groups. Instead of separating good from evil, we now come to separate light from shadow. And the shadow is indeed masterly, for Michelangelo is a master at drawing shadow. Aesthetically one has, into the bargain, quite as much delight from those who ``come forth unto the resurrection of judgment'' as from those who ``come forth unto the resurrection of life''; for the damned are masterly; the despairing, with these muscular contortions and these foreshortenings, are almost even more excellent than the blessed.

When the religious fantasy of the Middle Ages had grown weary, it composed a Day of Judgment, it painted a Day of Judgment, and, in its aesthetic enthusiasm, surely did not think that its own real Day of Judgment was close at hand.

There is, in the present, a good deal that reminds us of that dissolution; among other things, this: that aesthetic fantasy so strongly takes possession of the religious.

They are both good gifts, and can both be effective, each in its own sphere; but they must be thoroughly kept apart. To each its own: that is justice, and only justice stands firm before the judgment.

%% ============================================================
%%  III. THE ETERNAL LIFE: A PERSONAL NEED (printed pp. 23--32)
%% ============================================================
\chapter*{III. The Eternal Life: a Personal Need.}
\addcontentsline{toc}{chapter}{III. The Eternal Life: a Personal Need}
\label{ch:forel3}
\markboth{III. The Eternal Life: a Personal Need}{}

As a support for the belief that the soul is immortal, one is indeed often reminded of the need of the human being, of the longing of the heart.

There is a voice that sounds among the graves: ``the dead sleep, they shall wake again; the departed await us: there is a reunion. Believe, then, what your own heart says: you shall meet again; believe, then, the longing in your breast: there is a reunion.''

When we stand among the graves, and hear such testimony, we grant the speaker right. We do not doubt, we believe; we believe what the heart says, we believe the longing in our breast: ``there is a reunion.''

But the human mind is now so changeable. It is one thing to stand among the graves out there, and hear the solemn testimony; another to sit at home, and look upon these real, not always solemn, conditions of life. Here is another atmosphere, another temperature; and the human being is like an instrument that goes out of tune.

Among the graves out there feeling speaks, and then longing is given right: ``there is a reunion''; in everyday life the understanding comes, and gives doubt right: ``there is surely no reunion.''

Feeling and understanding contradict one another; which of the two parties, then, is right? for they wrangle and contend without
% --- p. 24 ---
ceasing.

Sober reflection says: the two must not thus contend; they must be fairly reconciled. True feeling can surely not be so unreasonable, and sound understanding can never contend against the better feeling. Now they must both be rational: they must be reconciled in reason. And so sober reflection calls upon reason, and so reason and sober reflection come to agree on mediating a reconciliation, and by what means? By setting up a kind of proof.

This proof is generally well known; it is very popular, for it is a half-proof, and half-proofs are always the most popular. The proof runs thus:

``Every need has its object, and testifies that there must exist something that can satisfy it. Hunger and thirst indeed testify that there must be food and drink. But now all human beings have a deep, natural need for an eternal life; consequently there must be an eternal life.''

This proof looks, at first glance, very convincing, and it is presumably this first glance that it has to thank for its popularity; for, examined more closely, it can satisfy neither feeling nor understanding.

As a chill can run through one when one suddenly touches cold metal, so a chill runs through feeling, when it must absorb the cold proof of the understanding. And as for the understanding, it is so far from bowing beneath the proof that it, on the contrary, dubs itself a knight upon the proof, and congratulates itself on its great acumen.

Hunger and thirst do indeed remind us of food and drink; but they remind us at the same time that what is to be eaten and drunk is not in the hungry one himself, but outside him; for, if he had within himself that which he needs, he would surely not hunger. Since, then, the hungry one does not have within himself that which he needs, but outside himself, it can easily happen that the hunger is not satisfied. A human being can indeed die of hunger; there are, alas, many who die the death of starvation. If this is now applied to the spiritual, then the hungry
% --- p. 25 ---
soul must indeed be supposed to have the eternal life outside itself; and it could then, for the same reason, happen that many souls too died of hunger.

Let us consider the Greeks. Through Greek fantasy the soul looked longingly toward Elysium, toward the Garden of the Hesperides, toward the Isles of the Blessed, indeed upward to Olympus itself. But from this the thinking Greek did not yet conclude that all who thus longed toward the heights had their longing satisfied. On the contrary, it was a rare exception that distinguished ones, like Hercules, reached Olympus; Agamemnon, Achilles, the souls of many heroes were sent to Hades, to the kingdom of shadows; these languishing shades did not, then, have their longing satisfied. Therefore an old thinker says: ``You shall not desire what belongs to the gods, but be content with what is human.''

Thus, then, does the understanding reason; it is the consequence of the half-proof; those unhappy half-proofs!

That a human being can die of hunger, we all know; but it is not therefore immediately evident to all how much is really said by this. Hunger and satiety, in a word: nourishment is not merely indispensable to the body, it is the very essence of bodiliness, in the plant as well as in the animal and the human being.

``Out of the eater came forth food;'' so runs the first part of Samson's riddle. That the Philistines could not solve this riddle proves precisely that they were genuine Philistines; for the riddle finds its interpretation everywhere in nature's great house of feeding. The plant eats first, in order to nourish itself; the animal in turn eats the plant; the human being then in turn eats both the plant and the animal.

Bodily nourishment is parceled out: it is an extensive parceling-out. My body can be nourished only on the condition that a preparatory nourishment has already taken place in another body; in order that I may be fed, another must already have fed himself --- for my sake; and now there arises anxiety over sustenance, strife and contention in the house of feeding.
% --- p. 26 ---
Each of the eaters, after all, really nourishes itself only for its own sake, and in order to preserve its own kind; nevertheless many of these eaters must submit to being eaten for the sake of others, to become prey and nourishment for other hungry creatures of another kind.

See the dove, pursued by the hawk: what haste, what flight, what whirling, what anguish in the poor dove! The dove, the free bird, is nonetheless, like a slave-girl, an unfree appendage to the bodiliness of the bird of prey.

See the tiger in the thicket, its eyes glinting, when it will spring even upon a human being. The Philistines perhaps think that the tiger hates the human being, that it burns with hatred and malice; that is Philistinism. It is not hatred and malice that shine out of those eyes; no, it is love and sympathy! The tiger looks upon the human being as Adam did, when he said: ``this is flesh of my flesh, and bone of my bones;'' the tiger thinks with Samson: ``out of the eater comes food,'' and so the tiger eats the human being --- out of sympathy.

In all this the many bodies show themselves as members of one great and universal bodiliness; it is one bodily essence, parceled out into the many living individuals of the many species. This piecemeal character is the natural character of temporality, finitude, and sensuousness; here is temporal hunger and temporal satiety.

This temporality is, moreover, recognizable in every body by this, that, even if the living individual escapes becoming prey to another, and everywhere finds occasion to still its hunger and nourish itself, it nonetheless ends with the eater eating itself to death.

When a human being grows old, one notices how the limbs stiffen, and the powers decline; this comes, after all, from the fact that the organs of nourishment are at last worn out and dulled under their own activity.

Thus: both the one who dies of hunger, and the one who becomes the beast of prey's quarry, and the one who dies of old age,
% --- p. 27 ---
all die essentially of the same cause; these are only different sides from which the body's mortality shows itself. Here one hungers temporally, and here one is satiated temporally; in temporal hunger the body proclaims its mortality: mortality is laid within hunger itself.

We are here at a turning point. For if anyone can now prove to us that the soul truly hungers and thirsts after the Eternal, the Personal, just as essentially, just as deeply, as the body hungers after the temporal; then it shall be an easy matter to prove to him that the soul is eternal, disposed toward an eternal life. As the being is, after which there is hunger, so also is the hungering being. As it was then said of the body, that its mortality is laid within its hunger, so it must now be said of the soul, that its immortality is laid within its hunger.

But there is one doubtful point here; this doubtful point cannot be resolved by any general consideration; it is of a personal nature, it is the spiritual character of the hunger. It is here that the half-proof fails.

Many deceive themselves, and reason thus: ``that we feel a need for an eternal life, that is quite certain; but, on the other hand, it is unfortunately very uncertain whether there really is an eternal life.''

For the thinking person it is settled that the matter stands the other way around: ``if it is only quite certain that your soul personally needs an eternal life, then there is quite certainly an eternal life.''

The uncertainty, then, turns upon the need: upon the need and the hunger falls the whole emphasis.

That the human being is essentially spirit, and that the spirit has its hunger and thirst, surely no one will doubt. The general need for instruction, the peculiar need for poetry, art, and science, testify sufficiently to this.

In the human being, spirit is personal. The human being who cannot in any way assert his personality lacks, after all, human independence. Scarcely, then, will anyone contradict us, when we assume without further ado that in every capable human being there is a strong drive and a great need to
% --- p. 28 ---
exercise the personal powers, and in one direction or another to assert his personality.

Personal self-feeling naturally expresses itself as a sense of honor. The sense of honor is therefore strong in every deeper individuality, not least in the younger person, who, without yet being able to assert the ideal claims, burns with desire to come to grips with reality and to test his powers against existence. There is in him a spirit that hungers and thirsts; there is in him a spirit's need for personal satiety.

And what does reality now do, when it catches sight of such a human being, and sees that there is ability in him? It offers him an occasion, it casts him a glance, as though it would say: just come!

Then the dreamer is awakened to battle for the ideal; it is a battle for honor, for honor is the natural ideal of personality. But what is honor without superiority and spiritual power? So the person must then see to setting himself a goal, and, as the saying goes, break himself a path.

The goal thus set has, at once, this double-sided character, that it is at one and the same time the self itself, and yet something entirely different. To contend directly for one's own interest, for one's own satisfaction, and declare it openly, is disreputable, dishonorable, and unsatisfying. To contend in such a way for a cause, so that the cause is advanced and everything personal is forgotten, is, however high-flown it sounds in public declarations, nonetheless still dishonorable, since it is not, after all, the cause, but the person, to whom honor in the proper sense can be applied. Personal striving thereby acquires, as everyone easily perceives, a very ambiguous appearance. From the one side it appears as though the one striving thought not at all of himself, but only of the cause, wanted nothing at all for himself, but everything for the cause; from the other side one sees clearly how the one striving nonetheless, at bottom, wants everything for himself, and, throughout the whole cause, essentially thinks of
% --- p. 29 ---
himself. It is this ambiguity, then, that is so conscientiously preserved in the worldly concept of honor.

Honor demands that I sacrifice myself for the cause; but precisely in that I sacrifice myself for the cause, honor also demands that I, in a higher sense, set the cause up as a means to my self-love. That the latter proposition is seldom or never brought out as strongly as the former proves only that the world is not entirely honest; wherefore everyone who wishes to live only for worldly honor is thereby also compelled to be a little dishonest. For the personality must see to preserving the appearance that most of it is done for the sake of the cause, while at the same time quietly seeing to it that self-love shall not lack its due encouragement. By means of the concept of honor, then, a many-sided and widely ramified equilibrium is brought about among the many worldly causes and the many worldly persons. The greater the cause, the greater the victory (when there is victory); the greater the victory, the greater the honor; the greater the honor, the greater the personal satisfaction.

It is then no wonder that here, in this world of spirits, there is hunger and thirst after honor.

The honor of which we here speak is a being that cannot do without its corresponding outward phenomenon. If the phenomenon fails to appear, then the drive for honor causes torment. Here not only personal merit is required, but recognition of this merit is also required: without recognition, or at least the hope of recognition, honor is nothing. The truly deserving one, who wants honor, but cannot attain recognition, becomes tragically unhappy; the one who, without ability and merit, merely craves the phenomenon of honor, and yet cannot attain the phenomenon, becomes comically unhappy.

It is then no wonder that here, in this world of personalities, there is hunger and thirst after recognition and phenomenon.

When the ambiguity in which the worldly concept of honor maintains itself is cleared up, the contradiction comes into view. The cause and the personal lie in secret strife with one another. To
% --- p. 30 ---
work for a cause, be it a scientific or a political or a civic-industrial cause, is, in other words, to work for the universal. The universal is, in its universality, impersonal. Over against this impersonal, then, stands the individual's private interest.

So long as private interest can, in some measure, be reconciled with the interest of the universal, honor can be preserved and advanced in all its ambiguity. Thus a great merchant can at once increase his capital and, in several respects, serve the public. Thus the Pope could, in his time, at once satisfy his lust for dominion and confer great divine benefits upon ``the holy universal Church.'' The great merchant and the great churchman thus both work for the sake of the cause, and both can, each for his own person, lay claim to much honor --- for the sake of the cause.

When the modern age so rather unreservedly extols the personal passions of ambition and the grand interests of self-love, this advance in shamelessness at the same time points to an advance in sincerity; it is a proof that the ambiguity has already been seen through.

Over against the particular interest, what we here call the cause, or the universal good, is the higher. Self-loving individuals must therefore be compelled, not only by force, but also by a natural sense of honor, into the discipline and order of the universal. This universal is nonetheless still not the highest. The universal is a worldly intermediary, through which personality develops; personality is the goal itself: only the personal can satisfy personally.

To offer the impersonal to one who deeply needs the personal is the same as giving a stone to one who asks for bread. To seek personal satisfaction in the impersonal is a contradiction; the more a human being loses himself in this contradiction, the more the hunger of the spirit grows: here is the insatiability. The one who craves the phenomenon of honor can, in his craving, never be satisfied; the one who craves power and
% --- p. 31 ---
eminence can, in his craving, never be satisfied. It is not an accidental exception that some ambitious and domineering individuals are insatiable; it is a law of the spirit, as certain as any law of nature: the impersonal can never satisfy the personal.

That this insatiability does not show itself to an equal degree everywhere one might otherwise expect it comes partly from the fact that reality sets a barrier for most people: thus far, and no further. So the worm languishes, though it does not die; so the ember smolders, and does not flare up into flame. It comes partly, too, from the fact that many who are thus compelled to limit themselves also become attentive to something more inward, something nobler. So the lover chooses his bride; so the husband loves his wife; so parents love their children; so human beings seek refuge here and there in the arms of personality, take shelter in the quiet harbors, and are saved from the great shipwreck.

But those who do not cast anchor in the grounds of personality, but boldly steer out into the gulf, and bare their breast to the cold impersonal powers, and satisfy themselves with air, are plagued by the fury of hunger. Everywhere craving to satisfy their own self, they toil as slaves under a foreign yoke. They labor for ideas, only for ideas, and call themselves men of ideas; they honor the race, only the race, and call themselves the strong ones of destiny; they look high, and they sit high; they sit, indeed, high at destiny's breast, --- and hunger.

Here is an error of the spirit, one of those which Scripture reckons among ``the strong delusions.'' But there is nonetheless some worldly truth in this semblance; honor is nonetheless a personal ideal; honor, says the old proverb, is still the fairest tree in the forest.

Insatiability shows itself differently when the spiritual, with which personality is to be satisfied, is without further ado degraded to a finite means, and loses the tension of ideality. So it is in the refined nature.

What is the refined nature? It is a nature that demands the very finest enjoyment of the spirit. On this altar of enjoyment
% --- p. 32 ---
nature, art, and poetry must sacrifice their noblest gifts. The most exquisite, the rarest, the most sublime, that which has cost ingenuity its greatest exertions, and talent many years of diligence to produce, is then consumed by the one who enjoys in a few moments, and in all this he enjoys, after all, only his own enjoying self. The refined nature does not enjoy out of natural need; here is a hunger that demands ever stronger stimulants, the longer it is indulged. To enjoy in this way is emptiness and torment. Here is famine; it is a spiritual consumption: the burnt-out personality consumes itself.

The refined nature has, admittedly, not taken shape in the ancient myths; this nature is too reflective for pictorial immediacy. But the passions of the Titans, the unrest of spirits, and the torments of the heaven-storming personality the ancients knew. To this every image in the Greek Tartarus expressly testifies: Ixion on the wheel, the Danaids who never rest, Sisyphus who rolls his stone, and Tantalus, indeed, above all Tantalus, tormented by an eternal hunger and an eternal thirst!

So deeply did the Greek see into the essence of personality. He did not see the hidden kernel; he did not see the holy depth (he lacked the light of revelation); but he saw the bottomless; he saw down into the abyss; here he saw Tantalus, and sensed an eternal hunger and an eternal thirst.

To wish to satisfy the personal spirit with the impersonal is injustice; personality demands its right, spirit demands justice.

\emph{Blessed are they that hunger and thirst after righteousness, for they shall be filled,} --- this is the holy voice from the deep. To hear this voice, and to believe this word, is more than a proof, it is a life, it is life itself, the eternal life.

%% ============================================================
%%  IV. PERSONAL HELP (printed pp. 33--42)
%% ============================================================
\chapter*{IV. Personal Help.}
\addcontentsline{toc}{chapter}{IV. Personal Help}
\label{ch:forel4}
\markboth{IV. Personal Help}{}

There has, then, formed a society of friends, whose aim shall be neither political, aesthetic, scientific, nor industrial, since it is, on the contrary, aimed solely and only at the personal satisfaction of the spirit's personal need. In this society, speeches are indeed also given, but never studied, never premeditated addresses calculated for effect; such addresses the society rejects as vain human handiwork. Instead, anyone to whom the spirit gives it in that very moment may rise and take the floor. One easily sees that this points to the Society of Quakers.

Let us then, for a moment, imagine that we were a kind of Quakers, and thus now gathered in the expectation that one among us, the one in whom the spirit had kindled the inner light, would stand up and deliver a speech. --- ``Will the spirit come? --- is it stirring in anyone? is there, perhaps, at this very moment, someone who will rise?''

Such an expectation has, especially the first time, something very exciting about it; but, if one must wait long, it becomes wearying and tedious. To avoid this tedium, let us then further imagine that the spirit came as in that first excitement, that, then, someone at this very moment stepped forward among us, earnest as a Quaker, inspired as a Quaker,
% --- p. 34 ---
solemnly dry as a Quaker, and began to speak. We fall silent then, for we are, after all, listening to the speech. But, so that this silence too shall not become tedious, it is surely best that, at the very last, we imagine that the speech has already been delivered, that we have already heard it.

What a speech, what an impression! Such an impression is never forgotten, for it is indescribable; and the indescribable I never forget. But the content? Well, a part of the content could unfortunately easily be lost, especially since the speech moved with such speed. The most important part of the content I shall then try to render from memory, while at the same time keeping as close as possible to the speaker's own words:

``My friends, guard yourselves against deceivers! There are found in the world many, many kinds of deceivers, many, many kinds of deceit. There are deceivers who cheat us out of money and goods; there are deceivers who corrupt our good name and reputation; there are deceivers who steal our heart's peace. And yet these deceivers, gross and subtle, are nothing in cunning, in craft and cleverness, compared with the deceiver I shall now name. This cunning, crafty, wily deceiver, where is she to be found? where does she dwell? She is in here --- in ourselves, she is our own sensuous imagination.

How dangerous this deceiver is can be seen when she juggles with a nothing as though it were a something, and beguiles human beings, so that they dread death, sometimes even more than pain and loss and sorrow. This comes from the fact that the deceiver has made them believe that the one who formerly lived is now --- the dead one. So they look at the dead, and say to one another: yes, it is he, still himself, with these pale, stiff features. So they lay the dead one in the coffin, and weep bitterly when they carry him away, as though it were truly he, he himself, whom they were to lay in the ground. So they carry the dead one away; and the solemnity of the funeral procession and the tolling of the grave-bell and the little earth that rattles upon the coffin, all this fills them with dread for the dead one's sake; for the
% --- p. 35 ---
deceiver uses death as a phantom image, that the beguiled ones shall stare down into the darkness of the grave, and not see into the light of the spirit.''

``And this deceiver does not remain sitting by the grave, motionless and silent, --- an image of grief upon the monument; no, no, she moves about in the busy life; she is all activity, she juggles everywhere in the world's turmoil. She steals about within language, she creeps through the word, so that even a truthful man's speech can deceive.''

``So human beings speak of sorrow, how it comes, and how it goes. They speak of the great sorrow and of the lesser sorrows, and measure sorrow by an outward measure. Is sorrow, then, something in itself; can any sorrow grieve itself; is it not, after all, the one who grieves who must himself grieve? But, blinded by the deceiver, human beings look at the sorrow, and forget the one who grieves. By this delusion they even become loud-voiced, and wish, on their own authority, to weigh the sorrow, how heavy it is, without asking the one who grieves.''

``And yet this game is still innocent compared with the deceit that is practiced with our conceptions of the good. The deceiver is obliging enough; she goes about like one who points the way, she points to the good things, the many great, greater, lesser goods that hang as on life's Christmas tree. Here the enchantress lights the many candles: it shines, and it glitters; the children of men look at all that glitter, and are beside themselves in craving and expectation. Yes, beside themselves; here is the right word, here the deceit is at its height, for this being-beside-oneself means precisely the essential evil,''

``The good is a light; but the light is only in the one who sees. When all your eyes close, then nature has no light; the ether then trembles in vain: all is wrapped in an eternal darkness. The good is, after all, the light of the spirit; the light is only in the one who sees: God is a light, in God is no darkness. Therefore, my friends,
% --- p. 36 ---
nothing is good that is not of God; and therefore, my friends, none is good, save one, who is God.''

``And now justice, this attribute in the essence of the spirit! Here the deceiver comes again, and separates the attribute from its essence, and displays the attribute as though it were an independent being, and makes it into an idol. So human beings now speak only of justice, and forget the just one. O, fools, what, after all, is all justice, if we forget the just one? Never has justice procured any human being's right, and never has injustice yet done any man wrong. It is, after all, the just one who does right, and the unjust one from whom all wrong comes. If the just one vanishes, then justice vanishes with him; if the just one comes, then justice indeed comes with him, for in itself alone it cannot exist.''

``But alas, the deceit has so worked that human beings look far more to a general justice than to the one truly just, far more to the good than to the good one. The deceit has so worked that human beings, who writhe under injustice, do not for that reason humble themselves before the one just one, but wring their hands like those who cry out for vengeance, and swear by the justice of heaven. O, you blinded idolaters, why do you not rather swear by the justice of the air!''

``It is the deceiver who has struck the unfortunate with such blindness. They cannot see that all of temporality's distinction between mine and thine, between right and wrong, vanishes, and dissolves into air, when we deny eternity's distinction between the just and the unjust. But the judgment awaits! The hour will come when the spirit shall demand its right, and all sinners tremble before the judge himself, the just one.''

With this, then, the Quaker speech ended.

It is one-sided to seek the good, the right, in what one calls the nature of the matter; the good, taken strictly, is not in the matter, but in the personality. Everything moral is
% --- p. 37 ---
personal; the moral ideas can by no means have a greater validity than can belong to personality itself.

If everything personal is finite, mortal, merely temporal, then everything too that we call the good, the holy, the just, is finite and merely temporal. The distinction between justice and injustice can by no means have any eternal validity, if the distinction between the just and the unjust has no eternal validity. That the unjust can convert, and come to share in the just one's lot, does not abolish the distinction; but if the just and the unjust dissolve in death, then death also dissolves every unconditional distinction between right and wrong.

All moral truth is, in its last ground, personal truth; personal truth is, in its last ground, religious; therefore moral legislation points to the will of God. To fear God and keep his commandments is the fundamental law of human personality.

In the speech itself, this fundamental law is now emphasized with a certain one-sided harshness. The spiritualist speaker has overlooked the natural side of our life, and has therefore also failed to account for the interaction between sensuous existence and personal formation. We will now take the natural conditions into account, in order to make clear to ourselves in what respects the present world offers help for personal development.

If the human being is destined for eternity, then it must indeed be a spiritually strong and indissoluble self that forms itself in the soul. This formation begins naturally: the human being is by nature selfish, just as the animal is. Such a natural selfishness is nothing one need be ashamed of; it is a foundation that conditions something higher. The greatest personalities, even the prophets and the apostles, were originally selfish natures. If we see an Elijah seeking Jehovah amid the flames of fire in the strongly rushing storm; if we hear the sons of Zebedee begging their Master that fire might rain down from heaven upon the ungrateful; if we meet Saul, the raging one, on the road to
% --- p. 38 ---
Damascus: then we cannot doubt that here, in these individualities, there must have begun a powerful natural self.

What, then, does existence do? It helps in a twofold way; it incites the natural self: it flares up, it rages, and foams over; it presses down the natural self, it thrusts it back, so that it must, in inward tension, stir according to its selfhood. Under these conflicting movements the spirit comes.

The spirit, in its naturalness, is, as regards the personal, to be characterized as the moral-practical reason. Reason shall, then, govern, educate, and form the natural self. This education is, according to the principle of nature, imperfect: the natural self is stronger than the natural law of reason, just as the personal is, at bottom, always stronger than the universal. Unable to reach the unconditioned, the morality of reason must confine itself to a finite equilibrium among the powers, an equilibrium that all too much depends upon circumstances.

According to circumstances, human beings then divide into different classes. In some the strong self rules: they swell with courage, they rise in noble pride. In others the weaker self is subdued, they are pressed by circumstances, and bow before superior force. The opposition between the proud and the bowed-down is outward: here are the proud, there the many bowed-down! The opposition is in the temporal.

Under this temporality the spirit becomes worldly: the proud become arrogant; the bowed-down, servile. Here one sees the perishable, the mortal, in personality; it is spirit in spiritlessness, for it is piecework. The spirit of personality is, in its truth, the spirit of God; the spirit of God is known by this, that it masters the piecework.

Only insofar as the spirit of God works in the human being does the outward opposition become transformed into an inward one, and the temporal made subject to the eternal. The spirit begins by bringing down the proud one and raising up the downtrodden. In the pressure
% --- p. 39 ---
under which the proud one feels himself brought low, and in the power by which the downtrodden feels himself raised up, eternity works personally. It begins with the one being brought low, and the other raised up; but it ends with the same one being at once brought low and raised up. In this turning point the eternal self comes to its breakthrough.

In worldliness and piecework it holds: the more self-feeling, the less humility; the greater the humiliation, the less the self-feeling. In religion it holds: the greater the humility, the higher the self-feeling; the higher the self-feeling, the greater the humility. Personal truth, true personality, has, in its revelation as a human being, the highest self-feeling in the truest humility. The same one who says: \emph{All power is given unto me in heaven and on earth,} the same one also says, indeed: \emph{Learn of me, for I am meek and lowly in heart.}

If the human being is destined for eternity, then an exalted goal has been set for us. If it is to begin naturally, then there must be in the human being an upward-striving drive, and there is. As the human being is formed to look upward, and thereby already distinguishes himself from the stooping animal, so too the drive for honor peculiar to the human being is a natural yearning toward the exalted. It is fantasy that helps the drive for honor to find its first alluring ideal, and Joseph is surely not the only youth who has, early on, dreamed that the sun and the moon and eleven stars bowed down before him.

Also in this respect, existence comes to the individual's aid by two mutually opposed effects. It awakens the drive for honor, provokes and incites it; it wounds the sense of honor, and presses with its whole weight, as though it would subdue the drive. It then becomes, once again, reason's task to bring about a certain equality; but the arbitrariness of temporality is stronger than reason.

The piecework, the fragmentation shows itself again: here, indeed, on the one side are the honored, the esteemed and recognized, on the other
% --- p. 40 ---
side the despised, the misjudged and dishonored. The world's judgment furnishes the standard: the greater the honor, the less the dishonor; it is human beings who honor and dishonor one another temporally.

If there is now a human being whose natural drive for honor is so pure that he seeks honor not for the sake of the phenomenon, but in personal truth for the sake of the good, then it cannot fail that existence will most readily press upon such a human being, and let him bear the burden of misjudgment.

What, then, does this mean? --- ``That it is an evil world we live in!'' --- Ah, well; but it also means that it is a good world we live in. The world itself confesses, after all, that it is shameful thus to belittle and misjudge the good one, and yet it misjudges him nonetheless; it approves and disapproves of the misjudgment in the same breath. The world cannot help it; it is set out to act in this way. In this respect it goes with the world as it went with Caiaphas, who prophesied without himself knowing it; in other words: the world does personality a great service, without itself knowing it.

When the human being is deeply grieved by the misjudgment, and shudders at the dishonor; then it is precisely a moment of eternity; then religion comes with a diadem, and the diadem is --- a crown of thorns. When the apostles were first scourged for the word's sake, they thanked God for that honor. It is an honor thus to be dishonored: here eternity has triumphed, personality breathes in God.

If the human being is destined for eternity, then it must have in its nature something that awakens an intimation of blessedness; and the human being has it. It has a drive to give life for life, to surrender its life, in order to win its life, a double life in personality. This drive is natural; for the mother's love for the child stands as close to natural sympathy as possible.

``To love and to be loved, yes, that is life, and that is blessedness!'' --- Here comes the poet; but the poet we shall, in this
% --- p. 41 ---
context, thank with the assurance that we have already heard him. All the blessedness that is in natural love is, however, essentially only a symbol, a sign, a forewarning. In these natural relations blessedness stands against unblessedness, as happiness against unhappiness. The one who is happy to the very verge of blessedness can, in a moment, become unhappy to the very verge of unblessedness. Existence puts these extremes together in the best sense; it is a combination that helps forward the personal.

To love naturally is natural enough, but it is not eternal. That Niobe loved her children was natural; these children were, after all, her life, her joy, her pride. But in this love the eternal was not yet present. The eternal gods grew wrathful with Niobe: Apollo and Artemis bent the bow; their arrows slew these children; and Niobe, who had just now been so blessed, so proud, became now so unblessed, so defiant and hardened, that she was turned to stone.

Here is an ingenious image of merely natural love, and of the powers of fate, which loose their arrows to annihilate its blessedness. Here is an image of the natural heart; in its blessedness it is as rich, as warm, as proud as Niobe's heart; in its unblessedness it is as defiant, dead, and cold as Niobe's heart. Niobe was made eternal in her petrifaction; this petrified eternity is the impersonal.

\emph{Thou shalt love the Lord thy God above all things, and thy neighbor as thyself.}

This commandment of the spirit is stern for flesh and blood, it is so cold for the warm heart; but it is nonetheless not so cold as Niobe's heart. Had Niobe known and grasped this commandment, then she, though bowed down like Job, would nonetheless not have been turned to stone. The cold commandment, which at first glance looks like death, would then have become for her a commandment of life, proclaiming blessedness.
% --- p. 42 ---
``Why, then, should these children have been given me, why should I then have pressed them to my breast, and felt this joy, this blessedness, if they were again to be so cruelly torn from me!''

Thus a Niobe might well complain; but if she complained thus before God, then religion would bring her relief, and set the riddle for her in a different light.

``They were to be given you, O Niobe, for blessedness's sake; they were to be taken from you, O Niobe, for blessedness's sake. You were to rejoice so deeply in these children, that you might afterward learn what it is to grieve. You were to have them, that they might be taken from you, that you might now feel what it is to lack them, and long for them.''

``You grieve deeply, O Niobe, and well for you, that you grieve deeply; for by your sorrow you shall now know how earnestly the joy was meant. You long deeply, so deeply, and well for you, that you long so deeply; for by your longing you shall know how earnestly the blessedness was meant. So shall your sorrow and longing nourish your hope for you, and this hope shall thrive, sound and strong, as one that draws its nourishment from your longing. Now you shall love God, in order rightly to love those you lack; for your children are, after all, with God. Now you shall love your neighbor as yourself; for every child you meet will, after all, remind you of those you have with God. So shall you get them back, each and every one: in this reunion is your life, your blessedness; so shall you lack no one, nothing, once you have learned the great commandment. But not before, O Niobe, no, not before you have learned the great commandment; for blessedness is that God is love.''

Personal help would be easy to understand, if we human beings only knew the languages. But one thing is the earthly language, another the heavenly language: what in the earthly language is interpreted as cruelty can, in the heavenly language, be interpreted as blessedness.

%% ============================================================
%%  V. COMPANY IN SOLITUDE: A PERSONAL TASK (printed pp. 43--54)
%% ============================================================
\chapter*{V. Company in Solitude: a Personal Task.}
\addcontentsline{toc}{chapter}{V. Company in Solitude: a Personal Task}
\label{ch:forel5}
\markboth{V. Company in Solitude: a Personal Task}{}

One sometimes tosses out, for social entertainment, the question: what can form a human being most, solitude or company? Since this question is surely just as profound as it is sociable, we may, for a change, turn the task around, and ask: what can torment a human being most, company or solitude?

In order to discuss this matter thoroughly, we dare not confine ourselves to thinking of a little solitude and a little company; we must see the effects on a large scale; we must first think of an uninterrupted, continuous solitude, an uninterrupted, continuous sociability. Only by thus taking it on a large scale may we hope that our inquiry will at last succeed in separating sociability from solitude, and penetrate into the peculiar essence of each.

As regards, then, first of all, a continuous solitude, it must surely be extremely painful, particularly for the reason that it is so closely bound up with silence.

How painful a continuous silence must be can be seen in the monks who have taken the strict vow of silence. When these tormented ones now and then get permission to speak together, the need for communication is so powerful that they
% --- p. 44 ---
together burst out in one great common stream of words; the many speeches break forth in one instant like jets from many hoses, and the total effect becomes so overwhelming that it recalls a sluice-gate more than a conversation.

If this is how it stands, what wonder, then, that many human beings shun solitude, because they are burdened by the silence, and feel a need for the living word.

But, on the other hand, it is also questionable enough with sociability, the incessant, ever-lively sociability, among other things for sleep's sake. To fall asleep in company, even if it were large company, may surely by no means be considered impossible, but to sleep sociably is a logical impossibility, since it contradicts the very concept of sleep.

If it is dangerous to swear an eternal silence, then it is, in this world, life-threatening to swear an eternal sociability; for some sleep a human being must, after all, have.

Just as difficult is the task, when one sets oneself the serious question: what, then, might be the most distinguished, to live very solitarily, or very sociably?

To live in company, and find oneself content in all sociability, there is in that, then, nothing distinguished; that, after all, every everyday person does. No, it then seems, rather, more distinguished to remove oneself from the sociable, and choose oneself an idealistic standpoint outside company.

On such a somewhat solitary, idealistic standpoint the emancipated one may then move about rather unconstrainedly, amuse himself with existence, play ball with the civic realities: ``a good trade, a good livelihood, a good brother-in-law-ship, a good friend,'' and so on.

When this game is practiced for some time with a little talent, it will always have its effect; anyone can easily see that a personality who is raised above such narrow-mindedness cannot possibly be any everyday person. The satirically sneering is, after all, precisely a sign of superiority; the superior one is already somewhat outside, or rather: beyond, company,
% --- p. 45 ---
thus somewhat solitary, and can, for that very reason, already demand no small recognition from company.

This, however, is not enough. The satirical nature is already fairly widespread. The solitary one could easily expose himself to falling into even very common company, for satire in our age is by no means distinguished.

Here a step forward must be taken, here a transformation must occur. The solitary one must see to cultivating a secret pain, a quiet melancholy; he must acquire a tinge of the melancholic, that is to say: the satirically melancholic; he must set his mind on the inexplicable, the incomprehensible, the ungraspable. Yes, the ungraspable! here is the high point.

The ungraspable one is always in solitude. What he thinks, what he loves, what he suffers, what he feels, what he wills, and whither he strives, no one feels, no one grasps, no one in the whole of company. This superiority in the ungraspable is the highest distinction; for keeping the sociable at a due distance there is no method that can be recommended like this one.

To withdraw directly from company and take a solemn leave, for instance in the newspaper, is not nearly so certain a method. That it caused a stir could hardly do harm; for thereby it would, after all, become known to all who it was who most wished to be unnoticed by all. But the danger shows itself from another side. The community might, namely, miss the solitary one, long for ``the unnoticed one,'' take a fancy to imitating the rare one, and thereby at last make it impossible for him to preserve his rarity.

It could then go here as it went, in its time, with the first hermits in the Egyptian desert. When one of the holy men let himself be seen in Alexandria, the whole population became so seized with awe that half of them ran after him out into the desert, and thus founded a new society, a society of solitaries: and with that it was then over with the rareness of solitude.
% --- p. 46 ---
No, in this respect it is far finer with the ungraspable. The ungraspable one can move about in the midst of the great world, and still be the solitary one. He may, for that matter, gladly accept an invitation to dinner every day, and appear at the theater every evening; if only he does not forget to bring along the melancholic-cheerful, the friendly-bitter, the ungraspable, he can always be certain of his rarity; one notices at once that here is a stranger; one feels at once that here something apart is concealed; anyone in the company who cares to use his opera-glass can plainly see that this one alone feels himself --- apart, so alone.

By these considerations I believe I have then proved that solitude is something very distinguished.

But, on the other hand, when I then look again at company, I see, nonetheless, sociability in the whole of its grandiose indispensability. What would all my distinction avail me, if there were no one here who could perceive how distinguished I am? what would it avail, if I became ever so eminent, if there were no one here to notice my eminence? what would it avail me, that I became distinguished, decorated, and got a title, were it even the title of the Solitary One, if there were no one here to address me by it, no one who knew my title, no one who could so much as write to ``S. T. the Solitary One.''

No, then I praise the multitude; the more there are here, the better; but it is, after all, precisely the multitude that constitutes the company, thus: the greater the company, the greater the distinction for the distinguished one. And --- let us now be honest --- what is it, then, that makes we rational human beings always so gladly go into company, and give parties, large parties? It comes, after all, from this, that we love the eminent, and so gladly wish to see the eminent. If there were nothing eminent about it, then it would surely not be worth the trouble to give such large parties.

With this I believe I have then established that there exists a sociability that is just as distinguished as solitude. Having
% --- p. 47 ---
brought the discussion to this point, I must here break off, recommending this important matter for further sociable discussion.

If we now take up the task from the personal side, and in earnest, then this law holds: that no human being, whether he be in the desert or in the city, may ever know himself to be entirely alone.

\emph{It is not good that man should be alone; I will make him a help, that may be with him.}

From this beginning, a beginning the spirit makes with the natural, proceeds the family, and from the family proceeds, in turn, human society in every direction and every form.

It is not good that man should be alone; for also in the spiritual it holds that the one human being begets, nourishes, raises, and forms the other: the personal is born of the personal.

In its natural beginning, the human being is only an individual. The individual must be raised, that is to say: disciplined, taught, and formed, so that it, according to its original endowment, can appropriate the universally human.

Merely humane formation can, however, never work up the individual and the universal into complete unity, and therefore cannot lead to true personality either. Sometimes the individual gains the upper hand; but, on the whole, universality asserts its harsh dominion. So it was among the ancients.

According to the view of the Greeks and Romans, persons were properly to be regarded as individuals, and these individuals had no rights over against the state as state. It was not only the slaves who were deprived of human rights; even the free, eminent men were, in idea, the community's slaves: it had not yet become earnest with the personal.

That the gods, these aesthetic personifications, could not liberate human beings is self-evident. The essence of the state was the true divinity; the true god was, then, the god of the community. Only insofar as individuals believed in the state did they believe in the gods: such a faith may then rightly be
% --- p. 48 ---
called a communal faith. The individual citizen was not permitted to have his own conscience, his own faith; the universal, common conscience resided in the state: here faith in the company was at the same time a faith in the company.

It is well known that Socrates was condemned to death because he, as it is said, ``introduced new gods.'' Socrates asserted the personal; the ethical thinker broke with the communal faith, and therefore had to die. According to Roman usage one might say that Socrates was condemned because he had made himself guilty of superstition, for according to the Roman concept every faith that deviates from what is prescribed by the state is a condemnable superstition.

On the natural ground of the idea of reason, classical antiquity is essentially right; so long as individuals are only individuals, the state is the highest, and so long as the state is the highest, faith is essentially a communal faith.

On the plane of revelation the relation is reversed. In the perfect revelation personality steps forth, the One, ``\emph{the single individual},'' over against the company, and says, not to the crowd as crowd, but to every single one in the crowd: ``\emph{Come to me; I am the Truth.}
% [printed: no closing quotation mark after "Truth." -- printer's defect
% (unclosed quote), transcribed as printed; accounts for the odd ``/''
% count in this fragment, mirroring the same defect logged in
% transcription.tex at this site.]

That this One --- ``\emph{the single individual},'' thus presents himself as the one who gathers the many by calling each single one personally to himself, is what abolishes communal faith in principle. To recognize this is an inescapable condition for recognizing personal truth.\footnote{That the thinker who, in our time, took aim, and struck ``the single individual'' categorically, struck --- the center, is not a flattering concession to be made to an individual; it is a personal concession that must be made to the Truth. If one wishes, out of an individual predilection for this and that, to steal past ``the single individual,'' and triumph in the company; then one may tinker with ``the Reforms'' as much as one likes: one does not thereby reach the community of the spirit, for one does not reach the personal.}

When true personality itself wished to reveal itself, he did not live as an isolated individual, as a hermit in the
% --- p. 49 ---
desert. On the contrary, wherever he appeared, human beings gathered; wherever he showed himself, there was usually a large company. So he once spoke to an assembly of listeners; but the speech was not to the company's liking: it was --- a hard saying.

Wherein, then, lay the hardness? It lay precisely in the great emphasis he placed on the personal appropriation of his personal life. In order to remove every notion of a merely sociable universal spirit; in order to ward off the misunderstanding that it was still a matter of a cause, a common cause, standing above the personal, he lays divine weight upon his human life, his life in flesh and blood, and says to the assembly that everyone who would have the eternal life must eat his flesh and drink his blood.

At this speech --- it was, after all, also ``a hard saying'' --- (one had not yet gotten around to saying anything solemnly general about it) the crowd was offended; it found that it was too personal, that it descended to personalities, and so it went its way.

The Master stood alone with the Twelve. He looks at them, and asks: \emph{Will ye not also go away?} Then there is one who answers: \emph{Lord, to whom shall we go; thou hast the words of .eternal life.}
% [printed: stray ink mark between "the" and "eternal" -- press defect,
% transcribed as printed, mirroring the same defect logged in
% transcription.tex at this site.]

At this answer, each single one among the disciples fixes his eyes upon the Master. They do not look first at one another, as though they would count one another; but first at the Master, and only then at one another. Had they begun by counting one another or consulting one another, then the answer would surely have had to sound thus: ``Lord, the large company is indeed gone; but we are still Twelve left; that is still a small company. We Twelve are agreed to remain; our number can, in time, be increased, and expand again into a larger company; for, of course, if we cannot first form a company, then we cannot believe at all.''
% --- p. 50 ---
Had such a communal faith hovered before the disciple's mind, he could not possibly have afterward said: ``Lord, though all others forsake thee, yet will I not forsake thee.''

By personal truth, communal faith is abolished. Thereby company is not violated; it has, on the contrary, received its true significance, for within the temporal company an eternal community has been founded.

The impersonal universal spirit, which otherwise everywhere rules in human society, must now, in faith, subordinate itself to the spirit of personality, which liberates the enslaved individuals. The personal communal spirit is therefore also called the spirit of love; the community it founds is a holy community; upon the word holy.falls the emphasis on the personal.
% [printed: "holy" and "falls" run together with no space and a stray ink
% mark between them -- press defect, transcribed as printed, mirroring the
% same defect logged in transcription.tex at this site.]

The true difference between temporal company and the personal community can be made clear through solitude. In a temporal respect, company and solitude stand in outward opposition: the more sociable, the less solitary; here again is the fragmentation. This fragmentation must vanish before the personal. The spirit of personality places company within solitude, solitude within company; in this union consists the holy community.

The one who personally appropriates personal truth is in solitude, insofar as he is with the truth in inwardness; nevertheless he is not, then, alone; he is, after all, in communion with the truth, in a personal relation to true personality. If now many are assembled, then the assembly, in its phenomenon, presents itself as company; but since each one in this company is as a solitary, the company is, in spirit, transformed into a community.

To be sociable is to be in distraction; to be gathered in the community of the spirit is to be in devotion.

When one speaks to an assembly as to a multiplicity of individuals, either in such a way that the speech expressly emphasizes the particular interests of the different ones, or
% --- p. 51 ---
in such a way that a general cause is maintained, which interests all, whether in pure generality, or only partially from a certain side: then one speaks to a company.

When, on the other hand, one speaks to an assembly in such a way that each single one feels his personal inwardness addressed so unconditionally that, in self-absorption, he forgets the multitude of the assembly; then the many individuals are essentially gathered into one, then one speaks, by the power of the spirit of personality, to the community.

Everything here depends upon insight into the personal. There may be a thousand present, and the preacher of the word may still speak to the community; there may be two or three gathered, and the religious speech, which misses the personal, be nonetheless as distracting as in a worldly company.

As the aim is, so too must the speech be. The devout speech has, as its aim, edification in faith.

How, then, do we come to faith in personal truth? An old Catholic has said: ``I would not believe, if the Church did not believe.'' That is a church-faith that culminates in the papacy. The papal communal faith is, under all its disguises, recognizable by this, that it always confuses the temporal conditions with the eternal unconditioned.

Among the temporal conditions belongs the preaching of the word: how should we believe, if we had not heard? how should we hear, if there were no preaching? But it is one thing to grant this, another thing not to be willing to believe, if the Church did not believe.

The Catholic can only believe in Christ insofar as he believes in the Church; the Church is a fixed, infallible intermediary that places itself between him and personal truth. The Catholic relates to the Church roughly as the Greek in pagan-classical antiquity related to the state; that is to say: in his faith in the Church he has only the Church's faith, and the Church has his conscience. The Catholic is, then, right, when he says: I would not believe, if the Church did not believe.
% --- p. 52 ---
But the Protestant is not right, when he says the same thing. The Protestant believes in the Church only insofar as he believes in Christ. The Protestant cannot let the Church believe on his behalf, he must, as the single individual, himself have faith in the inward appropriation of the truth: here the spirit is personal.

When the visible Church deviates from the truth, and forsakes Christ, then the church-believing Catholic must necessarily also deviate from the truth, and forsake Christ. The Catholic is. in this respect a genuine. company-man: he comes with the company and goes with the company.
% [printed: stray ink dots after "is" and after "genuine" -- press defects,
% transcribed as printed, mirroring the same defects logged in
% transcription.tex at this site.]

The Protestant, by contrast, has broken with communal faith; he is like a solitary, and precisely for that reason in the community of the spirit; he is, in personal inwardness, the single individual. Company, with its tradition, is for him a condition, but not the unconditioned. Should company, as company, then once take it into its head to change its view: the single individual does not for that reason change his view; he remains with the truth, he repeats the disciple's words: ``Lord, though all the others forsake thee, yet will I not forsake thee.''

If we look at the practical exercise of the good, we meet the same union of company and solitude. ``Thou shalt love God above all things, and thy neighbor as thyself.''

By the word ``neighbor,'' a word that in our language does not even have a plural, we are again referred to the relation between the single individual and the single individual.

All human beings in their totality no one can love; such a love is a sentimental fantasy. To love merely one's own, one's friends and relations, is heathenism. To love the universal good more highly than individual human beings is a worldly love of ideas, which coexists quite well with ambition and lust for dominion.

The word ``neighbor'' removes every egoistic consideration. It secures us against confusing the real human being with one of humanity's ideas; it secures us against showing partiality of persons, against asking whether friend or foe, but only about
% --- p. 53 ---
the human being; it secures, finally, that we, instead of falling in love with a certain multitude, a certain class, a certain party, each time have to do with one single, particular human being. By the neighbor's relation to the neighbor, the personal relation between the many individuals is expressly designated.

Love of the neighbor is subordinated to love of God, as the second great commandment is subordinated to the first.

If God did not himself, so to speak, step in between human being and human being, then individuals would relate to one another only as individuals, and the many individuals form one large company. But, in that God himself steps in between each human being and his neighbor, solitude comes, and transforms company into community.

When human beings deal directly with one another, they can easily, as the saying goes, encroach upon one another, in friendship as well as in enmity. Under such conditions the one human being will, only too naturally, direct his conduct according to the other human being's behavior. To love one's friend and hate one's enemy is just as sociable as it is heathen.

When God steps in between, those closely bound are separated as though by a whole eternity. The one no longer has directly to do with the other; for each of them has to do with God. For everything I think and do toward my neighbor, I shall be accountable to God; by virtue of this accounting, each single one becomes again a solitary.

But the separating is at the same time the connecting. If the one human being's disposition is to change because the other human being changes, then everything is changeable. Only through the relation to God does the relation between human being and human being acquire significance for eternity. Only those who are separated by eternity can, in truth, be united for eternity.

Merely sociable connections lack personal consecration; here the ground of eternity is lacking. The bond breaks, the company dissolves, the individuals are separated from one another: that is the mortal, the perishable.
% --- p. 54 ---
But through the relation to God, the connection, not only between husband and wife, but between human being and human being, is consecrated to the eternal personality.

I was naked, and ye clothed me; I was sick, and ye visited me; I was in prison, and ye came unto me. --- Inasmuch as ye have done it unto one of the least of these my brethren, ye have done it unto me.

So the individuals are separated by means of the truth, and remain, through the truth, the solitary ones; so the individuals are united by means of the truth, and live, in the community of the spirit, an eternal life.

%% ============================================================
%%  VI. GREAT MEN: PERSONAL SUPERIORITY (printed pp. 55--67)
%% ============================================================
\chapter*{VI. Great Men: Personal Superiority.}
\addcontentsline{toc}{chapter}{VI. Great Men: Personal Superiority}
\label{ch:forel6}
\markboth{VI. Great Men: Personal Superiority}{}

For personal development there is surely nothing more beneficial than to come into contact with a grand personality, and rightly feel what it means to bow before the superior one.

I once spoke with a good friend, in all confidence, about this point; the good friend granted me right, but also drew attention to how difficult it really was to discover true superiority, let alone come so near to it that one could truly receive the beneficial impression.

This remark led me to make some reflections upon human greatness and its relation to the ideal. ``We must, after all, know,'' I continued, ``what kind of personal superiority is being spoken of. If what is meant here is primarily the purely human superiority, as this, according to the conditions of the age, can take on a physiognomy in public life, then our fatherland indeed has a relatively considerable number of great men, and what we do not have among ourselves we may perhaps find abroad. That the ideal is unattainable goes without saying.''

``But do you not see, then,'' he said, ``that, when there is so wide a distance between the ideal and the real,
% --- p. 56 ---
it will indeed follow from this that the ideal impression cannot be real, and the real not ideal?''

``Do you mean, then,'' I objected, ``to maintain outright by this the impossibility of getting an ideal impression of a human being's personality, because that human being is not the ideal itself?''

``If I am to explain to you my view of this difficult point,'' he answered, ``then you must allow me to take a roundabout way, and, in order to take the personal as personally as possible, to relate to you some features of my earlier life. What I have to relate in this respect is surely by no means flattering to my personality; but to err is human, and if my errors could perhaps help you to a clearer insight into the relation between the ideal and the real, I shall not regret having touched upon a good deal that I otherwise do not wish to speak of.''

I did not become, at this beginning, exactly curious (for I knew very well that he would not relate anything calculated to satisfy curiosity), but, since it interested me to a high degree to have the disputed relation illuminated --- ``as personally as possible,'' I asked him to sit down with me, and let me profit from what he himself had gone through and experienced.

After a moment's reflection, during which a certain humor seemed to contend with some diffidence, he then began his tale, more or less as follows:

``Ever since my earliest youth I have been enthusiastic about great men, and dreamed of personal superiority. My ambition awoke early. `To distinguish oneself,' `to surpass one's equals,' `to become something,' these notions were deeply impressed upon me in my boyhood years, for they made up the principal ingredients of the paternal exhortation that spurred my diligence.''
% NB: a small superscript mark (resembling a tiny "e") appears directly
% before the opening quote below, in the left margin of transcription.tex at
% this site. Not a recognisable letter or footnote sign (this book's only
% confirmed footnote, p.49, uses "*)"); no footnote text appears anywhere on
% this page. Treated as a scan/page defect (ink speck or foxing) and not
% transcribed as a character, mirroring the transcription.
``Scarcely 20 years old, I had already finished my academic studies, and passed every examination with distinction. My father rejoiced to such a degree over this
% --- p. 57 ---
happy outcome that he, although such generosity was surely well beyond his means, granted me a considerable sum for a journey abroad.''

``I had already heard so much about public life, about the great movements in public life, and about great men, men whom the genius of liberty had endowed with such personal superiority that they could, by the mere word, govern and guide the otherwise unmanageable masses according to the highest direction of the idea of liberty; it will therefore not surprise you, when you consider my youthful levity, that I, with such liberty and mobility before my eyes, then also traveled with such speed that, in a single leap, I sprang over all the monarchies, and landed straight in the Republic.''

``Just as I arrived in the Republic (I arrived at the right moment), the whole Republic was illuminated: garlands, banners, crowds of people in every street! `Here you meet superiority, personal superiority,' I said to myself, and bared my head, for I became so enthusiastic.''

``Far off in one of the streets music sounded; everyone ran in that direction: I ran along, as best I could. A man who was behind me then also ran along, and nearly ran me down. He had, in the course of running, just thrown off his slippers, and was shooting along at a tremendous pace, when we collided, and I happened to grab him by his blouse.''

``Embittered at this stoppage, the man raised his arm: I was on the very point of receiving an impression of his personal superiority, when my purse succeeded in standing bail for me. I now begged the person, good-naturedly, by no means to strike me, nor to strike away my hand, but generously to accept my purse and be my guide.''

``By this equally agile and knowledgeable man's hand I was then led to the goal of my wishes, the object of my journey, and reached the street along which the procession of people was moving up toward the great square. We came to a halt, my guide and I, near the square, just as the procession approached. In
% --- p. 58 ---
this position I could then, undisturbed, survey the whole, and enjoy the grand impression.''

``At the head of the people's leader, that is to say: somewhat ahead of the drum major, walked a man, a noble, imposing figure, the foremost man of the age, of the Republic, of the century, up the street in quiet dignity --- with his hat in his hand.''

``That the personal superiority that came to meet me in this figure truly was --- the foremost man of the age, of the Republic, of the century, I could not possibly doubt, since thousand upon thousand throats proclaimed it from every side; and every time the extraordinary one was thus greeted from every side, the drum major struck up, and the people's rejoicing mingled with the music.''

``I felt the grandeur, the uplift. Meanwhile, however, once I came a little to my senses, there was a thing which, in the midst of all this greatness, struck me, I might well say offended me.''

``All the male persons in the Republic, the well-dressed as well as the ragamuffins, namely kept their hats on, even as they rejoiced over the personal superiority. Not only those who moved about in the streets, but even the gentlemen who filled the windows showed themselves with their hats on; indeed, I even expressly saw a man strutting in the middle of a parlor among the ladies go about there with his hat on.''

``As this now vexed me, I said to my guide: `It is really a shame that personal superiority here should be the only one without his hat on. Listen, should not the two of us agree to take our hats off for the great man who is now passing right by us?'''

``The guide, who was just catching his breath after a tremendous shout, when I got these words out, looked at me with a face in which I could plainly read his personal superiority, and said, as he slapped me on the shoulder with condescending familiarity: `Innocent parvenu, you have surely never been in the
% --- p. 59 ---
Republic before, since you do not know that the man we exalt so highly ought to take off his hat to public opinion.'''

``That my guide was indeed well versed in the Republic was confirmed most beautifully when the procession stopped at the flower-adorned tribune on the square.''

``The foremost man of the age, of the Republic, of the century, mounted the tribune, and lifted up his voice in a powerful speech. The whole speech I could not indeed follow (it went too deep for me in politics), but I did gather the most important part, and that was that the speaker --- loved the Republic, sacrificed himself for the people, and bowed to public opinion.''

``Imagine my admiration! for every time the speaker on the tribune loved the Republic, sacrificed himself for the people, and bowed to public opinion, the drum major struck up, and the people's rejoicing drowned out the music.''

``It was toward evening; I was weary in soul and body. The guide led me to a hotel; here I was then to lodge, and take some refreshment. I could enjoy nothing; I was so full, so full of gratitude. I thanked providence, I thanked my guide, I thanked the man for the Republic, indeed, I even felt that I ought to thank the drum major; more I could not do: I hung slumped in my chair, I had a fever, and began to rave.''

``Just as the doctor stepped in, I had a lucid moment. I assured him that the whole thing was a slight indisposition, which would soon pass; I most emphatically charged my guide that he must by all means come precisely the next morning, since I felt an irresistible urge to pay my respects to the foremost man of the age, of the Republic, of the century.''

``Of what followed now, I have only dim, confused notions. When I at last properly came to myself, I lay with a bandage around my arm, and the doctor sat by my bed.''
% NB printer's defect (logged, not corrected, mirroring transcription.tex at
% this site): this paragraph has no opening `` although it plainly
% continues the same first-person quotation as the paragraphs before and
% after it -- only the closing '' is printed.
In the doctor's calm, benevolent smile I saw not only something of a personal superiority, but also an assurance that my illness was not dangerous.''
% --- p. 60 ---
``I have surely been raving, and talked some foolish nonsense,'' I remarked, and looked questioningly at the doctor.''

``Yes indeed, you have been raving,'' he answered, ``but do be quite calm. You have not, upon my word, betrayed any state secret. All that you have said has so public a character that you would surely already have made your fortune publicly, if it had only been heard publicly; for you have, indeed, for three whole days done nothing but incessantly love the Republic, sacrifice yourself for the people, and bow to public opinion.''

``Your political understanding must undeniably be very firm, since it, as I shall gladly attest, is just as sensible with a fever as without one.''

``You are looking at your arm. Yes, I have indeed bled you; but I did that only out of concern for your bodily welfare, and then for fantasy's sake.''

``Take care with your fantasy, that is my advice to you; for it is, to be sure, not free from sometimes producing what you call some foolish nonsense. On the short political career you have covered here under my supervision, fantasy has, several times, run away with you rather alarmingly. Just think, that you, only last evening, and that even in full earnest, could put to a vote a proposal so impractical as this, that the men of the people ought to rise, and vote unanimously for the law that the whole Republic should take off its hat --- to the drum major! Now you see for yourself that I have had reason to bleed you.''

``As a convalescent I had to keep indoors for a few days. I had a notion that something extraordinary was afoot in the city; but what it really was, I could not find out: the waiter was silent, and the doctor had confiscated the newspapers.''

``When the doctor at last found me well enough that I dared go out again, my guide met me with a knowing face (in which I then again discovered glimpses of personal superiority), and
% --- p. 61 ---
offered that, if it were now convenient for me, he would gladly show me the foremost man of the age, of the Republic, of the century.''

``We now went round by another street. There was great commotion everywhere, just as before. To reach the place my guide indicated, we had to cross the river bridge; but here there was almost no getting through. People were crowding together; there was a man who had drowned himself in the river.''

``Like an arrow the guide shot into the crush, and came back with information. --- `Didn't I think as much,' he murmured, --- `the poor man!' --- `Which man?' I asked. --- `Why, the man you absolutely wanted to take your hat off for the other day,' he answered, and pressed down his hat with a superiority meant to prove to me that his hat sat firmly on his head.''

--- ``But that was, after all, the foremost man of the age, of the century . . . .''

``Ah, nonsense! no, I suspected as much; I really said it quite quietly to myself, when we stood together on the square the other day, that he probably wasn't the right man after all.' `But listen now! can you hear?' --- It was rejoicing again --- `there, surely, they have the right man, come, let us hurry!'''

``At the bridge people were occupied with the dead man. `So he did, after all, in a manner, sacrifice himself for the people,' remarked some. `He perhaps did, at least, sincerely love the Republic,' others observed. `He surely at least drowned himself sincerely for public opinion!' cried a young person who wished to be witty.''

``The dead man was carried away; I cast, in passing, a pained glance at that bared head; the hat floated on the water near the planks of the bridge.''
% NB doubtful reading in transcription.tex, carried over here: the Danish
% has "Admindelse," clearly printed as such (verified at high zoom, not a
% misreading of long-s/f), not a standard Danish word -- almost certainly a
% printer's slip for "Aamindelse" (archaic for "memento, keepsake"), which
% fits the sentence and the satire (the hat as the Republic's lawful
% "memento"). Translated as "memento," matching the intended sense.
``That's a good new hat, that's a memento which may be regarded as the Republic's lawful property; we must really see about getting it saved,'' said the young person, and bent down, to grab the hat. He did not succeed; he only gave it a push outward: the new hat floated a little way down the river, and sank in the current.''
% --- p. 62 ---
``Here you have the first fragment of my life's experience. It is a peculiar business, to travel abroad, to discover personal superiority.''

``A few hours later I sat in a traveling carriage, whose team faced toward Denmark. My heart pounded, as though I were pursued by wild demons; I yearned toward the fatherland, as a madman, who senses his condition, yearns toward sound reason. I had a presentiment that at home I should meet quiet greatness.''

With these words the narrator rose, and left me.

For personal development the impression of a superior personality is surely most beneficial; but the narrator is right, when he draws attention to the fact that it is no simple matter with personal superiority.

The man who is to be able to guide a movement in public life must, above all, be an independent, a self-reliant man. Given the requisite spiritual gifts, particular circumstances, such as rank and fortune, may indeed furnish favorable conditions for the feeling of independence; but these are, after all, only conditions; the unconditioned is in the personal. Only in true personality is there true independence,
% NB printer's punctuation oddity (logged, not corrected, mirroring
% transcription.tex): the paragraph ends here with a comma, and the next
% sentence begins a new paragraph.
To political independence belongs, after all, not merely a firm view of the goal that ought to be striven for, and a sure eye for the means that can and ought to be employed, but there belongs to it also, and this preeminently, that the one who governs is, from the outset, at one with himself as to how much he can and will venture, with what degree of self-sacrifice he can and dare work for the idea.

These are strong, tense oppositions that must be united and mastered within the ruler's inward being, if his dominion is to be said to be grounded in personal superiority.

Without a living interaction with the popular will as it expresses itself through the given institutions, a rational governance of the state is unthinkable. The one who governs must have a deep
% --- p. 63 ---
conviction of this; for before the idea he is only justified insofar as he dares regard himself as an organ for the essential universal will.

But the universal will expresses itself precisely in the interaction of the many particular individual wills: here opinion stands against opinion, claim against claim, interest against interest, party against party, passion against passion. In the finite phenomenon the universal will announces itself as the will of the crowd.

The opposition between the popular will, grounded in the nature of things, and the passions in the will of the crowd, is a tense opposition. The one who governs thus shows himself, in idea, as the one who, of all, must be the most dependent and, at the same time, also the most independent. To unite dependence upon the true universal will with independence from the passions of the day and the shifting mood of the crowd is an extraordinary task, whose solution is possible only for the superior personality.

To let oneself be raised up by the mass, to take the mood of the crowd as an expression of the universal will, and then to go at the head --- ``somewhat ahead of the drum major,'' is to deny oneself authority in the very moment one seems to appropriate it. It is no accidental mishap that toppled the man in our related tale; it is a categorical contradiction: he toppled himself, in the very moment he, as leader, began his procession --- ``with his hat in his hand.''

This political comedy, which is repeated often enough (though not always in such glaring colors), must necessarily contribute to undermining civic order: it brings those who govern into contempt, it bewilders the crowd, and demoralizes the people.

But the danger from the opposite side is just as threatening. To think meanly of individuals and despise the judgment of the crowd is surely tempting for the one who has power, and wishes only to enjoy it. But when the idea of the state is severed from that which
% --- p. 64 ---
ferments in the masses and wills to work its way forward in individuals, it becomes a fixed idea of power-holding, which lacks human justification. Such a fixed idea may indeed be flattering enough for one lusting after dominion: a reckless despot, a fortunate tyrant may well make use of it; but thereby the true task is not solved. The despotism that subjects the popular will to fate itself becomes, in turn, a tool of fate, again crushed by fate.

If we rightly consider the tense opposition of dependence and independence, an opposition that presses and weighs upon the one who governs at every step he takes: what power of inwardness, what personal superiority must he not then need, in order, even only to some degree, to solve his task!

One finds it natural that the crowd fears the one in power, and that he, in turn, fears losing power; but, understood personally, the one in power has the greatest reason to fear for himself: if there is anyone whom power can deeply alarm, it must surely be the one who holds it.

If the idea of the state were the highest, unconditioned idea, and the one who governs absolutely assured that his individual conception of the state's goal was infallible, then the choice and use of means would simply depend upon finite cunning, courage, and prudence; we would then, without further ado, have to admire Machiavelli's ``Prince.''

But the idea of the state is not the highest idea, just as truly as the personal is higher than the universal; still less can any mortal maintain that his individual conception of the state's goal should be so infallible that he could justify, in order to attain this goal, using any means whatsoever.

To work for a goal, as for one's life's highest task, and yet be conscious that this goal is a vanishing one; to venture the utmost, in order to realize a plan, and yet at every moment be prepared to abandon this plan, when
% --- p. 65 ---
regard for the highest demands it; to be feared by so many, and yet fear most of all for oneself: such directly conflicting determinations are surely apt to place a human being under tension. But the one who cannot, in personal inwardness, endure this tension, does not have the superiority to govern.

Inwardly moved by a thousand anxious considerations, the one who governs must nonetheless be like the resolute one, who has but a single consideration. Penetrated by respect for the universal will, the one who governs must nonetheless act as one who accomplishes his own will. Surrounded on both sides by supporters and opponents, by friends and enemies, the one who governs must nonetheless be just as wary of the love of the former as of the hatred of the latter. The greater the power, the greater the responsibility; the greater the responsibility, the greater the exertion. Thus the power that is otherwise the object of so many's envy can weigh upon the one who holds it as a crushing burden.

If the one in power will only hold to his task as to a cause, the solution is impossible.

To look at the cause is to weigh how much good and benefit can be achieved; how much one, in the given case, may risk; how many human lives one may sacrifice, in order to reach the goal. Such questions cannot, from the standpoint of the cause, be answered in advance, and afterward is too late.

So long as it is not yet full earnestness with responsibility --- whether it now be a sole ruler who uses power according to his own pleasure, in the fancy that he is accountable only to himself; or the several, among whom power is divided, likewise divide responsibility in such a way that each of them goes free, insofar as they are accountable only to human beings --- so long, then, one may indeed continue to keep the worldly cause in pure worldliness; but then it also goes with worldly causes as best it can: the ideal task remains unsolved.

If the task is to be solved, the cause must be transformed into a personal matter between the powerful one and the Almighty.
% --- p. 66 ---
This transformation does not happen because a ruler is so brazen as to exploit religious passions for an ambitious plan, and confuse the will of God with his own supreme egoism; no, the transformation happens through the ruler taking, in earnest, the worldly task as a matter of conscience, for whose solution he is accountable to the All-Knowing.

With this transformation comes responsibility, the unconditioned responsibility in personal truth, the eternal responsibility in divine earnestness. With the ideal impression of this responsibility comes counsel and strength and courage and wisdom: with this impression comes authority.

Viewed outwardly, one might suppose that responsibility and authority stood in inverse relation, so that the greater the responsibility, the more it must crush authority; but in idea just the opposite holds: the greater the impression of responsibility, the greater the authority.

The modern age prides itself on its political progress, and it is without doubt right, insofar as this progress has made it harder and harder to govern. These increasing difficulties, which show us the lot of those in power from a by no means enviable side, must, presumably, gradually force those who govern under the ideal, and teach them whence they are to take their authority.

The modern age places its progress chiefly in this, that those in power are made accountable. It is now certainly a great advance that those in power are made thoroughly accountable; the fault is only that accountability before the tribunal of human beings must, all too often, lack the requisite strictness and earnestness. The true progress will, then, without doubt, consist in this, that those who govern, prompted by their imperfect human
% --- p. 67 ---
accountability, take courage, and appeal to personal truth, in order to attain perfect accountability, and with it, due authority.

To have the highest authority in personal recognition of the highest accountability is, in spirit and in truth, to be the superior one.

%% ============================================================
%%  VII. INSIGNIFICANT PEOPLE (printed pp. 68--79)
%% ============================================================
\chapter*{VII. Insignificant People.}
\addcontentsline{toc}{chapter}{VII. Insignificant People}
\label{ch:forel7}
\markboth{VII. Insignificant People}{}

For personal self-feeling it is surely very disheartening to have to admit: ``You are, after all, at bottom an insignificant human being.'' The notion that most of the human beings with whom one lives are, after all, also insignificant affords only a moderate comfort. Wherein, then, does the difference between significant and insignificant human beings consist? Is there truth in any human being whose life is, in every respect, without significance?

Everyone at once senses how slippery this point is, since the expressions ``significant'' and ``insignificant'' can, after all, be taken in the most different senses. If the inquiry hereby indicated is to lead to any result, then the task becomes to make clear how the ideal can descend into the so-called insignificant, can hide itself behind the phenomenon, and precisely thereby transform even the trivial, and raise the everyday to a higher personal significance.

Such a movement can by no means be made clear in a general consideration; for the general consideration must either overlook the insignificant, or lose itself in insignificances. The peculiar interplay of the significant and the insignificant, which characterizes everyday life more than anything else (an interplay in which the most significant is often enough dismissed as an insignificance, while, conversely, the insignificant can again rise up, and acquire significance), this interplay can only be made clear in an individual portrayal of the individual.

In order to avoid useless prolixity, I shall then again permit myself to introduce my narrator, and let him relate what he calls the second fragment of his life's experience.

``As my first experience,'' he began, when we, on another occasion, had again brought up the personal, ``belongs to public life, so this following one now belongs to private life.''

``In my parents' house, where I continued to live after my homecoming, I at first gathered with my former companions. But, since, on account of my unfortunate adventure (something I preferred to avoid touching upon), I could not rightly share their naive enthusiasm for the new upsurges in politics, there set in ill humor, tension, and coldness. I forswore politicking; my politicking comrades withdrew: removed from all affairs of state, I now belonged entirely to the family.''

``Some profit I did, however, wish to have from my misfortune. `Even if I have not yet been able to discover where personal superiority is,' I thought, `then I shall at least grind my spectacles to see where personal superiority is not.' With this resolve I then began, in all quietness, to practice myself in the critical, the leveling, the negative.''

``All the persons in our rather extensive circle of acquaintance now had to pass in review before my critical spectacles; I sharpened the censorship in spite of liberalism's progress; I held muster over the insignificances. Anyone who could not stand before my ideal tribunal was, however decent and useful the person might otherwise be, mercilessly demoted, and reduced to the rank of a round zero. There was almost no one who could stand; most fell through pitifully. A whole crowd of acquaintances, distant and near,
% --- p. 70 ---
family friends of both sexes, uncles, cousins, and half-cousins, I had already rejected; they had all, rather good-naturedly, resigned themselves to being dissolved into decent nullities. Only one still remained, a single one in the shape of an old spinster, my mother's half-sister, the family aunt. This old spinster was obstinate, and would not let herself be dissolved at all. She stood firm in the claim that she was a real unit, and would by no means be reduced to zero.''

``It was, in my mind, no more spite or malice that drove me to level the aunt than it is spite or malice that moves the naturalist to wound an animal, to cut into a frog; on the contrary, it was the feeling of my spiritual calling, it was the pure zeal for the idea, that lay at the bottom of my negative endeavors. What provoked my criticism was, moreover, not any outright presumption on the aunt's part; she was so far from being presumptuous that she might even readily be called modest; no, what provoked, I might say: ideally irritated me, was a certain inexplicable, to me utterly incomprehensible emphasis, which she, in all she said and did, visibly placed upon her personality.''

``It is, nonetheless, a quiet pretension that lurks in this insignificance,'' I repeated to myself, when I was in a bad mood (and I was now almost always in a bad mood). Wherein, then, might such a pretension, according to her fancy, be grounded? In her rank and dignity? --- She is an aunt. Ah, well, but, more closely examined, she is really only a half-aunt. Besides, aunthood is, after all, a dignity attained solely through the virtue of others, without merit of one's own. Such a dignity can by no means ground any pretension.

``Her beauty? --- An old creature, tall and gaunt and sharp-shouldered, who shakes her head. Good heavens, here are, indeed, perhaps some ruins; but on such ruins one can surely not, at her age, wish to found personal pretensions.''
% NB printer's defect (logged, not corrected, mirroring transcription.tex at
% this site): this paragraph has no opening quotation mark although it
% plainly continues the same series of ironic rhetorical questions as the
% paragraph before it ("Her beauty?" ... "pretensions.") -- it is only
% closed later, by the closing mark after "marketplace." below.
% --- p. 71 ---
Her fortune? Well, if she had at least been a wealthy aunt, a rich aunt, an aunt who could remind one that there is such a thing in this world as a will; yes, that would be another matter. Then she might (at least in her own eyes) be entitled to appear with aplomb; for a will is an aunt's aplomb. A rich aunt, who has the great will in her pocket, has, at the same time, the whole great family in her pocket. That is a genuine aunt; she is the real thing. Such a genuine, solid aunt is far from being an insignificance: she is a real magnitude, even more, she is an ideal magnitude, a mythological figure, a goddess of fortune. She looks around her in the family circle, as though she would say: `Love me now, my children, consider who I am; I am, after all, the goddess of fortune, not the blind one, though, no, the seeing one. Let me now see your tenderness, your love, your devotion: none of you can know, after all, who will in the end be the fortunate one.' Such a genuine, solid aunt is, then, no tedious piece of furniture, far from it! she is an economic factotum, a household oracle, a personified family lottery; a lottery that promotes morality. Such an aunt, after all, sets the whole family to exercises in virtue: everyone stakes his contribution in obligingness, in fine attentiveness, in tender caresses; the one outbids the other, for none can know, after all, where the great prize will finally fall. Such a genuine, solid aunt is like a spirit, a sibyl: there is something prophetic about her. The older she becomes, the more prophetic she becomes: she is amiably prophetic, and prophetically amiable. The gaunter she becomes, the more amiable she becomes; the weaker she becomes, the more she shakes her head, the more strongly she draws, the greater the hopes she awakens. --- Had the aunt possessed such golden qualities, I should have forgiven her, if she made some claim to personal superiority, yes, I should have forgiven her; for this shabby world is, after all, once and for all, like that. But aunt and fortune! it is surely long since the two of them last saw one another; that is testified not only by the pinching
% --- p. 72 ---
in her private economy, but her old-fashioned, well-preserved finery is, indeed, nearly fit to be preached about in the marketplace.''

``Fermenting in these thoughts, I needed air. I tried a couple of times to get at the aunt by opening a dispute about the ideals of personality and about the difference between significant and insignificant human beings. In vain! The aunt never disputed about the ideals, and pretended as though she did not know the difference between significant and insignificant human beings at all.''

``Meanwhile I felt an irresistible urge to communicate my observations; I felt that these observations were too profound, too interesting, too instructive for me to keep them to myself: I had to get air, but how was I to get air? To speak of it to my university friends was impossible, I could not get a word in edgewise with anyone for sheer politics, and an attempt to bring up the aunt at a general assembly would only have made me ridiculous. A lecturer's rostrum I could not mount either; I had as yet no degree, and, besides, the aunt's subject matter was, of course, not scientific either.''

``So I had to see about getting air in another way. I gathered myself a relatively numerous and grateful private audience, consisting of my own siblings, brothers and sisters, all of whom were younger than I.''

``The aunt's name was never mentioned in my lectures; but, for our youthful entertainment, I sketched out various, in my mind quite satirical, portraits of a female individuality I called Old Dignity. In the most whimsical situations I could manage, I then depicted for my listeners, male and female, how Old Dignity combined the dainty with the thrifty; how Old Dignity marveled at her fourpenny presents, as Sultan Suleiman marveled at Morgiana's `winter apples'; how Old Dignity carried herself, as one about to enter the council of state; how Old Dignity sat
% --- p. 73 ---
in the council of state and shook her head, as though she would say: there is but little hope here, if you are expecting an heir to the throne.''

The older ones in the audience were quietly amused; but the youngest, a little sister of 10 years, could at last not keep her composure. She betrayed our secret with the plain declaration: ``why, that's just like Aunt Clara, it's certainly Aunt Clara you mean, Frederik!''

``The blood shot up into my cheeks; the childish declaration was a stab to my conscience. I felt it as a reproach, and now began, with a guardian's gravity, to rebuke the childish misunderstanding. `So respectable a person as Aunt Clara must on no account be confused with these jesting sketches. She might well have her little peculiarities; but it was nonetheless most improper to mock either her or any real human being's personality. Our entertainments aimed, after all, at a higher culture. It was the demand of the new age, to be at home in the comic, to see the contrast between the great and the insignificant in itself, to make acquaintance with the negative, and to raise oneself above narrow-mindedness.'''
% [printed: this paragraph opens two nested quotations ("The blood shot
% up..." and, mid-paragraph, "So respectable a person...") but only a
% single closing quotation mark is printed after "narrow-mindedness." --
% transcribed as printed, mirroring transcription.tex; the outer and inner
% quotes appear to share one closing mark.]

``This forceful admonition worked; my audience did not misunderstand me again. There occurred, however, a quiet transformation in the light in which we siblings gradually grew accustomed to seeing old Aunt Clara.''

``My mother noticed the change, but could not grasp how it hung together. She was a loving soul, but without personal superiority. She began to reproach us; a year and more passed, it grew worse. There was almost nothing in any single instance; every single instance always dissolved into something most insignificant, something which I, with my fine talent, even knew how to make ridiculous. My mother's sorrow, fruitless zeal, and increasing reproaches only made the evil worse.
% [printed: no closing quotation mark at the end of this paragraph --
% printer's defect (unclosed quote), transcribed as printed, mirroring
% transcription.tex; the quote is not closed anywhere in the remainder of
% this batch.]
% --- p. 74 ---
We siblings, who had, up to now, truly respected the aunt, now began properly to be annoyed with her, as though she deliberately caused us trouble, since it was, after all, for her sake that we had to suffer so many unpleasantnesses.

One day --- I still remember it so clearly --- when my mother had properly moralized to us, I spoke up boldly, and asked (though, as it seemed to me, with all due decorum) quite bluntly: ``What, then, has Aunt Clara complained of? what has she really reproached us with? She comes here to the house quite as usual, after all. I will not speak of our larger gatherings; but there cannot be three people here at our house of an evening without our also having the pleasure of seeing Aunt Clara. Do we not show her every possible regard, indeed the most punctilious attention? Do we not dutifully bore ourselves, when her thrift has once in a while brought it so far that she, as she puts it, truly looks forward to seeing all of us at her place one evening? Were we not, all of us, at her place the other day on the occasion of her sixtieth birthday, and that even twice, since we already called in the morning to offer our congratulations? Did that not take place with all possible decorum and solemnity? Did she not receive us with a mildness as condescending as though she were a queen? There was, indeed, a regular parade; we filed past her; she looked at us one by one, as though she wished to count us, and satisfy herself that we were all there, and look, we were all there! What, then, does she have to complain of, what does she have to lament over?''

``When I had said this, my mother rose from her chair, deeply composed, strangely transformed; it was as though she grew a head taller; I seemed to see personal superiority come into being before my very eyes.''

``You ask, my son,'' --- she began, choking back a tear --- ``you ask, Frederik, what Aunt Clara has complained of? She has not complained; she never complains. She has suffered much, gone through much; God forbid that
% --- p. 75 ---
any of you, my children, should suffer what she has suffered; but she never complains, she never laments.''

``She lives in narrow circumstances, does most of her work herself, and scarcely has her daily bread. It has often puzzled me how it really hangs together, for she did, after all, receive a not inconsiderable capital after her late mother. I must suppose that she has, in one way or another, been defrauded; I have often had it on the tip of my tongue, and wished to question her, but I have never been able to: she always brushes it aside. No, she never complains, she never laments.''

``It is I who complain on her behalf, my children, it is I who lament your unreasonableness. You have, after all, all of you formerly been so fond of her, and she is still, unchanged, the same: what, then, has happened here?''

``She often says a kind word to you; you act as though you were absent. She likes to show some small attention or other; you almost act as though you did not notice it. You assure me that you respect her, and that I cannot reproach you with any single instance; I will, however, mention a single instance, and show that I unfortunately have only too good reason to reproach you.''

``It was the other evening; you were going to the great ball. Aunt is, at this time, as you know, very unwell, and does not like to go out, but she still could not deny herself the pleasure of coming up here, and seeing you in your finery. She finds, you see, that you girls look well, and has a quiet joy in looking at you. You did indeed notice that she was here; but you certainly did not notice how deeply you wounded her. These are trifles, but in such trifles the deepest injuries often hide themselves. Do you remember, Sophie, she wanted to pour a little \textit{Eau de Lavande} on your handkerchief; you drew the handkerchief to yourself, as though you feared something old, that had lost its scent. Do you remember, she wanted to adjust your dress a little, Augusta; you turned away, so that she should not soil it with her fingers.''
% [Eau de Lavande printed in antiqua/roman type against the surrounding
% Fraktur -- the first confirmed \textit{} instance in this book, mirroring
% transcription.tex.]
% --- p. 76 ---
``You noticed nothing; nor did she let anything be noticed; but I noticed well that a tear glistened in her eye. Dear children, can you not see that you do her a bitter wrong, her, who loves you so deeply, so selflessly, when you deny her even so much as a kindly, friendly attention?''

``That was the last scene of that nature that occurred in our family; we now had other things to think about.''

``My father, who in his respectable position was esteemed as a capable civil servant, suddenly fell ill; it turned into a violent nervous fever, of which he died after a month's time.''

``In all that time we did not see the aunt at all. She wrote notes, she sent word: she was bedridden, and lamented greatly that she could not come herself. There was, however, no one who paid much attention either to her inquiries or her laments. My mother was so exhausted from keeping watch at night, and took my father's death so much to heart, that she sank into a lethargic faintness: we feared for her reason.''

``With scenes of lamentation I will not trouble you; but you can imagine my despair, when I discovered how matters stood. My father, who in his profession was a capable man, but rather too much subject to worldly appearances, left behind a debt whose annual installments had to consume two-thirds of the widow's pension.''

``The creditors came storming in from every side. A number of new pieces of furniture we had gotten on monthly installments; these were now gradually reclaimed by their rightful owners, and the rooms stood empty.''

``The room in which my mother used to sit, in our days of prosperity, and which we children jokingly called `Our Lady's Glory,' was now stripped of all its ornaments; a simple table with a few matching chairs had taken the place of the elegant furniture. There my mother now sat in the corner, staring at the emptiness.''

``It was toward autumn --- allow me to be a little
% --- p. 77 ---
circumstantial here --- it was an afternoon toward evening, I had just come home from one of my first tutoring lessons. I knew what sight awaited me when I went in, so I lingered a while in the corridor, and left the entrance from the stairs open, for the light's sake, and to breathe more freely.''

``Now I felt how insignificant I was. I, the eldest of so many siblings, could not even help myself, let alone the others. While my father lived, I had not troubled myself about any occupation. I cultivated natural science with success, and aimed at a high goal. I already had a book collection and a small cabinet of natural specimens; now the whole of it was advertised for sale in the newspaper. From a carefree, high-aspiring naturalist I was then suddenly transformed into a poor, insignificant tutor: to what insignificance a poor human being can, after all, sink!''

``I stepped into the great hall; it was desolate and empty here, as at a moving out. I opened the door to my mother's room: what a sight! there she still sat in the corner; my youngest sister lay kneeling on a stool, resting her head in her lap. My brothers had thrown themselves down among one another on the floor; the elder sisters sat resting on their arms over the table. No one stirred. There was no life here: I stood as at a funeral. My breath was about to stop; I felt a pressure, as though I would sink to my knees.''

``Then the door on the opposite side opens softly. I see a female figure, tall and dressed in black and pale as a corpse, shaking her head. In the strongly blazing evening glow she resembled a spirit-being more than a human being; but I soon recognized her: it was Aunt Clara. The involuntary shudder that ran through me in that first moment gave way to an indescribable trust, devotion, and hope.''

``She remained standing, and held onto the doorpost, as though she feared she might collapse; I approached her, pressed a burning kiss upon her hand, and led her in.''
% --- p. 78 ---
Without the least notion of what would now follow, I had involuntarily a feeling that we were saved, that this tottering, corpse-like figure still had the strength to rouse us from our despondency, gather us around herself, and free us from sorrowful dissolution.

Life now came into the dead. The brothers and sisters awoke; we surrounded the old woman, and regarded her with a reverence bordering on holy awe. She looked again at us one by one, as though she wished to assure herself that we were all there; and look, we were all there, but no one could say a word.

Then my little sister threw herself at her feet, embraced her knees, and burst out sobbing: ``We all beg your forgiveness, Aunt Clara, Frederik too begs your forgiveness.'' You can well imagine how this outburst affected me.

But the aunt, who never cared for what she called solemn scenes, turned quickly to my mother, who still sat motionless. She addressed the sick woman with a mixture of tenderness and severity, precisely suited to rouse her. When this did not at first seem to work, she said, in sharp earnestness: ``You do not know me, Maria; do you, then, not know him either, him there, the man with the crown of thorns, who bears his cross?'' and at the same moment held up, before the sick woman's staring gaze, a picture of \textit{Ecce homo}, visible in the clear evening light.
% [Ecce homo printed in antiqua/roman type against the surrounding
% Fraktur, mirroring transcription.tex.]

After repeated effort she succeeded, by her incomprehensible strength of spirit, in truly calling my mother back to some awareness.

She then passed imperceptibly into the more everyday tone of conversation, and spoke practically of practical matters. She now explained to us what we had never dreamed of, that she, far from being poor, was even quite wealthy. She had, she said, long feared what had now happened, and had always saved for our sake. Her capital had, through fortunate circumstances, grown to more than double, and it was now her only joy, the only thing she, as an old spinster, had to live for, to secure us a decent
% --- p. 79 ---
livelihood.

``I see, Frederik,'' she added, ``that you have advertised your book collection and your cabinet of natural specimens for sale in the newspaper; that mistake we must have corrected. You must not sell your books; you must continue to cultivate the science to which you have a calling. We hope, indeed, and look forward to your one day becoming a great and significant man.''

``If you now vividly put yourself in my place, my friend, you can well understand that I came to have other thoughts about the ideals, and about significant and insignificant human beings.''

With this he ended his tale, and with this I end this lecture.

%% ============================================================
%%  VIII. A TUTOR (printed pp. 80--92)
%% ============================================================
\chapter*{VIII. A Tutor.}
\addcontentsline{toc}{chapter}{VIII. A Tutor}
\label{ch:forel8}
\markboth{VIII. A Tutor}{}

Personality's essential tasks are ideal. Now, of ideal tasks the rule generally holds that a mere immediate application of the corresponding natural gifts is not enough; but that formation and instruction are needed, if anything is to come of it.

There are, for example, many human beings who have a good memory, good judgment, and a natural gift for narration. But without a definite direction and a corresponding study, they nonetheless do not become learned in history. The one who is to become thorough in history must begin by going to school, or having a tutor; otherwise nothing will surely come of it.

The one who has an ear and a good natural voice can, of course, sing. If there is any art that can be said to begin naturally, it must surely be the art of singing. And yet it is, in a higher sense, an art: the singer's task is more than merely natural, it is spiritual, ideal. As soon as one has seen that here is an ideal task, one sees at the same time that the aspiring singer must go to a singing school, or at least have a teacher, a tutor.
% --- p. 81 ---
In the world of culture, where there are so many both scientific and artistic tasks, there are, then, also many schools, many teachers, indeed many tutors.

One has, as is well known, tutors on different terms and at different prices, some at a mark, others at a specie-dollar. In London, where one must really make progress, one even has a singing teacher at five specie-dollars an hour.

That is natural, that is reasonable. If one wishes to advance in a subject, and bring it to something, one cannot prize an excellent teacher highly enough; even if one paid him ever so much, there is still always something that cannot be paid for.

But now I have often put to myself the question of how matters really stand, in this respect, with the personal.

Many human beings seem to think that here absolutely no special instruction is needed, that the personal is something that must come of itself, something wild-growing, which, at most, needs only the tending of general culture in order to thrive in the fairest luxuriance.

Should this now really be so? should we, perhaps, here meet with a quite peculiar ideal task, for whose solution no peculiar, corresponding instruction was required? This view is undoubtedly false. If the more ordinary ideal aims cannot even be rightly grasped and pursued without corresponding guidance, how then should it be conceivable that the spirit's highest and holiest ideal could?

If there really existed a tutor in the personal, a reliable and thorough teacher, who could assure the attentive and diligent disciple of rapid and reliable progress, then it seems to me that I could not pay him too dearly, even if I, were I only in any way able, should give him the same as one gives a singing teacher in London.
% --- p. 82 ---
But the question is whether such a teacher can be found. Lessons are otherwise offered, after all, in every possible subject, and from the different prices one can, with fair probability, infer the teachers' different competence: is there, then, anyone \emph{who offers lessons in the personal, and what does the lesson cost?}

This complicated question is not easy to answer directly, and, so to speak, off the cuff. The matter is veiled; if we are to succeed in clarifying this difficult point, we must, from the obscure beginning, see to spelling our way forward to a somewhat comprehensible concept.

With what is veiled a poet could now best help us along, if it should happen fortunately enough that he treated precisely the subject on which we seek enlightenment.

On this occasion I must then very much regret that I have not been able to bring along a book which I owned some time ago, but which has, unfortunately, entirely gone missing (one gets so many new, interesting books that one scarcely has room for the older ones); had I now had this book at hand, and been able to read aloud from it word for word, it would have helped us considerably, for it was, you see, such an ingenious book by an ingenious poet.

It was really a children's book, a book with pictures and text. But, as it is, after all, given to the poet to speak in the same language to both children and adults, the book was also furnished with a preface, from which one could see that the author had expressly set himself the task of uniting the useful with the pleasant, since his aim was precisely to amuse the children and benefit the adults. The book was, therefore, not only very poetic, it was also moral, it was instructive, I dare well say that, in this chapter, it would have satisfied --- even the old portrait in Andersen's fairy tale, so instructive was it. It bore precisely the title The Tutor.
% --- p. 83 ---
Since I now, unfortunately, do not have the book, nor do I know it by heart, I must confine myself to relating what I can remember. The consequence of this will naturally be that the poetic unfortunately gets lost; but there is nonetheless always something gained, if it only succeeds in retaining the instructive.

It is, then, a lost book, a rather large picture book, from which I here permit myself to show a few pictures.

Here is, then, for instance, a picture with the caption: ``The Student.''

At a large writing desk sits a young lady, pale and interesting, with a little ink on her finger. Here it is in black and white that this lady is --- a student. She is, then, also in a way a \textit{Prima Donna} (not yet the first, but very nearly the second \textit{Prima Donna}) in the Ladies' Academy of the Sciences. She does not yet give lessons, she takes lessons: she is arming herself. How bravely she arms herself can be seen from the large table hanging under the mirror: ``Seven foreign languages, History, Geography, Music, Singing, Drawing, Painting, Natural History, Mythology, Perspective, Fortification, and Astronomy.''

The student has, in this moment, lulled her scientific thoughts to rest and brought her little steel pen to rest upon an unfinished astronomical treatise; she then happens to come to think that perhaps yet another heading ought to be added to the schedule.

``There undeniably exists something called the personal. If one can learn all the rest, then surely one must also be able to learn this too. What, then, does one really mean by the personal? It could indeed be quite interesting to find out something about personality. To be sure, personality has not yet made it onto the schedule, but it might perhaps very soon make it onto the schedule, and then it is always good to be a little ahead of the others.''

The student would now really wish that there were some expert she could talk to, that she could get a lesson
% --- p. 84 ---
in the personal. As she sits in these thoughts, there is a knock at the door. It is a lean, narrow, pockmarked, lanky, but clear-eyed person in a long coat who steps in, and greets her. --- His errand? --- ``He offers instruction.'' --- His subject? --- ``He is a tutor in the personal.'' --- His fee? --- ``He does it for free, and is at her service whatever moment is convenient.''

However agreeably our student is surprised by the first part of the introduction, so questionable does she find the last: ``for free'' and ``always at your service.'' She turns to the side and looks out the window, like one who is having second thoughts. As she looks out the window, and has second thoughts, the person vanishes.

``But what is one to make of that? To come as one who is at your service every moment, and then not even be able to wait a moment, while one looks out the window. Surely there ought not to be anything illogical about personality? And he means to do it for free! That is proof that he lacks worldliness, and that there is nothing to him. There is, after all, no tutor who can be well served by doing it for free. Even if he were ever so excellent, and had means, and taught \textit{con amore}: he still would not dare do it for free, for he would risk losing his reputation. But his reputation? He surely has no reputation. Besides, I forgot to ask him about his references.''

With these concluding reflections our student turns once more back to the astronomical treatise. What she later got out of the personal, of that history reports nothing; but that was what was to be reported, that the tutor is willing to do it for free.

Another picture has for its caption: ``The Invalid.''

In a small, old-fashioned, grave, dim room, with large flagstones and mint plants in the window, sits the invalid. He is a worn-out old man in a dark blue coat with a red collar, --- without doubt an old sailor in the veterans' almshouse.
% --- p. 85 ---
This picture, then, promises to be dull, for life in the veterans' almshouse is dreadfully monotonous. Here there are, after all, neither general assemblies nor balls nor carnival nor masquerade, and so there is but little to live for.

The old man, however, seems to accept it. His gaunt hands he has folded upon a large open book: he sits like one coming from a long voyage, looking in toward the harbor; he is like one taking leave of memory, in order to follow hope.

``Yes indeed, there were, after all, many reefs here; there were so many a blow to be had; the hardest, --- that was a blow to the heart. . . .''

The old man would gladly have lost yet another leg, and would rather have gone down three times over, than that he should see that painful sight, when his son fell on the deck at his side.

``He was my only one; I had no one left but him. And yet, it was good too. He fell on the path of honor, he fell with distinction, as we fought for the flag, and won the victory. --- But now the other. . . . This? Ah, that was my own fault, yes, there is much that I have truly repented. --- Peace, it is very good! And He, up there, He who made it all very good, will surely settle the account in mercy.''

With these thoughts the old man perhaps believes himself alone, but he is not: there stands a figure behind the chair, it is --- his tutor. If I see rightly, the figure resembles the fallen son, the son who fell ``on the path of honor, and was painfully missed by his father.''

The old man does not see him; for he looks back no more. He is as though on the wide sea, as in a swell on the high sea, and yet, he sights land: he sees ``the bright harbor and the flag of the Cross and the ship with the great anchor.''

Life in the veterans' almshouse is undeniably monotonous; but for all that it is not dull, when one has --- a tutor.

A third picture has for its caption: ``The Forsaken One.''
% --- p. 86 ---
It is a young, rosy-cheeked little girl, who is just about to cry, and why? Because she feels only too well that she is, alas, an insignificant girl.

She had, indeed, also begun to receive lessons in languages, music, and drawing; but that was, unfortunately, sadly interrupted. Now she is in service with a noble household. She had once hoped to become something in the world, yes, she had really hoped that she would nonetheless manage to become a governess. Heaven willed it otherwise. She has not become a governess, not even a housekeeper, but the housekeeper's assistant, who is good with the children.

Of her conduct, diligence, and character there is nothing to be said against. She feels that she has come to a good household and into a good house. She is also grateful for this, and is already so fond of the little baronesses that she could gladly love them as her sisters, if only she dared.

Nevertheless she is seldom or never glad; but often has water in her eyes, and prefers to sit and cry, when she is alone. ``Why does she cry?'' It has, after all, already been said: she feels only too well that she is an insignificant girl.

When she really comes to think it over, she can quite properly feel pity for herself, and pet herself, and speak to herself as to the forsaken one.

It was a Sunday afternoon; she had just been to church. The family was out; the housekeeper was out too: our little assistant sat alone in her room.

She had just finished a letter to her mother, telling her how well off she was, and how grateful she was, and that in gratitude one could, after all, get through anything, if one were only good and pious, and held fast to dear God. It was a true letter of comfort to the poor mother, who lay sick at home in the poorhouse.

But the memories she had, on this occasion, refreshed, and the moral constraint she had had to impose upon
% --- p. 87 ---
herself, in order to conceal her melancholy, now had the sad effect that the melancholy became even stronger, and was on the point of overwhelming her. She folds the letter together, but can scarcely get it sealed; she is on the very point of drowning it in tears: only now does she feel what it is, to be the forsaken one.

As she now sits sunk in all this sorrowful mood (I almost believe she was on the point of imagining that she was a widow, a young, but poor, most sorrowful widow), she becomes aware that an old man is standing at her door.

To her tear-swollen eyes the man looks like a pious, an honest beggar. But he is nonetheless no real beggar; he is, rather, surely a holy man; he resembles, after all, St. Peter in the legend: it is the tutor.

Since the old man otherwise needed nothing, and was to have nothing, she offered him a chair, and now he sat down directly opposite her. He then asked her why she was crying. She assured him, to be sure, that it was nothing, that she was not crying about anything; but that could now not help at all. The old man now knew so well how to speak for himself, and to speak with her, that it almost seemed to her as though he had already spoken with her before, and was an old confidant who knew her thoughts.

So she had to confess, then, that she really cried, because she could not, like others, have lessons in languages, music, and drawing, and because she could never, never become a governess, and because she was and remained an insignificant girl, and because --- she felt only too well what it was, to be the forsaken one.

``You wish for lessons in languages, music, and drawing,'' remarked the old man with a friendly smile; ``you need not cry over that, my child, that wish could well be granted you. I am an old tutor, who knows languages, music, and drawing all three. I have practice in teaching, and shall gladly give you a lesson, when you are willing to learn.''
% --- p. 88 ---
And she was so willing to learn. And so the old man then began to speak the language. It was the pure, clear language, so calm, so moving, so distinct, so heartfelt, the language that penetrates into the soul. It was as though she heard her own silent thoughts, and every dark feeling got its own tongue. That she had to confess: he was a man who knew the language.

``And what, then, is music, my child?'' asked the old man. And he stressed his question so that she could herself answer this question; for the strings of the heart resounded at this language and this voice.

``And now drawing?'' he continued, and raised his forefinger. And before her gaze stood a world of images, lovely, life-giving images! It was a dream, a blessed dream, whose figures, as though by magic, took on all the power of the real, in light and shadow, movement and life and color.

The hours passed, it was already late: the tower clock struck twelve. The old man rose, and added a couple of grave admonitions: ``Be silent, my child, keep what I have entrusted to you in a beautiful and good heart. Never occupy yourself with things that do not concern you. Do not forget to watch and pray, and do your work. If you promise me this, then all things shall work together for your good; if you promise me this, then we shall see each other again at a convenient time.''

All this she then had to promise him; and, when she had promised him this, he went quietly out the door.

It was half past midnight, the family had still not come home; it seemed to the waiting girl as though it lasted an eternity. She, who had just now been so rich, felt herself again, all at once, so empty and poor. She tried to repeat something of what had been said, but everything was faded, dim, and hazy. The beautiful, the glorious, the old man had, after all, taken with him; whatever she had not been before, she was now, indeed, the forsaken one.
% --- p. 89 ---
An hour past midnight the family came home; our little assistant was now to be of service. But she was so absent, so confused, and went about everything so clumsily, that the gracious lady, who supposed that she must, during her waiting, have fallen asleep, and was still walking in her sleep, took pity on her, and ordered her to go to bed.

Weeks passed; it grew worse and worse. The little assistant wanted to know everything, meddle in everything, judge everything: she had moods. The housekeeper thought that there must surely be someone putting fancies into her head. She got ideas, she joined in about the children's instruction, she weighed up the music teacher, she even went so far as, at last, to dispute with the governess.

The gracious lady, who had kept an eye on her, and sent her many a warning look, found at last that it went too far, and had the little assistant called before her.

In a tone that betrayed motherly concern as much as mistress's displeasure, she set before the young girl the impropriety, indeed the unseemliness, of her conduct, and added, calmly, but with great firmness, that she would have to leave, if she did not change, and become again what she had been at first, an attentive, obliging, and modest girl. She might now go up to her room, and think it over.

As our little assistant steps into her room, the old tutor sits there, waiting for her. The bitterness she felt at the well-deserved reproof now turned into fierce anger against the old man. --- ``What did he really want with her; was she not unhappy enough already; it was, after all, he who had (as the housekeeper indeed also said) put fancies into her head, that she felt very well. If he wanted anything of her, he must at once declare whether he dared promise her that, by his lessons, he could bring her so far that she could, at the very soonest, become a governess (at this point she could no
% --- p. 90 ---
longer hold out); if not, then he, together with all his wisdom, was a poor, helpless wretch, and must leave her at once.''

``Vain child!'' --- began the old man --- ``vain, ungodly child, what was it you promised me? So thoughtlessly have you forgotten my admonitions. It is not my gifts, it is your faithless, self-willed disposition, that brings you to ruin. You have failed in your duty; you have wronged your superiors; you must repent, and ask forgiveness. I say unto you'' --- and now he raised his forefinger --- ``verily I say unto you: you must repent, you must at once repent, and ask your mistress's forgiveness. If not, then I will go away, and you will never see me again. Believe me, it is only reluctantly that I go away. I would most reluctantly leave you, my poor child; consider! when I leave you, then you are, in truth, the forsaken one.''

The young girl sank at these words like a bent reed. She threw herself at the teacher's feet; she wept such hot tears as she had never wept before. It was an hour in which she learned, for the first time, what it is to weep.

She humbly begged the old man's forgiveness; she begged her mistress's forgiveness, and it was made good again; and now all was well.

The tutor came and went, without anyone noticing anything. The only thing, though imperceptible, that could, in an inward sense, be noticed, was a quiet growth in the young girl's personal being. --- ``A lovely servant girl,'' --- ``wonderfully gifted,'' --- such a word was sometimes exchanged, quite quietly, between the master and the gracious lady; but otherwise there was no one who noticed anything.

See, then, it was very good; and we adults may now well let the story end. How long the servant girl remained in that position, and where she later came to in the world, is, after all, indifferent, so long as only the tutor gives her a good instruction. Should she then follow in her mother's footsteps, and end
% --- p. 91 ---
by being buried from the poorhouse; that then does not matter at all.

No, it does not matter at all; the old man will surely continue to keep an eye on her; he will surely see to it that she, even in the poorhouse, shall not feel like the forsaken one.

But, it is, after all, from a children's book that the picture is taken, that we must not forget. So, for the children's sake (and since, besides, it is supposed to be quite a truthful story), there was found, in the book, yet another picture of the forsaken one. That picture --- I remember it clearly, it stood precisely on the tenth page to the right; --- there were many unmistakable traces that it was the picture the children liked best of all, and for whose sake they read the whole story.

Here, then, is festivity and pomp at the castle. Here stands the forsaken one, adorned as a lovely bride: the old, gracious lady has herself set the bridal wreath upon her head. How splendid it looks! Everything now gleams in gold and purple. But this is, after all, also the forsaken one's wedding day, the day she becomes a baroness. Here is a throng and a bustle: great lords with stars on their breasts; their servants carry the cloak upon their arm. Both the King and the Queen wear crowns. One can plainly see the carriages outside the castle; there are surely more than seven thousand coaches.

But the tutor? Ah, that's true, wherever has the tutor gotten to; he is nowhere to be seen. --- ``He has surely crept into hiding.'' --- ``He is surely sitting inside the chamber, delighting in all the splendor.''

The most important thing is here too, then, the very most important thing for a bride: there are so many young bridesmaids.

It is a festive day for the forsaken one. The old lady takes the forsaken one by the hand: it is precisely at the moment when they are to go to church. The church, to be sure, one does not see; but the children assure that, if one puts one's ear to the picture, and listens closely, one can plainly hear it chiming.

So the tutor continues, throughout the whole book, to go about under many disguises to all kinds of
% --- p. 92 ---
people in all stations. He teaches those who are willing to be taught; he chastens those who are willing to be chastened; he comforts all those who are willing to be comforted.

He is with the warrior who dies on the field of battle; he is with the mother who kneels at her child's cradle. He stands in the prince's closet, and raises his forefinger; he sits in the madhouse, and breathes upon the confused one's brow. He draws near to the sickbed; he transforms the restless one into an image of patience, and the hour passes, and the struggle ends: the sick one folds his hands, and quietly falls asleep.

So I regret, then, that the book has been lost. Should, however, anyone among my honored listeners, male or female, happen to come across the book, and find my name in it; then I ask the finder to erase my name, and regard the book as his lawful property. Should it then be found that what has here been related is only a poor account of its rich content, that shall be a joy to me; for then it is precisely a proof that it is the right book that has been found.

But, should the tutor himself show himself to anyone at a convenient time, and the one to whom he comes should mention my name to him too, then I thank him sincerely. For it is not a book about a teacher, still less a picture book about a tutor, that I need; I need the teacher himself, the tutor in the personal.
% printed p.92 ends the lecture with a printed horizontal rule (ornament),
% not transcribed as text, mirroring transcription.tex.

%% ============================================================
%%  IX. GUILT IN FRAILTY (printed pp. 93--104)
%% ============================================================
\chapter*{IX. Guilt in Frailty.}
\addcontentsline{toc}{chapter}{IX. Guilt in Frailty}
\label{ch:forel9}
\markboth{IX. Guilt in Frailty}{}

The unity of the ideal and the real exists only in true personality.

To be human, to have the disposition of personality, is then the same as being disposed toward a continued striving, in order to appropriate this content-rich unity.

The feeling that feels only the real (real hunger and thirst) is a merely natural feeling, not yet transformed into the personal. The feeling that feels only for the ideal (the true, the exalted in its eternal universality) is a feeling of fantasy, which lacks reality: realization is in the personal.

The will that wills only the real in the present is not yet any true personal will. It is a will that works for finite ends; it is allied with prudence, endurance, and courage: it seems strong; it can, indeed, also be strong --- in the finite. But, since all its strength is the strength of finitude, it is only essentially strong in finitizing itself, and denying the personal.

The will that despises the real, because it wills only the ideal, is an upward-hovering will of fantasy, which craves the empty. Here content is lacking, here limitation is lacking; here earnestness is lacking: here the personal is lacking.
% --- p. 94 ---
To live personally is, then: to think, feel, will, and act in such a way that the ideal and the real meet and unite in every thought, every feeling, every act of will and action.

Where they do not meet, personality is paralyzed; where they meet in conflict, personality is striving; where they meet in peace, personality is victorious.

To sigh toward heaven is, indeed, surely to sigh for something ideal, but the merely ideal sighing does not help; there must be a pressure of reality here, an essential impression, that wrings from the soul such a sigh, if it is to reach heaven.

That mere reality presses does not help either. The material pressure can be so strong, the suffering, the misery so great, that it must awaken human beings' pity; but that is sensuous pity with the sensuous: such pity is not in heaven. If the misery lacks the ideal understanding, then the suffering lacks the baptism of the spirit, and heaven takes no heed of it.

When, on the other hand, the outward pressure meets the inner ideal striving, when the burden of reality is taken personally, and borne personally; then the sigh that is thereby wrung from the soul will not lose itself in the air, but will reach heaven. It is then a sigh of personality, and there is in heaven One who hears all personal sighs.

Something of this sort we human beings now grant readily enough, surely, in complete generality (insofar as we have not yet become so ``independent'' that we have entirely broken with the Eternal); but thus to acknowledge the ideal in its generality without realizing it in individual living is, at one and the same time, to acknowledge and to deny the claim of personality. It is a fundamental contradiction in our being. To own up to this contradiction in all sincerity is then the very least we can do for the truth. But how are we to interpret the contradiction?
% --- p. 95 ---
The usual interpretation generally comes down to this, that we are frail. That is now very true: we are frail; but the question then is whether, and to what extent, our frailty can be used as an excuse.

If we remain with the merely frail, then we remain, after all, with the mortal, and the immortal departs. The soul is then frail, perishable, and mortal, just as much as the body. If we are pure frailty, then we are pure perishability: the personal in us is then also perishable. Then we would indeed get an excuse, but an excuse that is rather superfluous; for, for pure perishability, it is not even worth the trouble to make any excuse.

If, on the other hand, the personal in us is the ideal, the immortal, then there must, after all, be something in our frailty that is not frail. When, then, the frail nonetheless triumphs over the not-frail: then there is here no longer any excuse; then there is guilt.

For personal development it is, then, an unwavering condition that a human being begin at once with \emph{recognizing his guilt, the guilt that is in frailty}.

But guilt has a double side; it is, in part, a common, general guilt belonging to the whole race, in part a particular guilt belonging to each single human being. These two sides are inseparable in the guilty personality.

If I here were to lay out the account of my life, reviewed and tested by personal truth, and the judgment that followed thereupon were to be made known to human beings: how, then, might it go with me?

--- ``Ah, it would go ill with you; human beings would gloat over you: you would be pitiably exposed to shame!'' --- some brash person might perhaps hit upon answering.

The first I grant. But as for the last, anyone who uses any reflection will at once see that human beings (the brash person included) would have other things to think about than to gloat over such an act of judgment.
% --- p. 96 ---
When the great criminal, condemned by human beings according to human laws, is led to the place of execution; then a spectator may well say: ``Thank God, it is not I''; and another may answer: ``it is no more than he has deserved, it is, after all, his own fault.''

But it is otherwise, when Truth itself applies its personal standard, and judges the guilt, the guilt that is in frailty.

The more ordinary, everyday, and, by outward appearance, insignificant the human being is who is called before the judge, the less there is here of extraordinary crime, the more perilous the act of judgment becomes in its effect upon those of like kind who are to be witnesses. The judgment that then befalls the one will, after all, befall all. For it becomes a dangerous reckoning for the otherwise insignificant human being, if Truth's gaze fastens itself upon the guilt alone, the guilt that is in frailty.

It is precisely this point that we wish to make clear to ourselves in this lecture. In order to make clear so personal a matter, we use the personal; but it would, after all, be dreadful, if now some real human being were to begin laying out his account, and demand that we, in his frailty, should see our common guilt.

Our presentation, which cannot use real human beings, must, then, confine itself to using imagined human beings, imagined persons, figurative representatives. It is precisely such representatives that are found indicated in the narrator's two fragments and the pictures that followed upon them. Since I could hardly borrow such figures, corresponding to my purpose, from a poet, I sketched them, as best I could, on my own, for use in an acknowledgment of personal guilt.

At dømme retfærdig --- to judge rightly of one's neighbor's state of conscience is altogether impossible, and therefore every attempt in this direction is improper; but to let imagined persons,
% --- p. 97 ---
personalities, whose inward being one forms oneself and consequently sees through, undergo a judgment, is an innocent matter. We are, in this case, like creators who judge their own creatures: only this must be maintained, that the judgment is just.

We shall then begin by turning our attention to a pair of figures in the narrator's fragments, namely father and son.

It is not only the men of public life and the heroes of history who are tempted by notions of worldly greatness; ordinary human beings can, after all, fall into the same temptation: father and son are both guilty.

The father, whose character has been indicated in passing by the son himself, may well have been, in his position, a respectable man, a capable businessman, and so on; but when this soul now goes in to the reckoning, and eternity asks him about the yield of his life: what, then, has he to answer?

``That he was a capable official.'' But for that he received his pay, after all; it balances out exactly: here is nothing. ``That in his whole conduct he showed himself an upright and benevolent man.'' But for that he also, after all, enjoyed the respect due ``an upright and benevolent man''; it balances out exactly: here is nothing. ``That he did not despise religion, and never had any dispute with the clergy.'' But what has he thought concerning the Church? what has he recognized in religion? Nothing. ``That he has had noble feelings and good principles, and now and then, when his business allowed him, thought a good deal about God and virtue and immortality.'' But is there, then, in his being any trace that the thought of immortality has ever been unconditionally decisive for his life? Nothing, nothing!

That he leaves his family in temporal helplessness is only a fault insofar as it is a consequence of his vanity; but what has he done, to prepare a way for his own toward spiritual help? Nothing. He was, otherwise, indeed active enough in other respects. He spared nothing on his children's
% --- p. 98 ---
education; he saw to their social life; he procured new furniture on monthly installments: he was a careful family father, a cultivated companion, an agreeable host; but a personality he was not. Since he himself was not a personality, he has, then, also been unable to leave his family any impression of personal strength; no telling trait, no fatherly exhortation, no last word from him could, in the hour of need, be a comfort and support to those left behind.

So this man has been able to make his way, and be active in the world, and go honored to his grave, without even suspecting the secret, that he, in an eternal sense, squandered his life.

Let us pause yet a moment before this picture!

If it is only a fixed idea, this matter of immortality, of the resurrection and the reckoning yonder, then this man is, in every respect, beyond reproach. He has, after all, been something in society, he has, after all, always managed to accomplish something: he has, indeed, come through the world quite well; he can, upon my word, rest in his grave just as secure and calm as the greatest heroes and the very strictest moralists.

But if, --- I say: \emph{if} there now nonetheless really is an eternal life, and if the Eternal One, who must, after all, know what he can demand of a human being, does nonetheless, after all, absolutely demand that the human being here set himself in a decisive relation to the Eternal, and prepare himself for the reckoning; must we not then say, when we look at this man: what frailty, what guilt; what guilt in frailty!

Hvad nu Sønnen angaaer --- as for the son now, the young man (for which he can also thank his father) is, after all, already enthusiastic enough about ``greatness.'' In order to see personal superiority, he travels abroad: he is, after all, still a minor. Existence, which at bottom means well by every human being, comes then also to the minor's aid: he is
% --- p. 99 ---
fortunate enough to reach the Republic at the right time, and see the right specimen of greatness.

But the help is only for the one who is willing to let himself be helped. After his return home, precisely at the turning point when he should have come to his senses, the person becomes guilty.

To lose oneself in bitter self-absorption, to judge human beings uncharitably, and seduce the innocent into doing the same, is an ideal crime, which does not find its excuse in youthfulness. The youthful one is here the frail one, and this frail one is a guilty one.

Existence has, once again, compassion for the frail one, and comes to the guilty one's aid. That human being, so very insignificant in his eyes, the lowliness in which the personal irritated him, stands at last before him in the shape of the superior one: now he is helped.

That it is not the aunt's fortune that helps him, everyone naturally sees. He is, after all, essentially helped only by this, that here is one who, in all frailty, has borne his offensive conduct, forgiven him his arrogance, and, without even outwardly humbling him, led him to go into himself, and acknowledge himself guilty. From the aunt he then inherited what he did not inherit from his father, a true personal impression, an impression he keeps, in repentance, in deep gratitude.

``But now the main character, she in whom the ideal works --- that she is old, peculiar, and frail is obvious --- does she, then, also have the personal in the recognition of guilt?'' I answer: yes, she has it. Such a severity toward herself, such a forbearance toward others, such a superiority in forgetting insults and preserving love, such personal qualities only one can have who, early and late, feels himself bound by the ideal of personality. The one who is bound by the ideal must humble himself before the ideal; but everyone who humbles himself under the ideal acknowledges not only his frailty, but his guilt in frailty.
% --- p. 100 ---
The aunt is, in her inwardness, religious (which she, indeed, does not conceal either, when she admonishes the sick woman), but she is not pietistic. Her strict self-discipline has taught her to handle religious ideas as one handles fire and light: the wordy prolixity of talkative godliness is not edifying to her.

Should, however, any of us be so satisfied with his own reasonable principles and his good works that he even reckoned on her admiration, then this old, experienced woman (though she is neither bishop nor priest, and does not presume to be anything either) might perhaps, in all confidence, speak a serious word with him about our common guilt, the guilt that is in frailty.

So much for the fragments; now something about the picture book.

Once it becomes evident that the ideal works even in everyday life, when the human being himself is willing to work along with it; then our frailty has only one excuse left, and this excuse runs thus: ``We cannot grasp the aim of the ideal will, and least of all do we know what it demands of us in each individual case. If we had the eternal will's express commands for each day, for each hour, for each occurring case, then indeed we could still work at our personal formation; but everything remains in hovering generality: we are befogged, and go astray in the particulars.''

This excuse transforms itself into an admonition: ``Then accept instruction in the particular, take lessons in the personal! Ideal instruction, with its corresponding inspiration, offers itself in every occurring case, on every occasion: the tutor does it, after all, for free, and is at your service whatever moment is convenient.''

Where, now, is our excuse? That the instruction is disregarded, and self-examination neglected, because the ideal, which hides itself in the everyday, comes and goes of itself, without our duly heeding it, can surely not, then, be any excuse: this negligence is, after all, a guilt.
% --- p. 101 ---
It is the picture of ``The Student'' that represents this guilt. It is a little colored picture, in \textit{camera obscura}, of modern versatility. This direction of culture does not lead to inward richness in the personal; it leads to bondage under the schema of knowledge, a brilliant spiritlessness in half-thoughts and learned-by-rote idealities.

From every side complaints are heard over this feminine bondage under the tyranny of the schema. The physician complains, he complains on behalf of health; the poet complains, he complains on behalf of femininity; the moralist complains, he complains on behalf of personality. Here again is frailty, and guilt in frailty.

It is a modern prejudice, that much knowledge and many-sided formal culture should constitute the foundation, or at least furnish the most indispensable conditions, for the formation of personal character. It is so flattering to regard one's theoretical preoccupation with the many ideas as a real living in the spirit; it seems so profound, to overlook reality, to abstract from oneself, and ``look out the window,'' when the personal announces itself.

The picture of ``The Invalid'' I will, in this context, pass over, and merely mention it, insofar as it reminds us of the difference between the two concepts: the tedious and the monotonous. The one who cannot endure the monotonous, but constantly needs outward variety, in order to keep himself animated, suffers from personal emptiness. But to be able to endure the monotonous without working in inwardness is, after all, also a proof of emptiness. To be tedious to oneself, and sit in the midst of the tedious, without even feeling strength and courage to break off the tediousness, is a sad sign of sluggishness. Through spiritual sluggishness a human being can become just as personally guilty as through the empty craving for diversion.

Er der nu en Skyld --- if there is, then, a guilt that shows itself, not only every time sluggishness gains the upper hand, but just as much every time
% --- p. 102 ---
a human being begins to dread the tedious in the conditions assigned to him; then it is seen from this what exertion would be required, if every hour (for I dare not say: every moment) of our self-conscious life were to have personal significance.

It is usually complained that the span of human life is so short; but, if we bring into the reckoning all the time that is, in part, spent in pure tedium, in part squandered on empty, childish diversions; then one might rather be tempted to suppose that the span of human life was, as a rule, too long. Here, then, again is frailty, and guilt in frailty.

In the last picture: ``The Forsaken One,'' several main points are indicated.

``The Forsaken One'' begins by feeling herself unhappy, and she is, in a way, also that. The unhappiness, however, is not, as she imagines, that her higher striving has been halted, and that she has come into a position of service: this is something inessential (she has, moreover, after all, come to a good family). No, the unhappiness is in the inward, the ideal. Her fantasy has been misled; ``The Student'' has become her model: she is, then, in the process of becoming guilty through a common guilt. She feels herself \emph{insignificant}! It only remained for her to have come into the official's house, and gained admission to the young man's ``private audience'': she might have become a talented pupil.

She is, however, helped by circumstances: she ferments in all quietness, she goes to church, she is reminded of the obligation to be grateful; thus prepared, she gets --- a lesson. In the moment of inspiration and mood, she feels herself richly gifted. She is no longer the insignificant one; she recognizes the spirit; she understands, in the bright moment, that the gifts of the spirit are granted her for personal formation, and promises to use these gifts rightly.

But the vain tempts the vain one: to be in such outward insignificance, and yet so inwardly
% --- p. 103 ---
entitled to something far better, seems to her a wrong, whereby existence injures her. It is this feeling that makes her bitter, presumptuous, and capricious: she forgets her good resolve, takes the gifts vainly, and becomes guilty.

Guilt follows upon guilt. On this occasion another main point is touched upon, namely the relation between master and servant.

This relation becomes distorted, when one holds fast merely to the real; for reality is so hard: here we find the hard master and the browbeaten servant. But the relation also becomes distorted, when we, all at once, become so sentimentally philanthropic that we disregard the barriers of reality, and grasp after the pure ideal. In the first moment master and servant then embrace, and give one another the kiss of brotherhood; in the next moment the servant is half master, and the master half servant: the half-master quarrels with the half-servant, and the relation is thrown into confusion.

Therefore it is surely best to take reality along, that the master as master and the servant as servant must, once again, go each his own way, and each take himself a lesson.

Once the tutor has instructed them, they can, after all, meet again, and in good understanding: the relation is then restored, a personal relation in the temporal, which can have its significance also for the eternal.

That, now, ``The Forsaken One's'' employers, ``the gracious lady,'' enjoys personal tutoring, is indicated by the union of forbearance and severity, firm bearing and motherly care, with which she treats those under her.

The little guilty one, by contrast, stands on the verge of becoming doubly guilty: with that her next lesson begins. It is an hour of repentance, an hour when she, who was otherwise so quick to tears, learns for the first time what it is to weep.

Only one more remark about the ending.

The conception of the personal can naturally not be the same among children as among adults. In childlike
% --- p. 104 ---
consciousness the notion of the good is innocently united with the notion of happiness and a good outcome. For the adult, that is to say: the personally fully grown, these notions are dissimilar. But where, now, is the boundary between the half-grown and the fully grown?

If I were a fully grown person, I could only rejoice in the recognition that the struggle for the good, far from leading to happiness, on the contrary, according to the ideal law, leads to what the world calls unhappiness. But, since I am still only half grown (and scarcely even that), I must hold fast to the childlike thought, that the Blessed One in heaven is not like a pagan god, who envies a human being his happiness.

All things, happiness as well as unhappiness, can, after all, work together for our good, when they only serve to awaken the consciousness of our \emph{guilt, the guilt that is in frailty}.
% printed p.104 ends the lecture with a printed horizontal rule (ornament),
% not transcribed as text, mirroring transcription.tex.

%% ============================================================
%%  X. THE STRONG AND THE FRAIL: A PERSONAL RECKONING
%%     (printed pp. 105--117)
%% ============================================================
\chapter*{X. The Strong and the Frail: a Personal Reckoning.}
\addcontentsline{toc}{chapter}{X. The Strong and the Frail: a Personal Reckoning}
\label{ch:forel10}
\markboth{X. The Strong and the Frail: a Personal Reckoning}{}

All who have enjoyed any instruction know, indeed, how it goes: namely, that the teacher proceeds according to certain rules, applying a method suited to the nature of the subject and the pupil's capacity for grasping it. That the pupil, for his part, is under express obligations as regards attention, diligence, endurance, and so on, goes without saying.

Of instruction in every particular branch of formal culture it is, therefore, easy to give a notion; it is, on the other hand, exceedingly difficult to present and make clear how it actually goes with instruction in the personal.

If I am, for instance, to be instructed in some science or other, one ordinarily begins here, indeed, with a teacher and with a book. As a beginner I must then learn much by rote, and take much on the teacher's authority. But, as the pupil makes progress, he is trained to study on his own. If the studies succeed, it ends, after all, with the pupil at last becoming his own teacher. I then no longer assume anything merely because a teacher has said it, or because it stands in a book: I myself must examine, myself test, myself see the truth. No human being, then, has leave to set his
% --- p. 106 ---
immediate authority between my eye and the truth. Anyone who attempts that, I must ask to step aside, for he stands in my light; I do not wish, after all, to look at a human being: I wish to see the truth.

If we turn now to the personal, instruction begins, indeed, with what we call, in the proper sense, upbringing. Upbringing likewise aims at this, that the minor shall become of age. With the age of majority, then, self-instruction should begin; but here comes the difficulty.

Self-education must, after all, aim at an appropriation of the truth. But which truth? The personal truth!

If personal truth really exists, then it must have its being in true personality. But where, now, do I find true personality, the model according to which I can order my conduct?

I look around me; I look to my teachers, my guardians, my superiors; they might still, in many respects, have the right to command me; but, when I ask them about true personality, they all point me to the ideal.

Do I, then, have this ideal in myself by nature in such a way that I can now, on my own, develop it out of myself and by myself alone? By no means! What my teachers and guardians quietly acknowledge, as they step back and let me stand alone as the one of age, I too, after all, must acknowledge: our frailty, our guilt in frailty.

So truth exists for me, then, on one condition only, namely, that I believe in Him who says of himself: \emph{I am the Truth}.

To the ideal man I cannot, then, possibly say: ``Get out of the way! You stand in my light, since you too, after all, are a real human being. I no longer look now at a real human being; I look only at the pure ideal; for you must know: I have, just these very days, come of age!
% NB printer's defect (carried over from transcription.tex, logged there, not
% corrected here): this quotation, opened with ``Get out of the way!'', is
% never closed in the Danish -- no closing quotation mark appears before the
% next paragraph on p.107 begins. Reproduced here with the same unclosed
% quotation, mirroring the source.
% --- p. 107 ---
If this talk of personal truth is not, as I have said, some arbitrary fancy, some whim, some fixed idea: then the real ideal man, the teacher himself from eternity, must be our essential guardian, educator, and guide. To acknowledge this is, then, the true, all-embracing fundamental condition for all true personal development.

But precisely as I acknowledge this (and my thoughts can now not get away from it), I meet a new difficulty, and it is, in reality, to me the very greatest. The teacher has not merely a certain superiority; he has, after all, the superiority of eternity; here is not merely some distance: no, here is, indeed, an immeasurable distance between this one of age and me, the minor, an immeasurable distance between teacher and pupil, an immeasurable distance between the strong and the frail.

My difficulty is not that the teacher's life and conduct should, as far as the accounts go, be obscure, ambiguous, and unreliable; on the contrary, it is, even in every single fragment of the accounts, so clear and manifest, that it, in every particular, reveals the ideal itself, the personal truth, the true personality.

Nor, either, is my difficulty that the guardian himself is no longer visibly present, so that I could speak with him directly, as with another real human being, and ask his counsel. No, on the contrary! I see expressly from his life and conduct that the matter stands the other way around.

As long as the disciples had the teacher himself among them, they could, with the greatest difficulty, understand him. Understanding was still so unclear and imperfect. Only when he had gone away did they understand him; for, when he had gone away, the Spirit came, after all, upon them, so that they became inspired. It belongs, then, to the bridging of the immeasurable gulf between teacher and pupil, that the pupil become inspired. Only the inspired one can personally understand the ideal man's personality; only he who is taught by the Spirit can teach himself according to the model's instruction.
% --- p. 108 ---
Here is my difficulty; for I am not inspired, and have never been inspired. I do, indeed, hear now and then, that there should still be personalities through whom the Spirit speaks directly, or at least almost directly. How this actually hangs together with these more recent, half-scientific, half-poetic, half-prophetic inspirations, I will not judge; only this can I say in truth: to me the Spirit has, as yet, never spoken directly, and neither can I speak directly with the Spirit, and ask it directly.

Nor can it help, either, that I directly appeal to the inspired, and repeat their words directly; for if I cannot understand the teacher himself without inspiration, how should I then understand the understanding of the inspired without inspiration? And, if I cannot understand the teacher himself, when I turn to him alone in all inwardness: how should I then understand him, when I turn to a company of thousands? If it is to be mere company, then the sociable is, after all, precisely the diverting; and if this company is, in spirit and truth, a communion: then I am, after all, most alone with him, with him alone, when I am in this communion.

From whichever side I look at my difficult point, I constantly find that I cannot possibly understand the ideal of personality, unless the Spirit is willing to help me in such a way that I, in the Spirit, can understand him himself --- without the intervention of others --- him himself alone.

Here there must be a reckoning, a personal guidance, that holds itself between the height of inspiration and the merely rational self-instruction. The same spirit of personality that can inspire the great, extraordinary messengers must, surely, also be able to work, by a weaker inspiriting, in the ordinary human being. Such an awakening and inspiriting I call here a reckoning, that can join the frail with the strong. It is to such a reckoning that I have to hold fast: further than that I have not yet come.
% --- p. 109 ---
The relation of inspiriting to self-instruction is difficult to present; in order, to some degree, to overcome this difficulty, I use a figure of speech: the Tutor.

By way of preliminary clarification, I shall, however, permit myself to preface this with a remark on the Church.

In the Catholic Church, as we all know, secret confession holds. The believing Catholic turns, in all matters of conscience, to his confessor as to God's representative, indeed as to God himself. It is not enough that the confessing one acknowledges his individual sins; he must also submit himself to the corresponding penance. The confessor prescribes for him what he has to do in each particular, and watches over it, that the right exertion is awakened and preserved in the penitent soul.

The evangelical Church has a sharp eye for the spiritlessness, the mechanicalness, and the falseness in this guidance. It points out the indefensibility of a human being's conscience and whole personality being thus surrendered into another human being's power; it draws attention, and rightly, to all the abuses that an arbitrary, self-interested, and domineering clergy can make, and really has made, of the believers' blind devotion; it points out the fundamental falseness in such a parceling-out of individual feelings, individual thoughts, and individual deeds, a piecework of spirituality that corresponds exactly to salvation-by-works, and puts faith, the one thing needful for salvation, in the shadow under the multiplicity of good works.
% NB printer's defect (carried over from transcription.tex): the Danish
% prints a stray hyphen wrongly joining "under" and "de" ("under-de gode
% Gjerningers Mangfoldighed") where two separate words are plainly meant.
% Not reproducible in the English word order; noted here rather than forced.

The evangelical-Protestant Church has abolished secret confession; and for the freedom of conscience and personal independence thereby gained, we cannot be sufficiently grateful.

But as regards the religious-ethical, our Church has undeniably found it easier to abolish the abuse of confession than to replace what was abolished with better practical means.

The truth originally aimed at in the commanded confession is to join the ideal with the real, not
% --- p. 110 ---
merely in general, but personally in each human being's life. This task the evangelical Church cannot possibly solve. It can, as a church, never reach, through the general, the personally punctual. In that it, according to its principle, makes, and must make, the layman spiritually independent, it must confine itself to admonishing each one, individually, to be his own confessor, or rather: to confess in secret to the all-knowing one, who sees in secret.

If, now, everyone conscientiously complied with this summons, then the ideal would take powerful hold upon the real, and then our Church would, at the same time, be infinitely superior to the Catholic.

But that, precisely, is not how it has gone. The Catholic Church has, in an outward, spiritless way, drawn the eternal down into the finite, even into the petty, and thereby wronged the ideal; but it has been shrewd, it has looked after its own interest. It has seen how human beings are, and how they conduct themselves with the divine, when they are merely admonished in complete generality, and for the rest left to an indeterminate inwardness in the personal.

The Catholic Church still surprises us, in many ways, by its worldly-wise, well-calculated tactics. Here, then, is something for a teacher in our Church, a profound theologian; but I, after all, am no theologian (not even a parish clerk, but only a schoolteacher), so I return, then, to my task, my figure of speech, that is, to my Tutor.

Inspiriting, the ideal admonition, is, by this figure of speech, made visible as something that is not merely above the human being in general, but also as something that offers itself to the human being on various occasions and under various circumstances; that is why the Tutor has so different an appearance.

The eternal and yet temporal, that is: the personal reckoning, puts on, so to speak, costume according to time and occasion. Sometimes it shows itself to the human being under the shape of
% --- p. 111 ---
something outward, or at least as something the person concerned shall and ought to turn to. It rests, then, upon the human being himself, whether he will make use of the occasion, take counsel with the ideal, throw open his inward being to the knower's gaze, and thereby become manifest to himself, or prefer to hold fast to the general and dissolve the particular into the diversion of thoughts.

Sometimes the ideal, the inspiriting, rises up within the soul as if of itself, before the human being rightly knows of it; it is the divine that seeks the human being: it comes in mood, and gives mood.

Thus, to recall the picture book, ``The Student'' forms a counter-picture to ``The Invalid.''

That in self-instruction which is of the Spirit, the ideal, which does not rest upon the human being, but must come to the human being's aid, when the personal is to be furthered, is, in these two pictures, made visible from two different sides.

For ``The Student,'' the Tutor shows himself outwardly, under the tedious disguise of everyday prose. ``The Student'' knows very well, when she wishes to know it, that the personal has a claim upon her most serious reflection; ``The Student'' knows very well, when she wishes to know it, that the personal is something one cannot get \textit{par renommée} from a tutor, any more than one can read one's way to it in some new book in some new language. The only book in which its features can be read is reality's own book, the detailed book of ethical concerns. To read a little with her in this book, the Tutor comes and offers himself, but, since he looks just as dry and long-winded as reality itself, she turns away from him, and so does not come to know his inward worth, his ideal riches.

For ``The Invalid,'' the Tutor is not visible at all. Here the relation is so arranged, that it is not the human being who seeks the ideal, but, conversely, the ideal that seeks the human being.

It holds true here, ``to work, while it is day, for a night comes, when no one can work.'' The old man, who rests his thin hands
% --- p. 112 ---
upon the great book, has ended his labor. Through this labor he has become neither hero nor martyr, but simply and plainly an everyday human being. Yet a human being he is, nonetheless, a real human being, who, in existence, has tried something, and learned something personal. He knows what it is to stumble, what it is to grieve, what it is to repent, what it is to hope. It is the personal yield of his life that accompanies him in the last hour; and then inspiriting is given him.

And in the last hour, when the body is so weak, when the soul is so faint, yes, in the last hour, then a human being needs help of every kind. Then a personal help is sought from the Highest, from the ideal itself: we cannot be helped by a mere reckoning; we need, unconditionally, the Highest.

This help is presupposed, but not pictorially represented; to represent such a thing under a figurative, self-chosen manner of speaking would, after all, be a mystification, a violation of the Most Holy. It is the human form of inspiriting that is pictorially indicated; for it is, after all, not the inspired one, but the ordinary human being, whom we here consider.

For the ordinary human being, the help from the Highest is like a ray of the Spirit's light, that clarifies the model, and in which the helper himself is present, and strengthens the frail one. But the ray fades again, and then it is so dark, and --- in the last hour! For the ordinary human being, in the last hour, there are so many dark moments. The last night, with the darkness of the last hour, can seem so long, so long, to the sufferer, longer than otherwise the longest lifetime.

Well is it for us that here is still a help at the very last, a help for the one who has striven personally, for the one who has used his time, and steadily applied the model to his real life, is surely not left helpless in the last hour. The eternal is not sensuous, and has no sensuous compassion for flesh and blood; but the eternal is nonetheless so human! Here is an inscrutable mercy, that takes up the frail one,
% --- p. 113 ---
and makes itself known in the purely human. This purely human is, then, the reckoning within inspiriting.

It is such an inspiriting that is represented in the picture; and that is why a purely human symbol has been chosen here as well.

That the old man does not see the figure is meant to signify that a finite remembrance, in its finitude, is not the help; the merely human is, after all, not the ideal; the merely human must not stand in the light and cast a shadow over the eternal.

The old man does not see the figure, but he senses the inspiriting; he is too weak to seek the help himself, so the helper comes to him. He is worn out by the long journey, and yet he is refreshed, and receives a lesson. His way led over the roaring sea: now the shapes of time fade from his sight, as the waves fade; but out of all that is vanishing, the personal rises up, which does not vanish. It comes as though out of what has vanished, it comes in ideal refreshment, it comes as if in a last breath from the roaring sea.

Here is something that can serve for personal instruction for us as well.

Had this old man performed great, brilliant deeds, and won rank with the most celebrated heroes: that would have helped him to a memorial in history; but even the most celebrated name would nonetheless not have saved him from the reckoning of inwardness, all the memorials of world history would nonetheless not have helped him personally in the last hour.
% First appearance of the "personally in the last hour" refrain that recurs
% (with variations) through the rest of this lecture (pp.113--117), mirroring
% transcription.tex's finding: checked individually at every occurrence in
% the Danish, none are letterspaced anywhere in pp.105--117 (unlike the
% different refrain in Lecture IX, which WAS letterspaced at its first/last
% occurrence). No \emph{} applied to any instance here, matching the source.
And had this old man, who mourned the fatherly grief over the son, ``the son who fell on the path of honor,'' had a quite different, great and bitter grief, the grief that his own sons had become traitors against the state; and had he defied this grief, and stifled the human feeling, and challenged justice, as a Junius Brutus stifled fatherly compassion, challenged justice, and stood --- an unbending spectator! --- at his sons' place of execution: that might perhaps have awakened the admiration of those who idolize the universal, it might perhaps have given him a name like a Brutus,
% --- p. 114 ---
a memory like a Brutus; but it would not have saved him from the sting of the self: it would not have helped him personally in the last hour.

And had this old man, in his vanished days, been the witty, the gifted, the happily endowed individuality; and had he then lived only for noble enjoyment, and practiced the fine art of measuring out enjoyment in every moment by the measure of shrewdness, prudence, and beauty: that would have helped him to much wistfulness in manifold memories; but it would not have saved him from the earnestness of eternity, it would not have helped him personally in the last hour.

And had this old man, in the years of his strength, had a stirred spirit, filled with surging, brilliant passions; and had he followed these passions boldly through anguish and joy, and had he felt a pain that only few feel, and a joy that only few know, the pain in the wild rending-apart, the inward jubilation in the bliss of the moment: that would have enriched him with manifold adventures, it might perhaps have attached strange legends to his name; but it would not have saved him from the wrath of the deep: it would not have helped him personally in the last hour.

Just as it is a lack of aesthetic fantasy, when the eye sees hazy shapes in the air, without any image taking on the form of reality; so it is a lack of ethical fantasy, when the individual cannot see the eternal in quiet inwardness, in ordinary reality. Those who talk much about the grandeur of grand passions and passionate catastrophes readily imagine themselves to be unusual, brilliant beings of fantasy, for whose sake a higher, special ethics must be invented; and yet that sort of human being not seldom betrays a lack of even the most modest ethical fantasy.

What is extraordinary in the phenomenon is, naturally, always conspicuous; but those who do not dare to count themselves among the extraordinary human beings must, after all, have the greatest interest in making clear to themselves the religious-ethical effects at work in the ordinary.
% --- p. 115 ---
If we turn again to the old man, then the meaning of his life lies, indeed, in this, that on ``the long journey'' he has directed himself by the model, the ideal helmsman on the roaring sea. Had he, however, merely in complete generality looked to the ideal, then his way of thinking might well have taken on a certain tinge of piety; but it would not have become earnest with the personal: the distance between the strong and the frail is, for the human being, immeasurable.

Had he, on the other hand, entrusted himself to his natural temperament, and steered by the compass of reason; then he would have missed personal truth. The compass of reason has just as many personal deviations as sensuous reality has rocks and reefs. These individual deviations cannot be corrected according to general concepts of reason; they can be corrected only according to the ideal itself. But, in order to correct accurately, a good reckoner is needed: such a reckoner is the Tutor.

The Tutor has a sharp eye; he looks into the pupil's inward being, and calls upon him to lay out the account exactly as it, in reality, is found to be, so that it may be revised according to the ideal's demand.

The Tutor has a sure eye for measure. He measures the distance between the ideal and this finite, petty, seemingly insignificant reality; he looks closely, that no point in this reality, which can in any way have personal significance, shall vanish as meaningless.

Joining the ideal with the real, the Tutor makes his notations everywhere in the account. Is there here a point that has been misunderstood, or a point that has been overlooked through the pupil's negligence, then the teacher at once sets his ``guilty'' beside this point. If the fault stands to be amended, an express notabene of obligation is added, in every case a mark of repentance.

Thus, we assume, the old man must have proceeded, by ideal instruction in self-instruction; only thus
% --- p. 116 ---
has he, through the insignificant deeds of his life, been able to win personal profit; only thus has he, under continued exertion with his deeds, been able rightly to learn to grasp the nothingness of these deeds, and to bow before the ideal in faith and humility.

Had the old man not thus made use of the Tutor, while it was still time, then inspiriting (if it had come at all) would have become dreadful to him in the last hour; it would then have gathered and compressed the ideal impression of personal truth into unspeakable anguish, into indescribable repentance.

Now, we do, indeed, have an example, that repentance in the last hour can be so strong, that it restores everything in a moment; it is a gleam of truth in the utmost darkness; it is the mercy of the Eternal One: ``today shalt thou be with me in Paradise.''

But that, after all, is the Savior's word to the repentant thief; it is, then, the ideal's utmost measure that we come to apply immediately, while our task was really to keep as close as possible to the ordinary. If the thief is to go to Paradise this very day, then his repentance must be extraordinary; for a thief is, after all, an extraordinary criminal. But the old man we here consider is, after all, no thief, no extraordinary criminal; we have him, indeed, precisely as a picture of the ordinary human being.

To live along in the everyday and regard the ordinary as a pretext for religious sluggishness and ethical slackness, and then now and then, on some solemn occasion, when the reminder comes, to trust to the extraordinary, and swing oneself up into the heights, is a pitiable self-delusion, a dangerous game with the personal.

To a life of ordinary, but nonetheless earnest, striving there corresponds, most nearly, the ordinary help in the last hour.

And in the last hour, when the body is so weak, so weak; and in the last hour, when the soul is so faint, so
% --- p. 117 ---
faint; yes, in the last hour, which can have such dark moments, it is, after all, nonetheless dreadful for the ordinary human being to have to receive inspiriting in unspeakable anguish, in indescribable repentance.

So it is, after all, far better to do what this old man is assumed to have done: to count the hours, still while they run, to seek the model, still while we are on the way, and to make use of the Tutor, --- while it is still before the last hour.

%% ============================================================
%%  XI. THE CONDITIONS OF PERSONALITY (printed pp. 118--131)
%% ============================================================
\chapter*{XI. The Conditions of Personality.}
\addcontentsline{toc}{chapter}{XI. The Conditions of Personality}
\label{ch:forel11}
\markboth{XI. The Conditions of Personality}{}

If we could hold fast to a straightforward separation of the purely religious and the purely worldly, then the question of the relation of ordinary human beings to the extraordinary ones would, in general, soon be answerable.

In a worldly sense we could then easily distinguish between the distinguished, the significant, and the insignificant, and in a religious sense we could, indeed, say: before God this distinction does not hold at all; before God we are all equal, that is: equally significant and equally insignificant.

But the religious and the worldly side do not let themselves be thus separated in a human being's personality. Spirit and nature are, indeed, opposed to one another, but nonetheless in such a way that the natural disposition furnishes the foundation and condition for the working of spirit. Insofar as the working of spirit is subject to particular conditions (individual natural conditions among them), we are already (even within the religious sphere) occasioned to distinguish between significant and insignificant characters. If, now, the divine freedom is also acknowledged, whereby the outbreak of spirit in different human beings at different times is determined; then we have, indeed, the distinction between the inspired and those who receive inspiriting only in the ordinary.
% --- p. 119 ---
That these different conditions of personality must belong to the divine plan can be seen from the ideal man's life and conduct.

He who knew what is in man says, indeed, to the one: ``Arise, and follow me!'' but to the other: ``Go home in peace to thy house!'' The one who was to go home in peace to his house was, surely, nonetheless also to appropriate the personal. Here one cannot think of an opposition between the elect and the rejected (the rejected could not, after all, go home in peace); here one must clearly think of those who are called in an extraordinary way to the extraordinary, and those who are, more nearly, appointed to hold fast to the ordinary.

Could we, even so, hold fast this separation straightforwardly, and, along with it, point to a fixed rule, a visible mark of distinction, according to which the two classes divide, then the matter would still not be so complicated; but with such a rule, such a mark of distinction, we stand again in the very midst of the Catholic Church.

In the Catholic Church the clergy (with monks and nuns) constitute, indeed, the class of the extraordinary; laypeople have, so to speak, a privilege of remaining with the ordinary. With this it hangs together, too, that the extraordinary can perform more good works than they themselves have need of for their own salvation, and that the Church can, in turn, grant the ordinary the privilege of enjoying the benefit of those others' superfluous works.

But such privileges do not hold with personal truth. Personal truth turns the relation, again, in such a way that the last become the first, the weakest the strongest, the lowly and despised the most exalted. The personal spirit distributes its gifts to everyone, according as it wills; this spirit is thus governed neither by the mere natural disposition nor by a merely outward order: it calls whom it will, when it will, and in whatever way it will. All must, then, each according to his ability, work with the utmost
% --- p. 120 ---
exertion, and hold themselves in readiness; for anyone, after all, can receive an extraordinary calling.

If, then, we are to arrive at any insight as regards this point, we must turn our attention to the difference of inwardness between the inspired and the ordinary, and to the intermediate stage that shows itself between the two.

In order to indicate a couple of definite points of view, let us, then, first consider:

\begin{center}
\textbf{The Breakthrough of Truth in the Personal.}
\end{center}
% Mid-lecture display sub-heading, mirroring transcription.tex's finding:
% set larger than body text and centred, not letterspaced. A second
% instance of the same pattern follows below (p.125, "Conscience.").

If we begin from the height, we have, indeed, the Apostle as truth's extraordinary messenger. Thus Paul is, then, a striking example of the extraordinary calling.

The young Pharisee begins by being zealous for the Mosaic Law; he sees in the crucified one an enemy of the Law and of truth, who received his deserved punishment, a criminal whose followers it is God's own work to root out. He sat by, when the first martyr was stoned; when they stoned Stephen, ``the witnesses laid down their clothes at a young man's feet, whose name was Saul.''

And it was not enough for Saul, that he ``harried the congregation at Jerusalem, and haled both men and women, and committed them to prison;'' no, he had to go further, further still: he breathed out threats and slaughter against the Lord's disciples, and went to the high priest, and desired of him letters to Damascus, to the synagogues, that, if he found any men or women who were of this way, he might bring them bound to Jerusalem.''
% Printer's defect (carried over from transcription.tex): this Acts 9
% citation closes after "prison;''" and then resumes with "no, he had to go
% further..." without a new opening quotation mark before the resumption,
% leaving the final "to Jerusalem.''" as an unmatched close. Reproduced here
% with the same unmatched close, mirroring the source.

Here is an extraordinary blindness, a personal raging against personal truth; and yet this raging one believes himself to be serving truth. Here is a man, who, in burning zeal for truth, rises up against truth, and would kick against the pricks: this, then, is the extraordinary. The prick is already in his inward being; truth has mercy upon him: the mercy is extraordinary, as the blindness was.
% --- p. 121 ---
The blinded one must first become so blind, that he sees his own blindness; then he becomes blind by the miracle: the light that shines about him from heaven makes him blind. --- ``He saw not for three days, and he neither ate nor drank.'' For three days! Who dares describe what can pass within such a man in such three days? Blind by a miracle, he became seeing again by a miracle. Such a breakthrough, such a calling, is, in the strictest sense, extraordinary.

We can now descend lower, and still have a height, when we look up to a Martin Luther.

The natural conditions of finitude in the merely human step forth more strongly here. It is not a snorting Pharisee, bent on choking truth in blood, not a wildly storming one, halted by a miracle; it is a young, modest master of arts, who was, indeed, by his father's express wish, soon to be a lawyer. But --- this master does not become a lawyer: he suddenly becomes, after all, afraid of himself, and afraid of the world, and flees by night into a monastery.

What does this fear mean, this flight? The young fugitive is so pale; one would think it was his last hour. He is like a forsaken one, a God-forsaken one, a dying man in fear of death. Yet he shall not die; no, he shall live! This fear means precisely that he is about to live, a life that is a life, an eternal life in the personal.

It costs anguish and struggle, when the personal is to come forth in earnest. If anyone would rightly consider what struggle it has cost here; then he should not hasten to go to Worms to hear the bold monk. In Worms, where the emperor, and where a whole world, already listens for the monk's word, yes, here in Worms, here the phenomenon might perhaps dazzle; history has already begun its drama here. No, it is far safer, in order to begin at the beginning, to go to the Augustinian monastery in Erfurt, and look through the window into one of those dark cells.

Here lies Brother Martin upon his face, unknown to the world and forsaken by the world: his weakness comes of much
% --- p. 122 ---
fasting. His door is shut; that door has already, for several days, been shut: the fasting man has forgotten to unlock the door, and can now not hear those who stand and knock.

Yet, all the same, things go humanly enough in this monastery. His anxious friends surround Brother Martin, and rouse him with the sound of the lute. Then he receives help of every kind; then the Tutor comes and spells out the letters with him; then he spells his way through the Gospel, and behold, he spells his way, in the end, to --- personal truth.

Luther is no apostle, and yet, what a height, when we look up to Luther!

But now the ordinary human being: how does it stand here with the personal turning point?

In everyday life the tension between the worldly and the holy cannot, after all, step forth so grandly. The Methodists' zeal for a conspicuous convulsion of penitence takes on the appearance of an enthusiastic excitement, a fanatical hunting after the extraordinary.
% Doubtful reading (carried over from transcription.tex): the Danish
% "Ivren" is printed with its medial e missing (expected full form
% "Iveren," "the zeal"); transcribed there as printed. Untranslatable as a
% spelling defect -- "zeal" above renders the word normally.

Under the general everyday conditions, one ordinarily begins with a religious-moral upbringing according to the guidance prevailing in Church and society. The first thing is, that the minor is gradually given the necessary knowledge of life's highest and most important concerns, and senses that there is here something called authority. The self-willed, arbitrary individuality must be bent into obedience under the law: the one who has not learned to bend never comes of age.

As long as the individual is in happy agreement with what has been handed down, no break is felt. The inward, the original, has not yet expressed itself: the germ of personal independence still lies as though dormant. If a breakthrough is to take place, there must come a collision between that which the individual carries within himself, and that which is offered in what has been handed down. Everyday life, too, has its hidden
% --- p. 123 ---
conflicts; here daily occur many stronger and weaker frictions.

Everything now depends on how deep the break goes, and how inwardly it is grounded. That the immature one criticizes and levels, that the minor presumes to judge what he does not understand, that the ill-bred one insolently breaks with his surroundings, presents only a sorry picture of the personal. It must always be presupposed, here, that it has cost the individual an effort of will to break. The greater the struggle, the more effort it has cost, the more inward and pure the break. For it is not only crystals that are known by their fracture; a human being's character, too, bears the mark of its original break; but natures differ.

The dry, but sensible nature generally breaks in such a way that the person, in all composure, learns to hold to himself, be his own counselor, and lay out his path according to certain unwavering principles. In such principles the will can, indeed, lay a certain firmness; but in mere principles, were they even nothing but orthodox dogmas, there is no idea, no inwardness, no depth: if inwardness is to come, further work must be done here.

In the hard, but deeper, nature the break proceeds differently. Such a human being's personal surface can, under the pressure of reality, become like a stone, and yet the inwardness can stir in the depths. On some single occasion this inwardness can even press so strongly that the petrification cracks, the hard man can weep: his tears are like drops that trickle down a cliff-face. Here is a good kernel; but if the shell is not to spoil it, further work must be done here.

Different again is the break in the soft, more fanciful nature. Here come moments, in which the individual, without being able to account for it to himself, breathes in melancholy, and feels himself so alone, so forsaken. Here is mood and fantasy-inwardness; it looks profound and romantic: here one must be careful.
% --- p. 124 ---
A young man sits beneath the shade of the tree; he looks down into the brook, so sorrowful, so grave, as though he came from a burial. ``Has he, then, some grief; has he crape upon his hat; are his parents, perhaps, dead to him?'' By no means! He knows, as yet, but little of real grief. Nor does he have crape upon his hat, either; he has, indeed, field violets and a flower in his breast: joy and happiness are in his home. ``Why, then, does he sit so sorrowful?'' Let us not disturb him, he does not know it himself. It is a moment of forsakenness, and there is, then, something trickling in there beneath the flower, as the brook trickles. ``What, I wonder, will it become?'' Well, that one does not know, for it can become anything at all. If it is only a little poet, it will surely become something very languishing; but it can, indeed, also become something quite different, when further work is done.

Eternity holds every human being responsible for how he uses every moment of inwardness; this responsibility is so much the greater, in that no human being, in such a moment, is or dares be entirely alone. Here is an invisible one, who looks upon the personal: the Tutor waits for the account to be laid out, and the ambiguous to be clarified.

In the picture book, ``The Forsaken One'' is, therefore, also represented in the moment, when a half-languishing fantasy-inwardness is in the process of appropriating the religious. It is a dangerous moment, for, if what is fermenting is not clarified, ``The Forsaken One'' can easily become romantic in more than one respect.

She might, indeed, become romantic on a Sunday afternoon, go out into nature and sit down beside the young man, and look down into the brook. Then we should, indeed, get something very languishing, but the personal we should scarcely get. Or she might, on a Sunday
% --- p. 125 ---
afternoon, become even more romantic, read the tear-jerking novel, go off to a pious gathering, and fall in love with the heavenly bridegroom: then we should, indeed, get something exceedingly languishing, we should, indeed, get the mystical, the pietistic; but the personal we should scarcely get.

So that this shall not happen, the Tutor comes, and offers his guidance. Under a disguise of fantasy, he begins, in the first lesson, by clarifying her fantasy. It then appears, that this piety, this devotion to God, is not, precisely, so entirely pure: here is a bit of self-love, a little ambition, a little worldliness, and so it is, indeed, good that she goes to confession, so that the dissimilar elements may be separated out. As for what is truly ideal, the Spirit has gifts enough, but for personal use.

Every human being has a disposition toward inwardness, and moments for inwardness; everything depends upon the exertion: here is no privilege for the ordinary.

As the next main point, I name

\begin{center}
\textbf{Conscience.}
\end{center}

What is conscience? Well, that is now a rather old-fashioned word. In the newer, more modern language, one prefers to use the word: consciousness. Of consciousnesses, again, there are different kinds, such as: self-consciousness, world-consciousness, God-consciousness, and so on. Those who have an altogether pure, transparent, and unclouded self-consciousness have even already come so far, that they can dispense with all God-consciousness without missing it.

When these pure self-consciousnesses now and then condescend to use the word conscience, we must, above all, not be so old-fashioned as to understand thereby any relation to God in the religious sense. ``Conscience demands the ethical: nothing further!'' --- ``The ethical is in human nature: nothing further! Roundly, nothing further!''
% --- p. 126 ---
Should this cocksure self-consciousness now really be the height of the ethical? Then I must give myself up as lost; for I simply cannot get beyond the notion, that, as far as conscience, moral earnestness, and such things are concerned, the apostolic must be the height.

An apostle is, surely, a man who can sacrifice himself for the ethical idea. When a Paul preaches, he does not preach with words: he preaches with his life; when a Paul chastises, he does not chastise with words: he chastises with his life; when a Paul comforts, he does not comfort with words: he comforts with his life.

``We are fools for Christ's sake, but ye are wise . . we are weak, but ye are strong. Ye are honorable, but we are despised. Even unto this present hour we both hunger and thirst, and are naked, and are buffeted, and have no certain dwelling place; and labor, working with our own hands. Being reviled, we bless; being persecuted, we suffer it; being defamed, we entreat; we are made as the filth of the world, the offscouring of all things unto this day. I write not these things to shame you, but as my beloved children I warn you.''

Here is personal superiority. If we ask this wonderful man, why he endures all this, why he thus exerts himself to the utmost, the answer sounds simple and short: in order to be able to have an untainted \emph{conscience before God and men}.
% Letterspacing (Sperrsatz) confirmed in transcription.tex by zoom: this
% phrase, the answer to the lecture's rhetorical question, is set off as
% the emphasised phrase in the Danish.

Before God! Yes, that is, after all, to him the one thing; for it is to him, for the rest, ``a very small thing to be judged by man's judgment, he does not even judge himself''; but --- before God! Such a conscience cannot, then, possibly transform itself into pure self-consciousness.

Luther was, indeed, no apostle, but nonetheless a character who could do something for the ethical. He was a man who knew life's joy and life's pain; he had, after all, self-feeling too, and a sense of freedom, and cheerful courage. He was a man who could comfort himself, and defy a whole world, a man who dared what few dare --- for the sake of the ethical; but, it goes without saying, there was, indeed, certainly something he did not dare venture: he did not dare to tear himself loose from God; he did not dare to transform his conscience into pure self-consciousness.

But --- this, too, must not be forgotten --- Luther was living, after all, still only in his sixteenth century, while we, indeed, live in our
% --- p. 127 ---
nineteenth: should we, ordinary human beings of the nineteenth century, perhaps be able, without further ado, to venture what the great man of the sixteenth held to be the most terrible of all? There is required, surely, even still in the nineteenth century, a quite extraordinary human being, in order to make God himself entirely superfluous.

Without the help of revelation it is impossible for a human being to arrive at true personal self-knowledge. Revelation corresponds precisely to its concept, insofar as it makes me ethically manifest before God and to myself, in that it makes manifest the eternal in personality.

Such a revealing does not happen, in the ordinary, by miraculous visions; for the visions by which the Pietists so readily edify themselves seem, indeed, not so much to be supernatural visions, as rather unnatural confusions of the dissimilar elements of inwardness.

The revealing activity takes place, as the human being, before the ideal's gaze, tests himself in his conscience.

If anyone is alone with his self-examination, by himself, without the impression of a co-knowing personal truth: then he cannot be raised above himself; but the one who cannot be raised above himself in the decisive moment can also not come to know the truth, least of all, when this truth is personally wounding.

To present this, we use the Tutor; it is the second lesson in the picture of ``The Forsaken One'' that I can here refer to.

The two sides, from which the emotion proceeds, are hereby clearly distinguished. The person in question has, after all, erred, and received a reprimand. A reprimand! Yes, so weak is ``our ethical consciousness,'' that it cannot, even in ordinary matters, correct itself properly, but needs a reprimand. Even those who, in ethics, have made God superfluous have, nonetheless, surely not yet made the reprimand superfluous.
% --- p. 128 ---
Is not our age precisely the age of reprimands? From private insignificance to public greatness, from the hearth-corner to the battlefield, from the cottage to the castle, from the nursery to the council chamber, from the nursemaid to the prime minister, the reprimand stretches, indeed, with ever-increasing emphasis, the public one far stronger than the private. If pure self-consciousness were really so sure of itself, what need would there be, then, of the many reprimands?

That an unjust reprimand must miss its effect is natural. But even the fairest, most deserved reprimand always produces a mixed, ambiguous effect. We are, after all, by nature, one and all, egoists; we can freely confess it: the pure self-consciousnesses will certainly not contradict us on this chapter.

That a reprimand now makes my fault, my transgression, clear to me, that another plainly tells me what I have neglected to tell myself, is undeniably something good; if I am a reasonable and otherwise more or less moral human being, I must, surely, admit this.

But now, with all this, I am also an egoist, so there is, then, no one who can reprimand me effectively, without humbling me, and wounding my egoism.

But what, now, is the one who reprimands me? ``A human being! So, then, surely also an egoist. The one reprimanding may perhaps be right on this one particular point; but then he is, in turn, wrong on so many other points. It is, after all, really hard, that the one human being should let himself be thus humbled by the other.'' Is it not, now, under such circumstances, quite natural, that the wounded egoism thinks of some small way out?

Yet, let us assume that the reprimand produces the intended effect: what follows, then? Self-reproach! The nearest, involuntary object of self-reproach is, then, the fault in the particular case in question. ``So I have, then, done wrong! in this matter I have erred; it must, then, be made good again, if it is possible.''
% --- p. 129 ---
A moral fear, an involuntary anxiety about something unknown, is, in the nobler nature, now surely at work along with such self-reproach; but, so long as the fault remains within the ordinary, self-reproach will cling to the particular, and lose itself in finite worry: how one can make it good again, how one can best get free of it, and so on.

But the ethical idea does not exist in such a parceling-out. The idea is not my right and your right; the idea, as idea, is justice itself, the unconditioned, according to which your right and mine are, in reality, to be determined.

If self-reproach is to have true, moral significance, it must reach the idea as idea; but the idea as idea it cannot reach through the merely ordinary.

An ordinary neglect of duty, a fault in some small matter, does not bring the human being left to himself to repent so deeply. ``The Forsaken One,'' who has received the reprimand, and ferments in self-reproach, is by no means prepared to go so deep: it comes to her, on the contrary, quite unexpectedly, that the Tutor is sitting right there in the chamber, waiting for her.

If one perhaps wishes to despise this ``insignificant one,'' and appeal to the genuine philosophical thinker, who has ethics in pure self-consciousness; then I answer, that to the genuine philosophical thinker, who has ethics in pure self-consciousness, nothing is so unphilosophical as self-reproach and repentance.

If the idea is to come as idea in repentance, it must break forth involuntarily; but if it is to break forth involuntarily, it must, beforehand, have been provoked by the crime. The deep anxiety then seizes the imagination: the idea comes, then, as fantasy's guest, as the Commander, when he knocks.

Yes, the Commander is a guest of fantasy, he is like a messenger of conscience from the grave. The dread and terror that announce him, and follow him, are the dread and terror of fantasy. He is a guest who does not show himself by the light of day; he comes only by the flash of lightning in the night.
% --- p. 130 ---
And yet the idea, even when it thus passes through fantasy, is still by no means certain of mastering egoism. The understanding can, after all, reason against fantasy's images of terror. The criminal can defy fantasy's guest and kick against the pricks. The understanding looks only to the real; the understanding is not, after all, afraid in the dark: it has not ``learned to dread the play of pale shadows.''

The criminal, who will not dread the shadows of imagination, is, at bottom, right; if he really holds himself convinced that there is, in the idea, no eternal will, then he must, after all, feel himself his own master, and see to mastering his own fantasy.

If the idea is to master the human being personally, it must itself be recognized in its personality. The involuntary movements of temperament vary with the natural differences of individuals; pure justice is only a theoretical concept: only with the recognition of the just one himself is the essence of conscience revealed.

If God is recognized in conscience, then the extraordinary crime is not needed, for fantasy's thunderstorm to shake and cleanse the inward being: he who reveals himself in conscience leads even the most ordinary human being into the depths of self-knowledge and repentance.

In everyday language the admonition runs, indeed: ``For God's sake!'' --- ``For God's sake, man, consider what you are doing!'' --- ``For God's sake, man, go into yourself, go and confess!'' With this admonition, then, the Tutor comes too, and joins the ideal with the real.

In closing, only one more point. If the personal relation to the eternal is taken away, then the life of the world divides into two pieces: the one part is pure calm and superficial security, the other part is gnawing worry, anguish, and struggle, and unrest. The mixture can, in different individuals, be highly different; but this piecemeal division vanishes only before personal truth.
% --- p. 131 ---
Let us, then, not forget the conditions of personality: here is no one so lowly, that he cannot come to the Highest, and receive, in weakness, his strength; here is no one so exalted, that he need not bow before the Highest, and have his strength in weakness.

Let us not forget the conditions of personality: here is no one so calm, that he is not troubled by some secret unrest; here is no one so restless, that he cannot come to rest.

Let us not forget the conditions of personality: over life's calm hovers a thundercloud of judgment, over the temporal struggle shines an eternal peace, like the star over a battlefield.

Let us not forget the conditions of personality: there were chosen individualities, whom the Eternal One himself planted so deep in the ground of history, that they should be like the tall trunks whose crowns roar in the storm; but here, too, was surely a rich multitude of humbler plants, which, sheltered beneath those others, grew a quiet growth; for the gardener ``came and went, indeed, without anyone noticing anything.''

%% ============================================================
%%  XII. PERSONAL STRIVING: A CONCLUSION (printed pp. 132--144)
%% ============================================================
\chapter*{XII. Personal Striving: a Conclusion.}
\addcontentsline{toc}{chapter}{XII. Personal Striving: a Conclusion}
\label{ch:forel12}
\markboth{XII. Personal Striving: a Conclusion}{}

``If there is any question that can rightly be called a question of life, it must surely be the question of the eternal life; if there is any cognition of which it indisputably holds that it concerns every human being personally, it must surely be the cognition of personal truth; if every human being in the race is truly disposed toward true personal development, then true personality must indeed be the prototype of each individual human being and, consequently, of the whole race.''

With these words we opened these lectures, and to this opening we now return, as to a conclusion.

Is there, then, really an eternal life, and is the eternal life our personal goal? We do not know it; there is no human being on earth who can know it. If anyone will seriously take the personal to heart, he shall surely learn how true, how true it is, that faith stands against knowledge, and knowledge against faith.
% [doubtful reading, carried over from transcription.tex: the Danish
% "Dersom Nogen i selv vil" is transcribed exactly as printed, flagged
% there as a likely printer's dropped "kke" from an intended "ikke"
% ("not"). Untranslatable as a spelling defect; the English above renders
% the evident sense ("if anyone will...") rather than the garbled surface
% form.]

But faith must, after all, have its ground. Wherein, then, shall we seek faith's true ground? In our infinite need for the life of the spirit, the life that is in the ideal. Has He, who is our model, been mistaken about the personal; has He, who
% --- p. 133 ---
says: I am the Truth, nonetheless, after all, been mistaken about truth, and squandered his life on an enthusiast's delusion, a fantastic exaggeration; then, indeed, his invitation: ``\emph{Come unto me!} is a false invitation, and then faith is --- shall I call it a happy or an unhappy error? It is quite immaterial. Happiness and unhappiness come with time and go with time; but the truth is, then, this, that in an eternal sense there is no personal truth.
% Printer's defect (carried over from transcription.tex): the opening
% quotation mark before ``Come unto me!'' (Matt. 11:28) is never closed --
% the sentence runs straight on into "is a false invitation" with no
% closing mark anywhere before the next full stop. Reproduced here with
% the same unmatched opening quotation, mirroring the source. "Come unto
% me!" is letterspaced (Sperrsatz) in the Danish, confirmed by zoom --
% rendered \emph{}.

But --- I believe --- there is, surely, no danger in this. I must, after all, believe, that He who said: I am the Truth, has known what truth is, and known, with eternal certainty, what he said. I must, after all, believe, that He who said: I am the Way, has known whence he came, and whither he went, and has not been mistaken about the way. I must, after all, believe, that He who said: I am the Life, and sacrificed his life for the sake of truth and to show the way, has had the source within himself, the source that wells up unto an eternal life.

No, from this side there is, surely, no danger here; but there can, therefore, indeed, be danger enough from another side. There can, indeed, be danger, that I, who am not the truth, shall misjudge the truth, that I, who am not the way, shall stray from the way, that I, who do not have the source within myself, shall not find the source of life.

It is now repeated, loudly and clearly, from various sides, that the clearer insight of the newer age into the laws of nature and reality makes it more and more difficult to reconcile the conceptions of revelation-faith with the rational concepts of science.

That I, for my part, belong among those who grant these difficulties to the very fullest extent, I have honestly and expressly acknowledged.

Here, in earnest, there can no longer be any talk of the difficulties themselves; the question of life presses, and the question of life is, what conclusions can here be drawn from the presupposition of the granted difficulties.
% --- p. 134 ---
One can draw from this the conclusion, that revelation \emph{is an untruth, and faith a self-deception}.

One can draw from this the conclusion, that revelation \emph{is the eternal truth, exalted above all human comprehension, and faith our highest task}.
% Both predicate clauses above are letterspaced (Sperrsatz) in the Danish,
% confirmed by zoom -- the two rival conclusions are set off from each
% other typographically.

Which of these conclusions each one will prefer rests upon himself: no reasoning helps here. What we, on the other hand, insofar as there can be talk here, in this matter, of general enlightenment, must strive to make quite evident to ourselves, is this, that whichever conclusion anyone prefers, he can by no means justify it scientifically.

Unbelief's conclusion rests, just as surely, upon an unprovable assertion, as does faith's conclusion. The conclusion is, here, a personal choice, and for this his personal choice every human being shall answer to truth itself.

When it now stands firm --- and it shall, indeed, stand firm, however one twists and turns --- that the choice between faith and unbelief is personality's ``Either --- Or:'' then it is no longer the high-flown question of the general state of the age and the world and science that is at stake, but a state within the human being himself.

As there is, in the human being himself, something that at any time and under the most various conditions can awaken the need for faith, so, too, is there, in the human being himself, something that at any time and under the very most various surroundings can awaken, nourish, and strengthen unbelief.

Corrupt a human being's fantasy thoroughly, if you can, and you have corrupted the whole human being; mislead a human being's fantasy, and you have misled the whole human being; give the human being's imagination a thoroughly sensuous, material, earthly-worldly direction, and you may be assured, that the human being, though he were given the richest knowledge, the clearest concepts, the finest culture, will, indeed, bow before material reality, and go the way of worldliness.
% --- p. 135 ---
The kingdom of nature is rightly called the kingdom of the five senses. Teach a human being thoroughly to esteem the testimony of the five senses higher than the testimony of spirit and conscience, and you need not even initiate him into the coherence of the laws, or burden him with the prolixities of natural science, to be sure that the person concerned will, indeed, of himself sense the ``coherence of the laws,'' and know how to be ``Nature.''

The question, whether I am to set the certainty of the senses higher than the demand of the spirit, can be decided even less by natural science than by pure logic. The decision is laid in the personal choice.

If it is a fantasy, through which a human being's feelings, thoughts, and will are directed toward the eternal, then it is also a fantasy, through which the human being receives the picture of the world corresponding to the sense-impressions. All forming of the ideal is a forming in the imagination; all error in the ideal is, therefore, also an error in the imagination.

It is a false imagining, that leads the sage of the East, when he, despising the temporal, denies his individuality, in order to behold the pure light, and lose himself in the infinite void.

It was a false imagining, that led the Greek of antiquity, when he deified the beautiful individuality, and fashioned for himself a harmony of virtue, in order to perfect his self-formation, as one perfects a work of art.

It was a false imagining, that led the pilgrims of the Middle Ages, when they sought afar what was infinitely near to them, when they searched in the outward for that which is found only in the inward; when they, forgetting the laws of reality, fashioned for themselves a heaven of sensuous images and pious legends.

It is a false imagining, that beguiles the thinkers of the present day, when, for sheer races and individuals and \emph{individuals and races, they cannot catch sight of personal truth}.
% Letterspacing confirmed by zoom in the Danish: begins exactly at the
% second occurrence of "Individer" (the reversed repetition) and
% continues to the end of the sentence; the first "Slægter og Individer"
% is plain weight. \emph{} applied over the corresponding English span
% only.
% --- p. 136 ---
If the rational ones of today could call out to the pilgrims of the Middle Ages: ``Why seek ye the living in the empty tomb, why seek ye the living among the dead?'' then those pilgrims might, in turn, call out to these rational ones: ``Why do ye measure the spirit by matter; why do ye seek the personal self in a heart of flesh?''

It is a false imagining, that at all times beguiles human beings, so that they either reject the Highest, or misuse the Highest, and fashion for themselves a connection between here and yonder, which has no validity.

It is a false imagining, that beguiles human beings, when they drag the eternal down from its divine height and transform it into an emergency aid for the temporal.

Where that imagining gains a foothold, the imagining, corrupting in its pious ambiguity, that He who is personal truth should have come into the world in order to set a little right our temporal ailments, our physical pains, our worries over livelihood, our civil unrest, and so on, where this imagining once gains real foothold, and nourishes itself on lofty phrases, it looks doubtful indeed for the recognition of truth.

It is, indeed, true --- true in faith --- that the real ideal has revealed its personal power, and has itself pointed to its wonderful works: ``the blind see, the lame walk, the lepers are cleansed, the deaf hear, the dead are raised up, and the gospel is preached to the poor.'' Yes, that is help! But is it to mean temporal help in the temporal, sensuous pity with the sensuous?
% Checked for letterspacing by zoom in the Danish: no Sperrsatz on this
% Scripture citation (Matt. 11:5), unlike some earlier lectures'
% quotations -- confirms each occurrence must be checked individually.

The same one, who makes it manifest, that there is a will that masters the laws of nature, makes it also manifest, indeed, that the sensuous individual is to be sacrificed for the law of the spirit, that flesh and blood are not spared. Or was He, perhaps, spared? Were his messengers, perhaps, spared?

It is, indeed, true --- true in faith --- that these messengers, according to his own word, who sent them out, could ``take up serpents, could tread upon vipers and dragons, could drink deadly poison, and it should in no way hurt them.''
% --- p. 137 ---
Yes, here is assistance! But is it temporal assistance in the temporal? sensuous help with the sensuous? It is, after all, a promise that contains the strictest demand for self-sacrifice.

Must this now be recognized in the extraordinary, and does the extraordinary have precisely this significance, that it clarifies what is otherwise obscure and ambiguous in the ordinary: should I, the ordinary human being, then have leave to expect sensuous pity with the sensuous, to smuggle in the interests of my self-love through the religious, and seek the comfort of softness in ``the hidden miracles''?

It is promised to the believer, who is willing to grasp and follow the demand of the spirit, that he shall receive strength in weakness, endurance in patience, and courage and power to go through all things. But it is, so far as I know, nowhere promised to the believer, that he shall receive supernatural assistance, in order to come well through the temporal in a temporal way.

If it is an instrument of God, a messenger of truth, who is to go forth among human beings, then let them do with him what they will: cast him into the sea, throw him into the lions' den, hand him the cup of poison; he shall, nonetheless, be preserved, so long as the Spirit wills; but flesh and blood is not spared.

It is, indeed, true --- true in faith --- that we are not only to work in the temporal, but also to pray for the temporal; for the one who prays for \emph{our daily bread} prays, after all, for something temporal. But it is just as true --- true in faith --- that the temporal is, and ought to be, a means for the eternal.
% Letterspacing confirmed by zoom in the Danish on "det daglige Brød" --
% the Lord's Prayer phrase.

That the matter stands thus, we all know; and if the natural in us were now, without more ado, on this point, in agreement with the spiritual, no particular caution would be needed here. We could then hold, straightforwardly, to the promise: ``Ask, and ye shall receive'' --- not only the spiritual, but
% --- p. 138 ---
also the temporal; for there is, surely, One who can give us both temporal and spiritual goods in overflowing measure.

But here is an impropriety that dare not be overlooked. The natural human being is by no means straightforwardly in good agreement with the ideal's demand. Instead of grasping the temporal as a means, solely and only as a means for the eternal, egoism, with all its dexterity, turns the relation about, and behold, then the temporal becomes, after all, the goal, the nearest and the real goal.

The confusion does not happen straightforwardly with the knowledge of the one praying (for those who want only the temporal rightly consider it a waste of time to pray); no, the confusion happens in the deepest cunning. Here it is a matter of being on one's guard; for the cunning of egoism, when it smuggles itself into the religious, is even more contrived than that of those who commit customs fraud.

``Ye ask, and receive not, because ye ask amiss;'' this word of rebuke does not apply only to those who, in a worldly sense, are given over to ``pleasures''; it applies, essentially, to all who, confusing the means with the end, straightforwardly suppose that they can ask the Eternal One for extraordinary help in the temporal.

But can a human being, then, expect no temporal blessing at all from prayer, the prayer that, in all piety, is sent up to God for the happy progress of his labor?

Let us consider truth's extraordinary messenger, in order to arrive at the recognition of truth. Paul was, after all, undeniably an industrious and diligent man: in order not to be a burden to others, he even worked with his own hands; and, besides, he was, indeed, also a man who knew how to pray. If prayer could really, straightforwardly, be supposed to raise the wages of labor: then we should, after all, undeniably expect, that Paul had received some extraordinary supplement, that through his praying and his manual labor he had become even a well-to-do, a substantial man. Did it, now, go thus? Did this praying man become a well-to-do man; did he receive a supplement in the temporal? Yes, in ``hunger and thirst, in fastings, cold and watchings.''
% --- p. 139 ---
It is a false imagining, that teaches a human being to excuse himself with the ordinary, and settle into ease in the everyday, under the pretext, that the ideal is and must be \emph{unattainable}.
% Letterspacing confirmed by zoom in the Danish on this first occurrence
% of "uopnaaeligt." A second occurrence of the same word, inside the
% quoted dialogue below, was checked individually and is plain weight --
% not letterspaced. Confirms each occurrence must be checked separately,
% not assumed uniform.

This imagining is nourished chiefly by this, that one treats conversion as a matter apart, as something discussed at a certain place in dogmatic ethics, and then, for the rest, takes everything as straightforwardly as possible, forgetting, that all relations and all concepts must, indeed, at conversion, be turned about.

Our personal striving toward the ideal of personality cannot, after all, be a straightforward striving of approach, but a striving of conversion.

Before we go further here into the religious, I will try to illuminate the difference between a straightforward and a converted striving by an intelligible example.

That the art of music has its ideal, and that there is a striving of approach toward this ideal, we know. This striving is, despite all preliminary missteps, temporary setbacks, and delaying detours, essentially a straightforward striving.

The untalented one, who does not even have an ear for music, should not waste his time on any continued approach. The untalented one is, then, right, when he says: ``For me the ideal of music is unattainable, and therefore I do not strive at all.''

The talented one, who has the desire and the opportunity, ought to strive. Every stage he passes, every difficulty he overcomes, brings him nearer to the goal. If he is fully talented, then he will, indeed, under his striving, eventually reach a point, where the ideal, inexhaustible in itself, gives its content form in his actual production. Here we have, then, a master in the art of music: the goal is reached straightforwardly, by a straightforward striving.

Now, in order to make the converted striving recognizable, let us place a middling figure between the wholly untalented and the
% --- p. 140 ---
fully talented. This is, then, a human being, who is talented enough to be able to strive, and yet, besides, so little talented, that he can never satisfy himself by his striving. Here, too, is an approach, but an approach of a different kind.

Instead of approaching the ideal straightforwardly, this striving one will see himself more and more removed from the ideal, but nonetheless, under this removal, attain an ever purer insight into the essence of the ideal, an ever increasing enthusiasm for the glory of the ideal.

Such a striving is, now, surely, a folly; for the striving one never, indeed, in reality, comes to anything. Nevertheless, it is impossible, even for the wisest man, to stop this striving one, to bring this fool to reason. If we hear him himself, he explains to us his folly as the highest wisdom.

``That in reality I accomplish nothing and attain nothing, means nothing; I nonetheless, all the same, attain the highest thing a human being can attain here on earth, an ever purer insight into the essence of the ideal, an ever increasing love for the ideal itself, an ever growing joy in the glory of the ideal. For you must know, you human beings: I love the ideal for its own sake. You clever folk, who accomplish so much, must know: I believe in the ideal; I believe in a god of tones; I believe that there is, yonder in heaven, a master above all masters, a master in whom the ideal lives. Yes, when he comes, the heavenly master, then I shall rejoice, as no other; then I shall worship him, as no other; then I shall love him, and the ideal in him; then he shall grant me, the rich God, for my love's sake, what I lack in gift; then my diligence, my exertion, shall be infinitely repaid in my infinite joy.''
% spacing.py flagged several scattered short words inside this long
% quoted speech; checked by zoom against the image in transcription.tex
% -- all plain weight, no Sperrsatz. Noise, not emphasis.

If there now really were such a god of tones; if this god would then disguise himself as a shepherd, and descend
% --- p. 141 ---
from Olympus, --- not to listen to the music of human beings, but to let himself be heard, and thereby learn, whether there was anyone among mortals who loved the ideal for its own sake, --- and now saw, that the earthly masters loved the ideal only in their own productions: then he would, surely, pity them; and if anyone even dared to enter into a contest with him, then the angry god would, indeed, punish the presumptuous one, as Apollo punished the flute-player Marsyas; but that striving one, who had everything in conversion, that converted striving one, who recognized him under his disguise, and loved him for the ideal's sake, verily, this striving one the divine one would not reject, but grant him his grace, lead him with himself up to Olympus, consecrate him into the art of tones, and gladden his soul among the blessed gods.

So we have, then, an example. But for an ideal of art, the converted striving does not hold. Here it is a matter of either striving straightforwardly or stopping.

For the ideal of personality, by contrast, only the converted striving holds. The extraordinary messenger cannot, after all, even strive straightforwardly: the Apostle himself strives in conversion.

``\emph{Not as though I had already attained, either were already perfect: but I follow after, if that I may apprehend that for which also I am apprehended.}''
% Phil. 3:12-13. Letterspaced (Sperrsatz) throughout in the Danish,
% confirmed by zoom -- entire citation rendered \emph{}.

The more zealously he strives, the more works he does, the more inwardly does he feel himself the unprofitable servant, who receives everything of grace. On these terms he holds out ``in labor and toil, often in watchings, in hunger and thirst, in fastings often, in cold and nakedness.'' So he does not become, straightforwardly, stronger and stronger; no, he becomes frailer and frailer, and then, when he is frail, then he is \emph{mighty}.
% 2 Cor. 12:10 echo. "mighty" letterspaced in the Danish, confirmed by
% zoom -- rendered \emph{}.

So he does not reach the ideal straightforwardly; had he reached it straightforwardly, he would, after all, have had to stop. But he is unstoppable,
% --- p. 142 ---
he follows after --- forgetting what is behind, and reaching forth unto what is before, he presses toward the goal.

The remoteness from the goal, that hastens this striving, is again, conversely, a higher expression of the higher approach. The extraordinary messenger does not straightforwardly chase after something straightforwardly unattainable; he is not a reckless one, who strives in haste, not an enthusiast, who works himself up over the impossible.

Here is calm in this unrest, rest in this restlessness, assurance in this striving. For the striving one is, indeed, ``persuaded, that neither death, nor life, nor angels, nor principalities, nor powers, nor things present, nor things to come, nor height, nor depth, nor any other creature, shall be able to separate him'' --- from the ideal itself, personal truth, true personality.
% Rom. 8:38-39. Note the transcription's observation of the Danish's
% spelled-out "ei heller" negation used repeatedly here, for contrast
% with the single bare, doubtful "i" on p.132 flagged above.

Such an assurance is there, then, in such a striving.

May we, then, not be beguiled by the false imagining! If the ideal's messenger cannot receive the assurance straightforwardly, but must strive in conversion, in faith, in order to receive faith's assurance: how could we human beings, then, think of arriving at the recognition of truth by the help of a general thinking, without personal striving?

May we, nonetheless, not be beguiled by the false imagining! If the believing striving is a striving in conversion, then, indeed, no human being can excuse himself, and say: I am willing enough to suppose, that in eternity there is a goal of personality; but it does not help, that I strive: the goal is too high for me.

May we, nonetheless, not be beguiled by the false imagining! When inspiriting is to correspond to exertion, and exertion to inspiriting, then there is here, indeed, no respect of persons, no arbitrary distinction between the extraordinary and the ordinary, but a personal distinction between the striving and the non-striving.

May we, nonetheless, not be beguiled by the false imagining! It looks so strict, so cold, so poor, this striving,
% --- p. 143 ---
in its beginning; but it has, nonetheless, within it an inward richness, a warmth of life, a promise and an assurance, of which the sensuous human being, who sees all this only outwardly, can have no notion.
% The fourfold refrain "May we ... not be beguiled by the false
% imagining!" is checked individually at each of its occurrences in the
% Danish (four on p.142, one more here on p.143) -- plain weight every
% time, no Sperrsatz anywhere, consistent with the earlier finding
% (pp.105-117) that a repeated refrain is not automatically emphasised.

May we, nonetheless, not be beguiled by the false imagining! We love this life, this daylight, these surroundings, this peace, this coziness in this circle among these dearest to us. Existence threatens and troubles, and most of all those who most of all love the present. The present shall be dissolved, this chain shall be broken, this daylight shall be extinguished. May we, nonetheless, not be beguiled by the false imagining! The present is given, only to be taken away again; but it is, after all, taken away, so that the highest may be given. And what is the highest? Not a cold, general being, but an inward life in the kingdom of personality.

All the connections formed here in spirit and truth are, indeed, personal connections. The natural bonds can break; but the personal bonds can never break, as surely as there is an eternal personal truth. The happy must vanish, that the unhappy, too, may vanish; the temporal must yield, that the eternal may triumph; the imperfect must be done away with, that the perfect may come.

The perfect! But who can describe the perfect? --- ``There is one glory of the sun, and another glory of the moon, and another glory of the stars: for one star differeth from another star in glory.'' So is the victory of personality, so is the resurrection of the dead: ``it is sown in corruption, it is raised in incorruption; it is sown in dishonor, it is raised in glory; it is sown in weakness, it is raised in power.'' So testified He, the great messenger, who had not yet attained it, but ``followed after, that he might apprehend it;'' now he has, surely, attained it, but --- the assurance is only for the one who strives.
% 1 Cor. 15:41 and 15:42-43. Checked for letterspacing by zoom in the
% Danish: both citations plain weight, not emphasised.

So let this conclusion, then, be a reminder of the real beginning; so may this assembly, each one especially
% --- p. 144 ---
among my hearers, both men and women, receive my sincere thanks for the kind attention with which you have followed a lecture that has, for me, any significance only insofar as this striving in presentation might point toward an honest will, and, if possible, contain some guidance for a real personal striving.
% Printed p.144 ends the lecture -- and the whole book -- with a printed
% horizontal rule (ornament), not transcribed as text, mirroring
% transcription.tex. Below the rule the rest of the page is blank;
% nothing else follows (no Indhold, no Efterskrift). This is the last
% page of the book.

\end{document}
